%!TEX root = ./../../springer-mono-format/book.tex
%!TEX root = ./../../thesis-trier.tex

% https://english.stackexchange.com/questions/186434/is-notations-a-proper-english-word
Here, we introduce additional concepts and notations %that is only 
needed in \cref{part:constrsync}, and give two new results in coding theory.
Specifically, in \cref{constr_sync:subsec:pfun}, we introduce partial functions. 
\Cref{constr_sync:subsec:pdfa} introduces
partial automata (PDFAs), which generalize our automata model by allowing
the transition relation to be partial function.
These automata give a convenient specification mechanism
for constraint languages (see the discussion in \cref{constrsync:notions:sec:discussion}).
\Cref{constr_sync:subsec:levenshtein-romanov-map} introduces
the Levenshtein-Romanov mapping.
In \cref{constr_sync:chap:classes_constraint_languages} we use prefix codes in some hardness conditions
and solid codes for a statement on the preservation of the computational complexity.
\Cref{constr_sync:subsec:codes} introduces the relevant notions
related to codes, various subclasses of codes as, for example, solid codes,
and overlap-free languages. We also present two new characterizations of solid codes.
 % that have not been noted before.
Later, in \cref{thm:constr_sync_hom_strongly_self_sync}, 
we will prove, using the Levenshtein-Romanov mapping, 
that the computational complexity
of the constrained synchronization is preserved for reencodings
of the constraint language by solid codes.
\Cref{constr_sync:sec:bounded,constr_sync:subsec:monoid} introduce bounded and monoid languages, respectively.
%both will be later investigated 
%as classes of constraint languages in \cref{sec:constr_sync:letter_bounded,constr_sync:sec:monoid}, respectively.
Lastly, in \cref{constr_sync:sec:growth}, we investigate
the growth of languages, and state connections to bounded and monoid languages.
In particular, we classify the growth of monoid languages.

\section{Partial Functions}
\label{constr_sync:subsec:pfun}

%!TEX root = ./../../thesis-trier.tex

\label{def:partial_function} 
\IN{partial function}
A \emph{partial function} $f \colon X \nrightarrow Y$
% is given by a unique set $D \subseteq X \nrightarrow Y$
%such that $f \colon X \to Y$ is a function
either associates with $x \in X$ a unique value, written $\applyfunc{x}{f}$, 
or is undefined for $x \in X$,
and $f(x) = y$ for $x \in X$ and $y \in Y$ means %denotes the fact
that $f(x)$ is defined and equals~$y$.
Correspondingly, $f(x) \ne y$ means 
that either~$f(x)$ is defined but it does not equal $y$ or that $f(x)$ is undefined.
% When writing $\applyfunc{x}{f} = y$ for $x \in X$
%and $y \in Y$ we mean that $\applyfunc{x}{f}$ is defined and equals $y$.
Setting $D = \{ x \in X \mid \applyfunc{x}{f} \text{ is defined } \}$,
we obtain a (total) function $f \colon D \to Y$. 
Every  function can be seen as a partial function defined everywhere.
Given $f \colon X  \nrightarrow Y$ and $g \colon Y  \nrightarrow Z$,
their \emph{composition} is $\concat{f}{g} \colon X \nrightarrow Z$
where
$\applyfunc{x}{(\concat{f}{g})} = z$ for $x \in X$ and $z \in Z$
if and only if $\applyfunc{x}{f}$ is defined with $\applyfunc{\applyfunc{x}{f}}{g} = z$,
%are defined and $\applyfunc{\applyfunc{x}{f}}{g} = z$,
and otherwise $\applyfunc{x}{(\concat{f}{g})}$ is undefined. Other notions also carry over from functions without difficulty,
for example the image of a subset $S \subseteq X$
by setting $\applyfunc{S}{f} = \{ \applyfunc{x}{f} \mid x \in X \text{ and } \applyfunc{x}{f} \text{ is defined } \}$.
\IN{composition!of partial functions}

\section{Partial Automata (PDFAs)}
\label{constr_sync:subsec:pdfa}

%!TEX root = ./../../thesis-trier.tex

\IN{partial deterministic automaton}
\index{partial automaton|see{partial deterministic automaton}}
\index{PDFA|see{partial deterministic automaton}}
%\IN{state set in PDFAs}
%\IN{final state set in PDFAs}
\IN{state set!of a PDFA}
\IN{final state set!of a PDFA}
\IN{start state!of a PDFA}
\IN{transition function!of a PDFA}
\IN{input alphabet!of a PDFA}
Here, we generalize the notion of a complete and deterministic automaton.
A \emph{partial deterministic and finite
	automaton (PDFA)}
is a $5$-tuple $\automat{A}{Q}{\Sigma}{\delta}{q_0}{F}$
where, as for complete and deterministic automata, $Q$ is the \emph{finite state set},~$\Sigma$
the \emph{finite input alphabet}, $q_0$ a \emph{start state}, $F \subseteq Q$ a set of \emph{final states}, but where
 the \emph{transition function} is allowed to be a partial function $\delta \colon Q \times \Sigma \nrightarrow Q$.\footnote{We only consider PDFAs as language acceptors and hence always endow them with a start state and a set of final states.}
So, every complete and deterministic automaton is a partial automaton as well.\footnote{This is actually a rather unfortunate naming. Usually,
	one would expect that when adding modifiers to a concept
	this would narrow down the possibilities, which also allows the description of subtypes by combining modifiers,
	but here adding	the modifier ``partial'' gives a more general concept.
	For this reason, some authors, for example~\cite{DBLP:journals/tcs/CaiF22},
	deviate from the standard terminology by calling the most general model simply finite automaton.}
The modifier \emph{complete}
refers to the fact that the transition function is a total function, i.e., defined for each argument. 

\smallskip
\begin{center}
	\fbox{ 
		\begin{minipage}{28.5em}								
			When we just say ``automaton'' we always mean a complete and deterministic semiautomaton, i.e., 
			the automaton model
			previously considered with totally defined transition functions, and we explicitly
			say PDFA to denote a partial automaton with a start state and set of final states. In the context
			of this work, we only use PDFAs to describe constraint languages. Input automata
			to our problems are always assumed to be complete and deterministic.
		\end{minipage}
	}
\end{center}
\smallskip

\IN{accessible!PDFA}
\IN{accessible!state!in PDFA}
\IN{coaccessible state!in PDFA}
\IN{reachability order!PDFA}
\IN{accepted!language!by PDFA}
\IN{accepted!word!by PDFA}
\IN{PDFA!as language acceptor}
The PDFA $\mathcal A$ can be made complete by adding a non-final trap state
and by adding transitions to this new state for each~$q \in Q$ and $x \in \Sigma$
where $\delta(q, x)$ is undefined.
The notions accessible states, coaccessible states, accessible PDFAs, the reachability order on the states $\lereach$, accepted words,
the accepted language~$L(\mathcal A)$, the notation $\mathcal A_{q,S}$ for $q \in Q$ and $S \subseteq Q$, 
the extension of the partial transition function to words, i.e., 
$\delta \colon Q \times \Sigma^* \nrightarrow Q$, and the extension of~$\delta$ to subsets of states 
$\applytrans{\delta}{S}{u} = \{ \applytrans{\delta}{q}{u} \mid \applytrans{\delta}{q}{u} \text{ is defined and } q \in S \}$
all carry over to PDFAs without difficulty.\footnote{Martyugin~\cite{Martyugin14} defines the extension of the transition function to subsets differently, requiring that $\delta(q,u)$ be defined for every $q \in S$, and leaving $\delta(S,u)$ undefined otherwise. A similar approach for nondeterministic automata was used in~\cite[Sec.~8]{DBLP:conf/aaai/ArrighiFO023} in the context of \emph{conformant planning}. However, apart from these cases, this convention appears to be uncommon in the literature.}
We call~$S \subseteq Q$ an upper set if $\delta(S, x) \subseteq S$ for each $x \in \Sigma$.
We define $\mathcal A_{|S}$, $\mathcal A_{|\Gamma}$,
and $\mathcal A_{|\Gamma,S}$ by restricting the partial transition
function to inputs from $S$,$\Gamma$, and $S \times \Gamma$, respectively,
analogously to the corresponding constructions for the previous automaton model.
The product, power, and pair automata constructions also carry over and will be used in the following
for PDFAs.
A PDFA is called \emph{trim} if every state is accessible and coaccessible. Every non-empty regular language is accepted by some trim PDFA; the empty language $\emptyset$ is the only language that is not accepted by a trim PDFA.
\IN{trim}

\begin{proposition}
	\label{constrsync:prop:trim}
	Let $\automat{A}{Q}{\Sigma}{\delta}{q_0}{F}$
	be a PDFA. Then $\mathcal A$
	is trim if and only if~$[q_0]_{\eqreach}$
	is a minimum in the order induced by $\lereach$
	on the strongly connected components
	and $[q]_{\eqreach} \cap F \ne \emptyset$
	for every maximal strongly connected component $[q]_{\eqreach}$.	
\end{proposition}
\begin{proof}
	Every state is accessible if and only if $q_0 \lereach q$ for all $q \in Q$, i.e., $[q_0]_{\lereach}$
	is a minimum.
	Further, every state is coaccessible
	if and only if for each $q \in Q$ 
	we have $q \lereach p$ for some $p \in F$,
	which is fulfilled if every maximal strongly connected component contains a state from~$F$.
\end{proof}

Minimal PDFAs are characterized by adding the additional requirement that every state is coaccessible
to \cref{prelim:thm:char_minaut}.
The minimal PDFA differs from a minimal complete
automaton if and only if the minimal complete automaton has a non-final trap state,
which can be discarded in this case to obtain the minimal PDFA. 
Note that in a minimal complete automaton for the set of synchronizing words of a given complete and deterministic automaton
every state is coaccessible, so that  it makes no difference 
whether we consider complete or partial automata when investigating minimal deterministic automata for the 
for the set of synchronizing words.

\IN{exactly synchronizing word}
\IN{carefully synchronizing words}
\IN{careful synchronization of automaton}
\IN{exact synchronization of automaton}
Next, we introduce two generalizations of synchronization to PDFAs.
Let $u \in \Sigma^*$.
We call $u$ a \emph{carefully synchronizing word}~\cite{Martyugin14,Rystsov1980} for $\mathcal A$
if there exists $p \in Q$ such that $\delta(q, u) = p$
for all $q \in Q$, i.e., $|\delta(Q, u)| = 1$ and $\delta(q, u)$ is defined
for each $q \in Q$.
We call $u$ an \emph{exactly synchronizing word}~\cite{Shabana_2019,TraverCrutchfield2011}
if $|\delta(Q, u)| = 1$.
Further, $\mathcal A$ is said to be \emph{carefully synchronizing} (resp.~\emph{exactly synchronizing})
if there exists a carefully (resp.~exactly) synchronizing word for~$\mathcal A$.

\section{The Levenshtein-Romanov Mapping}
\label{constr_sync:subsec:levenshtein-romanov-map}

%!TEX root = ./../../springer-mono-format/book.tex
%!TEX root = ./../../thesis-trier.tex

Levenshtein~\cite{Lev64} and Romanov~\cite{Rom66} introduced the following mapping.

\IN{Levenshtein-Romanov mapping}
\begin{definition} \label{def:levenshtein-romanov-mapping}
	Let $L \subseteq \Sigma^*$ with $L \ne \emptyset$. %be a non-empty language.
	The \emph{Levenshtein-Romanov mapping} is $\sigma_L \colon \Sigma^* \to \Sigma^*$
	where $\sigma_L(u)$ for $u \in \Sigma^*$ equals the longest suffix of $u$ that is also a prefix of a word in $L$.
\end{definition} 

For example $\sigma_{\{a,bba\}}(ababb) = bb$. 
Note that $\sigma_L$ is only defined for non-empty languages~$L$.

Let $L \subseteq \Sigma^*$ be non-empty and $u,v\in \Sigma^*$. The following relations
are readily verified:
\begin{align}
	\sigma_L(u) & \in \suff{u}\cap \pref{L}  \label{constrsync:eqn:sigmaL-in-suffpref} \\
	\sigma_L(uv) & \in \Sigma^*v  \text{ for } v \in \pref{L} \\
%	\sigma_L(u) & \in \Sigma^* L \text{ for } u \in \Sigma^* L \\
	\sigma_L(u) & = \sigma_{\pref{L}}(u) \label{constr_sync:prop:sigma_L_prefix:1} \\	
	\sigma_{L}(v) = v & \Leftrightarrow v \in \pref{L}  \label{constr_sync:eqn_sigma_pref1}  \\
	\sigma_{L}(v) = v & \Leftrightarrow v \in \pref{L\setminus \{v\}} \text{ for } v \notin L \label{constrsync:prop:suffsigma:5} \\	
	\sigma_{\{u\}}(v) = u & \Leftrightarrow u \in \suff{v} \label{constr_sync:eqn_sigma_suff0} \\
	\sigma_{\{u\}}(v)  & = \sigma_{\pref{u}\setminus \{u\}}(v) \text{ for } u \notin \suff{v}
\end{align}
By \cref{constr_sync:eqn_sigma_suff0,constr_sync:eqn_sigma_pref1} we have
$\sigma_{\{u\}}(v) \in \{u,v\}$ if and only if $v \in \pref{u} \text{ or } u \in \suff{v}$.
Note that $v \in \pref{L\setminus\{v\}}$ if and only if $v$ is a proper prefix of some word
from~$L$.

We need the following two properties
from~\cite[Prop.~2]{DBLP:journals/tcs/JurgensenS21}. For the convenience of the reader, we include their proofs.

\begin{proposition}[{\cite[Prop.~2]{DBLP:journals/tcs/JurgensenS21}}]
	\label{constr_sync:prop:sigma_L_properties}
	Let $L \subseteq \Sigma^*$ be non-empty and $u,v \in \Sigma^*$.
	\begin{enumerate}
	\item  \label{constr_sync:item:sigma_L_respects_suffix_relation}
		$\sigma_L(v)$ is a suffix of $\sigma_L(uv)$.
		So, $\sigma_L$ preserves the suffix relation, i.e.,
		if $u$ is a suffix of $v$, then $\sigma_L(u)$ is a suffix of $\sigma_L(v)$. 
		
   	\item \label{constr_sync:sigme_L:item:concatenation} $\applyfunc{uv}{\sigma_L} = \applyfunc{\applyfunc{u}{\sigma_L}v}{\sigma_L}$.
	\end{enumerate}
\end{proposition}
\begin{proof}
\ref{constr_sync:item:sigma_L_respects_suffix_relation}: Since $\sigma_L(v)$ is a suffix of $uv$
in $\pref{L}$ and $\sigma_L(uv)$ is the longest suffix of $uv$ in $\pref{L}$,
the suffix $\sigma_L(uv)$ is at least as long as the suffix $\sigma_L(v)$ of~$uv$.
Thus, $\sigma_L(v)$ is a suffix of $\sigma_L(uv)$.

\ref{constr_sync:sigme_L:item:concatenation}: In general, we have
\begin{equation}\label{eqn:sigmaL}
	\sigma_L(uv)
	= 
	\begin{cases}
		\sigma_L(v) & \text{if } \suff{u}v \cap \pref{L} = \emptyset, \\
		wv          & \text{otherwise, where $w \in \suff{u}$ is maximal with $wv \in \pref{L}$.}
	\end{cases}
\end{equation}
(this equation also gives \cref{constr_sync:item:sigma_L_respects_suffix_relation} since $v = x\sigma_L(v)$ for some $x \in \Sigma^*$.)
Now, first suppose $\suff{u}v \cap \pref{L} = \emptyset$.
Then $\suff{\sigma_L(u)}v \cap \pref{L} = \emptyset$,
and so we obtain $\sigma_L(uv) = \sigma_L(v) = \sigma_L(\sigma_L(u)v)$.
Next, suppose $\suff{u}v \cap \pref{L} \ne \emptyset$.
Then $\sigma_L(uv) = wv$ for the longest $w \in \suff{u}$
such that $wv \in \pref{L}$.
As $w \in \pref{L} \cap \suff{u}$,
we find that $w$ is a suffix of $\sigma_L(u)$.
So, $wv$ is also the longest suffix of $\sigma_L(\sigma_L(u)v)$
in $\pref{L}$, i.e., $\sigma_L(\sigma_L(u)v) = wv$.\footnote{For an alternative argument:
	Note that, by \cref{constr_sync:item:sigma_L_respects_suffix_relation}, we have
	$
	\suff{\sigma_L(v)} \subseteq \suff{\sigma_L(\sigma_L(u)v)} \subseteq \suff{\sigma_L(uv)}.
	$
	So, for $\sigma_L(uv) = \sigma_L(v)$
	we get $\sigma_L(\sigma_L(u)v) = \sigma_L(uv)$.
	Otherwise, we have $\sigma_L(uv) = wv$ for some $w \in \Sigma^*$.
	As $w \in \pref{L}$ this yields that $w$ is a suffix of $\sigma_L(u)$
	and so $\sigma_L(\sigma_L(u)v) = wv$.}
\end{proof}

\begin{proposition} \label{constr_sync:prop:sigma_L_prefix}
	Let $L, L' \subseteq \Sigma^*$ be non-empty and $u \in \Sigma^*$.
	\begin{enumerate}		
	\item \label{constr_sync:prop:sigma_L_prefix:2} If $\pref{L}\subseteq \pref{L'}$,
	then $\sigma_{L}(u)$ is a suffix of $\sigma_{L'}(u)$.
		
	\item \label{constr_sync:prop:sigma_L_prefix:3} $\sigma_{L}(u)$ is a suffix of $\sigma_{L L'}(u)$.	

	\item \label{constr_sync:prop:sigma_L_prefix:4a} We have
	\begin{equation} \label{constrsync:notions:eqn:sigmaL-union}
		\sigma_{L \cup L'}(u) 
		=
		\begin{cases}
			\sigma_{L}(u)  & \text{if } \sigma_{L \cup L'}(u)  \in \pref{L}, \\
			\sigma_{L'}(u) & \text{otherwise.}
		\end{cases}	
	\end{equation}
	
	\item \label{constr_sync:prop:sigma_L_prefix:4} We have
	\begin{equation}\label{constrsync:notions:eqn:LR-mapping-equiv-union}
	\sigma_{L \cup L'}(u)  \in \pref{L}
	\Leftrightarrow
	\sigma_{L'}(u) \in \suff{\sigma_L(u)}
	\Leftrightarrow
	|\sigma_L(u)| \ge |\sigma_{L'}(u)|
	\end{equation}

	\item \label{constr_sync:prop:sigma_L_prefix:5} We have
	$
	\sigma_L(u) \in \suff{L'}
	$
	if and only if $\sigma_{\{v\}}(u) \in \suff{L'}$ for all $v \in L$.
	
	\item \label{constr_sync:prop:sigma_L_prefix:6} We have
	$
	\suff{\sigma_L(u)} = \bigcup_{v \in L} \suff{\sigma_{\{v\}}(u)}. 
	$
	
	\item\label{constr_sync:prop:sigma_L_prefix:7}
	We have $\sigma_L(u) = \varepsilon$ if and only if $\sigma_{\{v\}}(u) = \varepsilon$ for each $v \in L$.
	\end{enumerate}
\end{proposition}
\begin{proof}
\ref{constr_sync:prop:sigma_L_prefix:2}: Every suffix of~$u$ that is a prefix of a word from~$L$
is also a suffix that is a prefix of a word from~$L'$. Thus, $\sigma_{L'}(u)$
is longer than $\sigma_{L}(u)$, which yields that $\sigma_{L}(u)$
is a suffix of~$\sigma_{L'}(u)$.

\ref{constr_sync:prop:sigma_L_prefix:3}: We have $\pref{L} \subseteq \pref{L L'}$.
So this follows from~\cref{constr_sync:prop:sigma_L_prefix:2}.	

\ref{constr_sync:prop:sigma_L_prefix:4a}: Suppose $\sigma_{L \cup L'}(u) \in \pref{L}$.
By \cref{constr_sync:prop:sigma_L_prefix:2}, $\sigma_L(u)$ is a suffix of $\sigma_{L\cup L'}(u)$.
Since $\sigma_{L \cup L'}(u)  \in \pref{L}$, $\sigma_{L\cup L'}(u)$
is a suffix of $\sigma_L(u)$ by the definition of~$\sigma_L$.
Thus $\sigma_{L\cup L'}(u) = \sigma_{L}(u)$.
%
% Then the longest suffix of~$u$ in $\pref{L\cup L'}$ 
% is contained in $\pref{L}$, and thus equals
% the longest suffix of~$u$ in $\pref{L}$, i.e.,
% we have $\sigma_{L\cup L'}(u) = \sigma_{L}(u)$.
 Otherwise, suppose $\sigma_{L \cup L'}(u) \notin \pref{L}$. 
 Since $\sigma_{L \cup L'}(u) \in \pref{L \cup L'}$,
 this yields $\sigma_{L \cup L'}(u) \in \pref{L'}$.
 With the same argument as before, we find 
 $\sigma_{L \cup L'}(u) = \sigma_{L'}(u)$.

\ref{constr_sync:prop:sigma_L_prefix:4}: First, suppose $\sigma_{L \cup L'}(u)  \in \pref{L}$.
 Then $\sigma_{L \cup L'}(u) = \sigma_L(u)$ by \cref{constr_sync:prop:sigma_L_prefix:4a}.
 Thus, by \cref{constr_sync:prop:sigma_L_prefix:2}, we have $\sigma_{L'}(u) \in \suff{\sigma_L(u)}$.
 
 Next, suppose $\sigma_{L'}(u) \in \suff{\sigma_L(u)}$. 
 Then the longest suffix of~$u$ in~$L$ is longer than the longest
 suffix of~$u$ in~$L'$.
 This gives $\sigma_{L \cup L'}(u) = \sigma_L(u) \in \pref{L}$.
 
 The last equivalence follows easily since $\sigma_L(u)$ and $\sigma_{L'}(u)$
 are both suffixes of~$u$. This shows \cref{constrsync:notions:eqn:LR-mapping-equiv-union}.
 
\ref{constr_sync:prop:sigma_L_prefix:5}: First, suppose $\sigma_L(u) \in \suff{L'}$.
By \cref{constr_sync:prop:sigma_L_prefix:2}, $\sigma_{\{v\}}(u)$ is a suffix of~$\sigma_L(u)$ for $v \in L$.
Thus we find $\sigma_{\{v\}}(u) \in \suff{L'}$.

Conversely, suppose $\sigma_{\{v\}}(u) \in \suff{L'}$ for all $v \in L$.
We have $\sigma_L(u) = \sigma_{\{v\}}(u)$ for some $v \in L$.
This yields $\sigma_L(u) \in \suff{L'}$.

\ref{constr_sync:prop:sigma_L_prefix:6}: By \cref{constr_sync:prop:sigma_L_prefix:2}, we have $\sigma_{\{v\}}(u) \in \suff{\sigma_L(u)}$ for $v \in L$.
 Thus $\bigcup_{v \in L} \suff{\sigma_{\{v\}}(u)} \subseteq \suff{\sigma_L(u)}$.
 
 Since $\sigma_L(u) \in \suff{\sigma_L(u)}$,
 by \cref{constr_sync:prop:sigma_L_prefix:6} we find $\sigma_{\{v\}}(u) \in \suff{\sigma_L(u)}$
 for all $v \in L$. This gives $\bigcup_{v \in L} \suff{\sigma_{\{v\}}(u)} \subseteq \suff{\sigma_L(u)}$.
 
\ref{constr_sync:prop:sigma_L_prefix:7}: This follows from \cref{constr_sync:prop:sigma_L_prefix:5} 
by considering $L' = \{\varepsilon\}$.
\end{proof}

\begin{proposition}
	Let $L \subseteq \Sigma^*$ and $u \in \Sigma^*$ such that $L \setminus\{u\} \ne \emptyset$.
	Then we have 
	\begin{equation}\label{constrsync:eqn:sigma_prefL}
		\sigma_{\pref{L} \setminus \{u\}}(u) 
		=
		\begin{cases}
			\sigma_{L\setminus \{u\}}(u) & \text{if } \sigma_{\pref{L}\setminus \{u\}}(u) \in \pref{L \setminus \{u\}}, \\
			\sigma_{\pref{u}\setminus\{u\}}(u) & \text{otherwise.}
		\end{cases}
	\end{equation}
\end{proposition}
\begin{proof}
	We have
	\[
	\pref{\pref{L} \setminus \{u\}} = 
	\begin{cases}
		\pref{L \setminus \{u\}} \cup ( \pref{u} \setminus \{u\} ) & \text{if } u \in \pref{L}, \\
		\pref{L \setminus \{u\}} & \text{otherwise.}
	\end{cases}
	\]
	
	First, suppose $u \notin \pref{L}$.
	Then $u \notin L$
	and $\pref{L} = \pref{L \setminus \{u\}}$.
	By \cref{constrsync:eqn:sigmaL-in-suffpref}, we have $\sigma_{\pref{L}\setminus\{u\}}(u) \in \pref{L}$.
	Thus, in this case, the first case from \cref{constrsync:eqn:sigma_prefL}
	applies.
	Since we have $\sigma_{\pref{L}\setminus\{u\}}(u) = \sigma_{\pref{L}}(u) = \sigma_L(u) = \sigma_{L\setminus \{u\}}(u)$,
	we find that \cref{constrsync:eqn:sigma_prefL} holds.
	
	Next, suppose $u \in \pref{L}$.
	Then we find	
	\begin{align*}
		& \sigma_{\pref{L} \setminus \{u\}}(u) \\
		& = \sigma_{\pref{\pref{L} \setminus \{u\}}}(u) & \tag{Eqn.~\eqref{constr_sync:prop:sigma_L_prefix:1}} \\
		& = \sigma_{\pref{L \setminus \{u\}} \cup ( \pref{u} \setminus \{u\} )}(u) \\
		& = \begin{cases}
			\sigma_{\pref{L \setminus \{u\}}}(u) & \text{if } \sigma_{\pref{L} \setminus \{u\}}(u) \in \pref{L \setminus \{u\}} \\
			\sigma_{\pref{u}\setminus\{u\}}(u) & \text{otherwise}		
		\end{cases}	& \tag{Eqn.~\eqref{constrsync:notions:eqn:sigmaL-union}}
	\end{align*}
	Since, by \cref{constr_sync:prop:sigma_L_prefix:1}, we have $\sigma_{\pref{L \setminus \{u\}}}(u) = \sigma_{L \setminus \{u\}}(u)$,
	the claim follows.
\end{proof}

The Levenshtein-Romanov mapping can be used to describe the Myhill-Nerode equivalence classes
for a left ideal as shown next. 

Note that for a non-empty left ideal $I$, 
we always have $\sigma_I(u) = u$ since $\pref{I} = \Sigma^*$.

\begin{proposition} \label{prelim:formlang:prop:left-ideal}
	Let $L \subseteq \Sigma^*$ be non-empty and consider the left ideal $I = \Sigma^*L$.
	Then for $u \in \Sigma^*$ we have $u \equiv_I \sigma_L(u)$.	
\end{proposition}
\begin{proof}
	Let $u \in \Sigma^*$. Next, we show that $uv \in I$
	if and only if $\sigma_L(u)v \in L$ for all $v \in \Sigma^*$,
	which is equivalent to $u \equiv_L \sigma_L(u)$.
		
	First, suppose we have $uv \in I$ for some $v \in \Sigma^*$.
	If $v \in I$, then $\sigma_L(u)v \in I$ since~$I$ is a left ideal.
	Next, let $v \notin I$.
	In this case, since $uv \in I$,
	there must exist $x,y \in \Sigma^*$
	with $y \ne \varepsilon$
	such that $u = xy$ and $yv \in L$.
	Since $y \in \pref{L}$, $y$	is a suffix of $\sigma_L(u)$.
	Thus we also have $\sigma_L(u)v \in I$.
	
	Conversely, suppose $\sigma_L(u)v \in I$.
	Since $\sigma_L(u)$ is a suffix of~$u$,
	this clearly gives $uv \in I$. 	
\end{proof}

\begin{corollary} \label{prelim:formlang:cor:left-ideal}
	Let $I = \Sigma^*L$ for some $L \subseteq \Sigma^*$.
	Then $\{ [w]_{\equiv_I} \mid w \in \Sigma^* \}
	= \{ [w]_{\equiv_I} \mid w \in \pref{L} \}$.	
\end{corollary}
\begin{proof}
	This follows from \cref{prelim:formlang:prop:left-ideal} 
	and since $\sigma_L(u) \in \pref{L}$ for all $u \in \Sigma^*$.
\end{proof}

Next, we give some equalities involving suffix sets.

\begin{proposition} \label{constrsync:prop:suffsigmaB}
	Let $L \subseteq \Sigma^*$ and $u \in \Sigma^+$. Then the following holds:
	\begin{enumerate}
		\item \label{constrsync:prop:suffsigma:1B} $\suff{\sigma_L(u)} \cap \pref{L} = \suff{u}\cap \pref{L}$.
		
		\item \label{constrsync:prop:suffsigma:B4} $\suff{\sigma_{\pref{L}\setminus\{u\}}(u)} \cap \pref{\pref{L} \setminus\{u\}}
		=
		\suff{u} \cap \pref{\pref{L} \setminus \{u\}}$.
		
		\item \label{constrsync:prop:suffsigma:2B} $\suff{\sigma_{\pref{L}\setminus\{u\}}(u)} \cap L = \suff{\sigma_{L\setminus \{u\}}(u)} \cap L$ if $L \setminus \{u\}$ is non-empty.
				
		\item \label{constrsync:prop:suffsigma:3B} $\suff{\sigma_{\pref{L}\setminus \{u\}}(u)} \cap L = ( \suff{u}  \cap ( L \setminus \{u\} ) ) \cup( \{u\} \cap L \cap \pref{L\setminus\{u\}} )$.
	\end{enumerate}
\end{proposition}
\begin{proof}
	\ref{constrsync:prop:suffsigma:1B}: Since $\sigma_L(u)$ is a suffix of $u$,
	we have $\suff{\sigma_L(u)} \subseteq \suff{u}$
	and so $\suff{\sigma_L(u)} \cap \pref{L} \subseteq \suff{u} \cap \pref{L}$.
	Next, let $v \in \suff{u}\cap \pref{L}$.
	Since $\sigma_L(u)$ equals the longest suffix of~$u$ that is in $\pref{L}$,
	this suffix has length at least $|v|$. Thus~$v$
	is a suffix of $\sigma_L(u)$, which gives $v \in \suff{\sigma_L(u)} \cap \pref{L}$.	
	
	\ref{constrsync:prop:suffsigma:B4}: This follows from \cref{constrsync:prop:suffsigma:1B} 
	applied to $\pref{L}\setminus\{u\}$.
	
	\ref{constrsync:prop:suffsigma:2B}:
	Since $L \setminus \{u\} \subseteq \pref{L}\setminus\{u\}$, by \cref{constr_sync:prop:sigma_L_prefix:2} of \cref{constr_sync:prop:sigma_L_prefix} we have
	\[
	\suff{\sigma_{L\setminus \{u\}}(u)} \cap L \subseteq
	\suff{\sigma_{\pref{L}\setminus\{u\}}(u)} \cap L.
	\]
	Next, let $v \in \suff{\sigma_{\pref{L}\setminus\{u\}}(u)} \cap L$.
	Suppose $v = u$.
	Then $u = \sigma_{\pref{L}\setminus\{u\}}(u)$.
	This yields that $u$ is a proper prefix of a word from $L$
	and thus $u = \sigma_{L\setminus\{u\}}(u)$.
	If $v \ne u$, then~$v$ is a suffix of $u$ in $L \setminus\{u\}$
	and thus it is a suffix of $\sigma_{L\setminus\{u\}}(u)$.
	In both cases, we find $v \in \suff{\sigma_{L\setminus\{u\}}(u)} \cap L$.
	
	\ref{constrsync:prop:suffsigma:3B}:
	By \cref{constrsync:prop:suffsigma:B4} we have 
	$\suff{u}  \cap ( L \setminus \{u\} ) = \suff{\sigma_{\pref{L}\setminus\{u\}}(u)} \cap ( L \setminus \{u\} )$.
	By \cref{constrsync:prop:suffsigma:5},
	$\sigma_{\pref{L}\setminus\{u\}}(u) = u$
	if and only if $u$ is a proper prefix of a word from $\pref{L}\setminus\{u\}$.
	This is equivalent to~$u$ being a proper prefix of a word from~$L$, i.e,
	$u \in \pref{L \setminus \{u\}}$.
	Thus $u \in \suff{\sigma_{\pref{L}\setminus\{u\}}(u)} \cap L$
	if and only if $u \in L \cap \pref{L\setminus\{u\}}$,
	which gives the claim.
\end{proof}

\begin{proposition} \label{constrsync:prop:sigma-uvA}
	Let $L, U, V \subseteq \Sigma^*$ and $v \in \pref{L}$.
	Then the following are equivalent.
	\begin{enumerate}
			
	\item \label{constrsync:prop:sigma-uvA:1} $\sigma_L(uv) \in U$ for some $u \in \pref{L} \cap \suff{V}$.
		
	\item \label{constrsync:prop:sigma-uvA:2} $\sigma_L(uv) \in U$ for some $u \in \suff{V}$. 
		
	\item \label{constrsync:prop:sigma-uvA:3} $( \pref{L}\cap U ) \cap ( \pref{L} \cap \suff{V} ) v \ne \emptyset$.
		
	\item \label{constrsync:prop:sigma-uvA:4} $( \pref{L}\cap U ) \cap \suff{V}v \ne \emptyset$.
			
	\item \label{constrsync:prop:sigma-uvA:5} $\sigma_L(uv) = uv$ and $uv \in U$ for some $u \in \suff{V}$.
	\end{enumerate}
\end{proposition}
\begin{proof}
$\ref{constrsync:prop:sigma-uvA:1}\Rightarrow\ref{constrsync:prop:sigma-uvA:2}$: This is clear.

$\ref{constrsync:prop:sigma-uvA:2}\Rightarrow\ref{constrsync:prop:sigma-uvA:3}$: Since $v \in \pref{L}$, we have $\sigma_L(uv) = wv \in \pref{L}$ for some $w \in \Sigma^*$.
We have $w \in \pref{L}$ and the word~$w$ is a suffix of~$u$, and so $w \in \suff{V}$.
 Thus we have $wv \in \pref{L}\cap U$ and $wv \in \suff{V}v$.
 
$\ref{constrsync:prop:sigma-uvA:3}\Rightarrow\ref{constrsync:prop:sigma-uvA:4}$: This is clear.

$\ref{constrsync:prop:sigma-uvA:4}\Rightarrow\ref{constrsync:prop:sigma-uvA:5}$: Let $u \in \suff{V}$ with $uv \in \pref{L}\cap U$.
 By \cref{constr_sync:eqn_sigma_pref1} we have $\sigma_L(uv) = uv$.

$\ref{constrsync:prop:sigma-uvA:5}\Rightarrow\ref{constrsync:prop:sigma-uvA:1}$: Since $\sigma_L(uv) \in \pref{L}$,
 we get $uv \in \pref{L}\cap U$ and $u \in \pref{L} \cap \suff{V}$.
\end{proof}

\begin{proposition} \label{constrsync:prop:sigma-uvpref}
	Let $L, U, V \subseteq \Sigma^*$ with $v \in \pref{L}$
	and $U \subseteq \pref{L}$.	Then the following are equivalent.
	\begin{enumerate}
		\item \label{constrsync:prop:sigma-uvpref:1} $\sigma_L(uv) \in U$ for some $u \in \suff{V}$.
				
		\item \label{constrsync:prop:sigma-uvpref:2} $uv \in U$ for some $u \in \suff{V}$.
		
		\item \label{constrsync:prop:sigma-uvpref:3} $uv \in U$ for some $u \in \pref{L} \cap \suff{V}$.
	\end{enumerate}
\end{proposition}
\begin{proof}
$\ref{constrsync:prop:sigma-uvpref:1}\Rightarrow\ref{constrsync:prop:sigma-uvpref:2}$: Since $v \in \pref{L}$ 
we have $\sigma_L(uv) = wv$ for some $w \in \suff{V}$. This gives the claim.

$\ref{constrsync:prop:sigma-uvpref:2}\Rightarrow\ref{constrsync:prop:sigma-uvpref:3}$: Since $uv \in \pref{L}$ 
we have $u \in \pref{L}$, which gives the claim.

$\ref{constrsync:prop:sigma-uvpref:3}\Rightarrow\ref{constrsync:prop:sigma-uvpref:1}$: Since $uv \in \pref{L}$ 
we have $\sigma_L(uv) = uv$.
\end{proof}

\begin{remark} The following are equivalent: $L = \{\varepsilon\}$, $\sigma_L(y) = \varepsilon$
	for all $y \in \Sigma^*$, and $\sigma_L(xy) = \sigma_L(y)$
	for all $x,y \in \Sigma^*$.
\end{remark}

\begin{remark} For $v \in \Sigma^+$, the following are equivalent:
	$\sigma_L(uv) = \varepsilon$ for all $u \in \Sigma^*$,
	$\sigma_L(uv) = \varepsilon$ for all $u \in \pref{L}$,
	$\sigma_L(v) = \varepsilon$
	and $wv \notin \pref{L}$ for all $w \in \Sigma^+$,
	$\sigma_L(v) = \varepsilon$
	and $wv \notin \pref{L}$ for all $w \in \pref{L}$.	
\end{remark}

\section{Codes and Overlap-Free Languages}
\label{constr_sync:subsec:codes}

%!TEX root = ./../../springer-mono-format/book.tex
%!TEX root = ./../../thesis-trier.tex

Here, in \cref{constrsync:sec:codes:basic-notions} we introduce codes and overlap-free languages
and state some basic relations. Some of our later results will use prefix codes and solid codes. 
A few conditions using prefix codes use constructions related to maximal prefix codes. This
class of codes is discussed in \cref{constrsync:sec:max-prefix-codes}.
In \cref{constrsync:sec:more-on-prefix-codes} we collect properties of prefix codes that we need for those results.
\Cref{constrsync:sec:codes:levensthein-romanov} states some further characterizations of our classes of codes and overlap-free languages and relations to the Levenshtein-Romanov mapping.
\Cref{constrsync:sec:codes:solid} collects characterizations of solid codes (some of them have
not been observed before).

\subsection{Basic Notions}
\label{constrsync:sec:codes:basic-notions}

Here, in \cref{constrsync:def:codes}, we introduce overlap-free languages,
codes, and various subclasses of codes as, for example, solid codes, comma-free codes, and infix codes.
%Here, in \cref{constrsync:def:codes}, we introduce codes, various subclasses thereof as, for example, solid codes, comma-free codes, and infix codes,
%and overlap-free languages.
\Cref{constr_sync:prop:selfsync_implies_avoiding} 
proves that every solid code is a comma-free code, and that every comma-free code is an infix code.
In \cref{constrsync:ex:binarysolidcode} we also show that overlap-free languages
are not codes in general. The remaining inclusions among infix, bifix, prefix, and suffix codes
are obvious. \Cref{fig:code-hasse} gives an overview of the inclusion relations among
the introduces code classes and overlap-free languages.
As formally stated in \cref{constr_sync:prop:coding_morphisms},
codes are used to encode symbols from some source alphabet.
%Codes are used to encode symbol from some source alphabet,
%a formal statement of this fact is given
%in \cref{constr_sync:prop:coding_morphisms}.
%In \cref{constr_sync:prop:coding_morphisms} we formalize 
%the idea that codes encode symbols from a source alphabet.
%the relation
%between a source alphabet and codes as an encoding of symbols from such an alphabet.
\Cref{constr_sync:lem:two-words-code,constrsync:prop:conjugate-words} 
give two combinatorial results on word equations.

% L = emptyset erlauben? macht für präfixcode bei einigen aussagen sinn.
\begin{definition} \label{constrsync:def:codes}
	Let $L \subseteq \Sigma^+$ be non-empty and $u,v,w \in \Sigma^*$.
	
	\IN{code}
	\IN{overlap-free code}
	\IN{infix code}
	\IN{prefix code}
	\IN{suffix code}
	\IN{bifix code}
	\IN{solid code}
	\IN{comma-free code}
	\index{unique factorization property|see{code}}
	\begin{enumerate}%[itemsep=-2pt]	
		\item \label{constrsync:item:code} $L$ is a \emph{code} if, for $u_1, \ldots, u_n, v_1, \ldots, v_m \in L$,
		the equality
		\[
		u_1 \cdots u_n = v_1 \cdots v_m
		\]
		implies $n = m$ and $u_i = v_i$ for all $i \in \{1,\ldots,n\}$ (the \emph{unique factorization property} in $L^*$ w.r.t.~$L$).
		
		\item  $L$ is \emph{overlap-free}		
		if $uv, vw \in L$ with $u \ne\varepsilon$ and  $w \ne \varepsilon$ implies $v = \varepsilon$, i.e., no proper non-empty suffix
		of a word from~$L$ is also a proper non-empty prefix of some word in~$L$.\footnote{The condition
			that~$u$ and~$v$ are non-empty is missing in~\cite{DBLP:journals/tcs/JurgensenS21},
			but it is essential since without it,
			for $v \in L$, we have $\varepsilon\cdot v, v \cdot \varepsilon \in L$
			and thus $v = \varepsilon$, which
			contradicts the fact that $L$ is a subset of $\Sigma^+$.}		
		
		%\item $L$ is a \emph{block code} if it is a code and there exits $n > 0$ such that $L \subseteq \Sigma^n$.
		
		\item $L$ is an \emph{infix code} (or \emph{infix-free}) if $u, vuw \in L$ implies $v = w = \varepsilon$.
		
		\item $L$ is a \emph{prefix code} (or \emph{prefix-free}) if $uv, u \in L$ implies $v = \varepsilon$.
		
		\item $L$ is a \emph{suffix code} (or \emph{suffix-free}) if $vu, u \in L$ implies $v = \varepsilon$.
		
		\item $L$ is a \emph{bifix code} if it is a prefix and a suffix code. 
		
		\item $L$ is a \emph{solid code} if it is an overlap-free infix code.
		
		\item $L$ is a \emph{comma-free code} if $L^2 \cap \Sigma^+ L \Sigma^+ = \emptyset$.\footnote{The name comma-free refers to the fact that, when facing a stream of code words,
			we can uniquely decipher the next code words when 
			starting to read from somewhere in the middle of the stream
			without the need of a special separation symbol.}
	\end{enumerate}
\end{definition}

\Cref{constrsync:fig:freeness-fails} gives graphical depictions when some of the defining
conditions from \cref{constrsync:def:codes} fail.

\begin{table}[htbp]
	\centering
	\caption[Failure of defining conditions for various language classes]{Situations where the defining 
		conditions for overlap-, infix-, prefix-, suffix-, and comma-freeness fail. 
		Boxes labeled $u,v,w$ represent words or factors, and their relative positions indicate the overlaps.}
	\label{constrsync:fig:freeness-fails}
	
	\setlength{\aboverulesep}{6pt}
	\setlength{\belowrulesep}{6pt}
	
	\begin{tabular}{m{4cm} m{6cm}}
		\toprule
		not overlap-free & \tikzcell{	
			\node[normal] (A1) {$u$};
			\node[overlap, right=0pt of A1] (A2) {$v$};
			
			\node[overlap, below=-0.5pt of A2] (B1) {$v$};
			\node[normal, right=0pt of B1] (B2) {$w$};            
		} \\
		\midrule
		
		not infix-free & \tikzcell{
			\node[draw, minimum height=6mm, minimum width=4cm] (u) {$u$};
			\node[normal, below=0pt of u] (v) {$v$};
		} \\
		\midrule
		
		not prefix-free & \tikzcell{        
			\node[draw, minimum height=6mm, minimum width=4cm] (u) {$u$};
			\node[normal, anchor=north west] (v) at (u.south west) {$v$};        
		} \\
		\midrule
		
		not suffix-free & \tikzcell{
			\node[draw, minimum height=6mm, minimum width=4cm] (u) {$u$};
			\node[normal, anchor=north east] (v) at (u.south east) {$v$};        
		} \\
		\midrule
		
		not comma-free & \tikzcell{
			\node[normal] (u) {$u$};
			\node[normal, right=0pt of u] (v) {$v$};
			
			\node[normal, anchor=north] (w) at ($(u.south)!0.5!(v.south)$) {$w$};        
		} \\
		\bottomrule
	\end{tabular}
\end{table}


By definition a code is a subset of $\Sigma^+$, but as $\varepsilon = \varepsilon \cdot \varepsilon$
a language fulfilling \crefitem{constrsync:def:codes}{constrsync:item:code}  
never contains the empty word and we could 
replace $\Sigma^+$ by~$\Sigma^*$  in the definition.
%equivalently define a code as a subset of $\Sigma^*$. 
However, the defining conditions 
for overlap-free languages, infix, prefix, suffix, bifix, solid, and comma-free codes
are also fulfilled for $L = \{\varepsilon\}$ and the exclusion of the empty word makes
a difference in these cases.

\begin{example} \label{constrsync:ex:binarysolidcode}
	\begin{enumerate}
		\item The set $\Sigma$ is a code, as $\Sigma^*$ fulfills the unique factorization property w.r.t.~$\Sigma$.
		
		\item For $n \ge 1$, the set $\Sigma^n$ is a code. A subset of~$\Sigma^n$
		is a solid code if and only if it is overlap-free. Thus, $\Sigma$
		is a solid code.
		
		\item \label{constrsync:ex:comma-free-single-word}  $\{ w \}$ is a code for $w \in \Sigma^+$. 
		Here, we have that 
		$\{ w \}$ is not a comma-free code if and only if $w = uv$
		for $u,v \in \Sigma^*$ such that $uv = vu$.
		This follows since $uv = vu$ implies $uvuv \in \Sigma^+ vu \Sigma^+$
		and, conversely, $ww = x w y$ for $x,y \in \Sigma^+$
		implies that we can write $w = xu = vy$ for $u,v \in \Sigma^+$
		with $uv = w$. Now $xu = uv$ implies $|x| = |v|$
		and similarly we obtain $|u| = |y|$. Thus, since $x$ and $v$
		are equal-length prefixes of~$w$ we have $x = v$
		and similarly $u = y$ and thus $vu = w = uv$.
		
		\item $\{a,ab\}$ is an overlap-free code that is not a prefix code.
		
		\item $\{a,ab,aab\}$ is overlap-free, but not a code since $(aab) = (a)(ab)$.
		
		\item $\{abc,b\}$ is an overlap-free bifix code that is not an infix code.
		
		\item $a^*b$ is a prefix code that is overlap-free and not a suffix code. 
		
		\item  \label{constrsync:ex:binarysolidcodeitem}
		%solid code
		%abab
		%aaabbb
		%aabab
		$\{ aba^ibb \mid i \ge 1\}$ is a solid code. 
		
		\item The language~$C$ from \cref{prelim:lem:alphabet-reduction} forms a solid code.
		
		\item \label{constrsync:ex:comma-free-not-solid} $\{abb,aba\}$ is a comma-free code, but not a solid code since
	    it is not overlap-free.
		
		\item \label{constrsync:ex:infix-not-comma-free} $\{ab,ba\}$ is an infix code, but not a comma-free code.
		
		\item Let $C \subseteq \Sigma^+$ be a prefix code. Then $C\Sigma^n$ is a prefix code for each $n \ge 1$.
			
		\item For $L \subseteq \Sigma^*$ with $L \cap \Sigma^+ \ne \emptyset$,
		the language $L \setminus ( L\Sigma^+ \cup \{\varepsilon\} )$ is a prefix code.
		
		\item \label{constrsync:ex:constr-solid-code} Suppose $C \subseteq \Sigma^+$ is a finite prefix code
		such that $c\Sigma^* \cap C = C \cap \Sigma^*c = \emptyset$
		for some $c \in \Sigma$.
		Set $k = \max\{ l \ge 0 \mid \Sigma^* c^l \Sigma^* \cap C \ne \emptyset \} + 1$.
		For example, when applying the construction to the prefix code $C = \{bcb, bccb\}$
		we obtain $\{cccbcb, cccbccb\}$.
		Then, $c^k C$ is a solid code  as we show next.
		
		\underline{1. $c^kC$ is an infix code:} Assume $xc^kuy = c^kv$ with $u,v \in C$ and $x,y \in \Sigma^*$.
		As $c^k$ is not a factor of $v$, 
		we have $|x| < k$, which gives $x \in c^*$.
		Since $v \notin c\Sigma^*$
		we conclude $|x| = 0$.
		Hence we have $uy = v$ and as $C$ is a prefix code we get $y = \varepsilon$.
		So $c^kC$ is an infix code. 
		
		
		\underline{2. $c^kC$ is overlap-free:} Assume $xy = c^ku$ and $yz = c^kv$
		for $u,v \in C$ and $x,y,z \in \Sigma^*$.
		%with $x,z \ne \varepsilon$.
		For $|y| \le k$, as $yz = c^kv$ we have $y \in c^*$
		and, as $u \notin \Sigma^*c$, from $xy = c^ku$ we get $|y| = 0$.
		For $|y| > k$, from $yz = c^kv$ we find $y = c^kw$ for some $w \in \pref{v}$.
		Then, from $xy = c^ku$, as $c^k \notin \factor{u}$, we must have $|x| < k$.
		Further, as $u \notin c\Sigma^*$, $xy = c^ku$ then gives $|x| = 0$, which 
		implies that $c^ku$ is a prefix of~$c^kv$.
		As $c^kC$ is an infix code, we find $|z| = 0$ and so $y = c^ku = c^kv$.
		Hence $c^kC$ is overlap-free.
		% Assume $y \ne \varepsilon$.
		% As $u \notin \Sigma^*c$, we have $y = c^kv'$ for some $v' \in \Sigma^+$.
		% Then, as $u \notin c\Sigma^*$ we have $|x| > k$.
		% However, this implies $c^k \in \fact{u}$, which 
		% is not possible. Hence, we have $y = \varepsilon$
		% and so~$c^kC$ is overlap-free.
		
		\item By the previous item, for a finite prefix code $C \subseteq \Sigma^+$
		and distinct $c,a \in \Sigma$ with 
		$k = \max\{ l \ge 0 \mid \Sigma^* c^l \Sigma^* \cap C \ne \emptyset \} + 1$
		we have that $c^k a C a$ is a solid code
		and $c^k aC$ and $c^k Ca$ are prefix codes
		if $C \cap \Sigma^*c = \emptyset$ and $c \Sigma^* \cap C = \emptyset$, respectively,
		since $aC$, $Ca$ and $aCa$ are prefix code as well. 
		
	\end{enumerate}
\end{example} 

\begin{remark}		
	We actually came up with the construction from \cref{constrsync:ex:constr-solid-code}
	in \cref{constrsync:ex:binarysolidcode} by analysing
	a result from~\cite[Thm.~29]{FernauGHHVW19}
	and extending it in~\cite[Thm.~17]{DBLP:conf/fct/Hoffmann21} (see also the discussion on page~\pageref{discussion:strongly-self-sync-codes}).	
	However, a closely related variant of this construction is much older and was introduced 
	by Gilbert~\cite{Gilbert1960Synchronization} in the context of synchronization between sender and receiver during message transmission
	and by Levenshtein~\cite{Lev64,levenshtein1970maximum_en} 
	to construct maximal solid codes of equal-length words.\footnote{Levenshtein referred to them as non-overlapping codes. In fact, 
		codes of equal-length words
		are infix codes. Thus, such a code forms a solid code if and only if it is overlap-free.}
	Levenshtein considers codes of the form $c^k C$, where
	$C = \{ u \in \Sigma^{n} \mid u \notin \Sigma^* c^k \Sigma^* \text{ and } u \notin c\Sigma^* \text{ and } u \notin \Sigma^* c \}$
	for $1 \le k < n$ and $c \in \Sigma$.
	Further generalizations appear in~\cite[Construction~2]{Blackburn2015NonOverlappingCodes}. It was also independently rediscovered in~\cite{CheeKiahPurkayasthaWang2013}.
\end{remark}

A prefix code is also characterized by the property
that for code words $u,v$ and any words $w,w'$
the equality $uw = vw$ implies $u = v$
since $uv = vw$ implies that either $u$ is a prefix of~$v$
or vice versa by the unique factorization property of~$\Sigma^*$
w.r.t.~$\Sigma$ and so both must be equal.
Now, as the naming anticipates, this yields that prefix 
codes (and similarly suffix codes) are indeed codes: The equation $u_1 \cdots u_n = v_1 \cdots v_m$ implies
$u_1 = v_1$ and $u_2 \cdots u_n = v_2 \cdots v_m$ for prefix codes (resp.~$u_m = v_m$ and $u_1 \cdots u_{n-1} = v_1 \cdots v_{m-1}$ for suffix codes).
Continuing, we ultimately find $n = m$ and $u_i = v_i$ 
for all $i \in \{1,\ldots,n\}$. 
Thus, as subclasses, bifix and infix codes are also codes.
Furthermore, as we show next, solid and comma-free codes are infix codes,
yielding that all the mentioned classes except overlap-free languages
are codes. See \cref{fig:code-hasse} for an overview of the inclusion relations.

\begin{proposition}[\cite{DBLP:reference/hfl/JurgensenK97,DBLP:journals/tcs/JurgensenS21}]
	\label{constr_sync:prop:selfsync_implies_avoiding}
	\begin{enumerate} 
		\item \label{constr_sync:prop:item:sc-in-cf} The class of solid codes is properly contained in the 
		class of comma-free codes. 
		
		\item \label{constr_sync:prop:item:cf-in-ic} The class of comma-free codes is properly contained in the class of infix codes.
	\end{enumerate} 	
\end{proposition}
\begin{proof}
	A comma-free code that is not a solid code
	and an infix code that is not comma-free
	are given in \cref{constrsync:ex:comma-free-not-solid} and \cref{constrsync:ex:infix-not-comma-free} 
	of \cref{constrsync:ex:binarysolidcode}, respectively.
	Next, we show the claimed inclusions.	
		
	\ref{constr_sync:prop:item:sc-in-cf}:
	Let $L \subseteq \Sigma^+$ be non-empty. Suppose~$L$ is not comma-free.
	Then there exist $u,v,w \in L$ such that
    $uv = xwy$ for some $x,y \in \Sigma^+$.
	Then $|x| \ge |u|$ or $|y| \ge |v|$ 
	implies $w \in \factor{v}$ or $w \in \factor{v}$, respectively, so that $L$
	is not an infix code.
	For $|x| < |u|$ and $|y| < |v|$, we have 
	$w = u'v'$ for a proper non-empty suffix $u'$ of $u$
	and a proper non-empty prefix~$v'$ of~$v$,
	so that $L$ is not overlap-free.
	Hence, if $L$ is a solid code, we must have $L^2 \cap \Sigma^+L\Sigma^+ = \emptyset$.
	
	\ref{constr_sync:prop:item:cf-in-ic}: Let $C \subseteq \Sigma^+$ be a comma-free code.
	Suppose $u, vuw \in C$ for $u,v,w \in \Sigma^*$.
	For $w \ne \varepsilon$, we have $uvuw \in C^2 \cap \Sigma^+ u \Sigma^+$,
	and for $v \ne \varepsilon$, we have $vuwu \in C^2 \cap \Sigma^+u \Sigma^+$.
	Whence, by comma-freeness, we find $v = w = \varepsilon$. So $C$ is an infix code.
\end{proof}

\begin{figure}[ht]
	\centering
	\begin{tikzpicture}[
		node distance=1.7em,
		every node/.style={align=center},
		every path/.style={-}
		]
		
		\node (solid) {solid};
		\node (comma) [above=of solid] {comma-free};
		\node (infix) [above=of comma] {infix};
		\node (bifix) [above=of infix] {bifix};
		\node (prefix) [above left=of bifix] {prefix};
		\node (suffix) [above right=of bifix] {suffix};
		\node (general) [above=of bifix, yshift=10mm] {code};
		
		\node (overlap) [
		right=3em of comma,
		yshift=6mm
		] {overlap-free};
		
		\draw (solid) -- (comma);
		\draw (comma) -- (infix);
		\draw (infix) -- (bifix);
		
		\draw (bifix) -- (prefix);
		\draw (bifix) -- (suffix);
		
		\draw (prefix) -- (general);
		\draw (suffix) -- (general);
		
		\draw (solid) -- (overlap);		
	\end{tikzpicture}
	\caption[Relations between different code classes and overlap-free languages]{Inclusion relations between different code classes and overlap-free languages.}
	\label{fig:code-hasse}
\end{figure}

\IN{rooted tree} \index{root|see{rooted tree}}
\begin{remark} \label{syncmax:rem:root-tree-prefix-code}
	Prefix codes can be visualized as rooted trees.
	A \emph{rooted tree} is an undirected connected and acyclic graph
	with the selection of a distinguished node, called the \emph{root}.
	A prefix code $C \subseteq \Sigma^+$
	induces the rooted tree $(C,  \{ \{ u, ux \} \mid u,ux \in \pref{C} \})$ with root $\varepsilon$ and 
	where each node has outdegree at most~$|\Sigma|$.
	For example, $\{00,01,10,110,111\}$ from \Cref{fig:intro:encoding}
	is a prefix code with its corresponding
	rooted tree shown (with added directions).
	Conversely, by labeling the outgoing edges 
	of each node of a rooted tree where each node has outdegree at most~$|\Sigma|$
	by a distinct letter from~$\Sigma$
	we uniquely obtain a prefix code.
	In this sense, up to a reordering of the alphabet,
	prefix codes over~$\Sigma$
	correspond uniquely to rooted trees with outdegree at most $|\Sigma|$.
\end{remark}

\IN{intercode of index one}
Comma-free codes are also known as \emph{intercodes of index one}~\cite[Thm.~3]{DBLP:journals/iandc/Jurgensen08}.\footnote{There appear to be two typographical errors in~\cite{DBLP:journals/iandc/Jurgensen08}. The definition of intercodes of index one omits the exponent~$2$, i.e., it states $L \cap \Sigma^+ L \Sigma^+ = \emptyset$ instead of $L^2 \cap \Sigma^+ L \Sigma^+ = \emptyset$. Furthermore, the condition $u,w \ne \varepsilon$ appears to be missing from the definition of an overlap-free code. Without this condition, only $L=\emptyset$ satisfies the definition, for if $v \in L$, we have $\varepsilon \cdot v = v \cdot \varepsilon$, and so $v=\varepsilon$. But this is impossible because $L \subseteq \Sigma^+$.}
Originally, comma-free codes were defined to be subsets of $\Sigma^n$ for some $n \ge 1$~\cite{golomb_gordon_welch_1958}, but this restriction was subsequently dropped~\cite{DBLP:books/daglib/0025093,DBLP:journals/iandc/Jurgensen08,DBLP:journals/tcs/JurgensenS21}.

\IN{coding morphism}
\IN{channel alphabet}
\IN{source alphabet}
The (algebraic) approach regarding %(variable-length) 
codes as certain subsets of $\Sigma^*$ goes back to Schützenberger~\cite{DBLP:journals/tit/Schutzenberger56}.
Prior to this, the theory of codes had its origin in Shannon's theory of information~\cite{Shannon:1948:MTCa,Shannon:1948:MTCb}.
Both viewpoints are linked via \emph{coding morphisms}, i.e., injective morphisms $\varphi \colon \Gamma^* \to \Sigma^*$
such that $\applyfunc{\Gamma}{\varphi}$ is a code.
Here, $\Gamma$ is called the \emph{source alphabet} 
and $\Sigma$ the \emph{channel alphabet}. 
Intuitively, this  expresses the fact
that codes  may be used as communication tools via a correspondence (encoding)
between the letters  $a \in \Gamma$ of the source alphabet
by code words $\varphi(a)$ over the channel alphabet~$\Sigma$,
where injectivity guarantees that a message $\varphi(u)$, $u \in \Gamma^*$,
can be uniquely deciphered as $u$ by the receiver.
%the decoding of a will never lead to ambiguities. 
% used for transmission
% of a word over $\Gamma$.
%The unique factorization property 
%guarantees that the decoding will never lead to ambiguities. 
%The link is as follows: codes may be used as communication tools by defining a correspondence (encoding)
%between the letters of a ``source'' alphabet $\Gamma$ 
%and the words of a code $C \subseteq \Sigma^+$ over the ``channel'' alphabet~$\Sigma$. The unique factorization property 
%guarantees that the decoding will never lead to ambiguities. 
%This is formalized by \emph{coding morphisms}: injective morphisms $\varphi \colon \Gamma^* \to \Sigma^*$ such that $\applyfunc{\Gamma}{\varphi}$ is a code where $\Gamma$ is called the \emph{source alphabet} 
%and $\Sigma$ the \emph{channel alphabet}.
For instance, in the introductory
example on error-prone communication from \cref{intro:sec:examples:communication}, 
the source alphabet is $\{a,h,l,o,r\}$
and the channel alphabet is $\{0,1\}$.
%For reference we state the following relation between coding morphisms and codes~\cite[Prop~2.1.1]{DBLP:books/daglib/0025093}.
The following statement~\cite[Prop~2.1.1]{DBLP:books/daglib/0025093} summarizes the relation of codes to coding morphisms.

\begin{proposition}[{\cite[Prop.~2.1.1]{DBLP:books/daglib/0025093}}]
	\label{constr_sync:prop:coding_morphisms}
	To every code $C \subseteq \Sigma^+$ there exists a coding morphism $\varphi \colon \Gamma^* \to \Sigma^*$ with $\varphi(\Gamma) = C$.  
	Moreover,
	any bijection from some alphabet~$\Gamma$ (here possibly infinite) onto a code $C \subseteq \Sigma^+$ 
	extends (as in \cref{prelim:prop:UMP}) to a unique injective morphism from $\Gamma^*$ to $\Sigma^*$ with image~$C^*$. Conversely, for injective morphisms $\varphi \colon \Gamma^* \to \Sigma^*$, $\applyfunc{\Gamma}{\varphi}$
	is a code.	
	%	 = C$ for $C \subseteq \Sigma^+$, then~$C$ is a code.\todo{originalformulierung?}
\end{proposition}

%\todo{großschreiben Universal...?}
%By the universal mapping property (\cref{prelim:prop:UMP}) the extension 
%in \cref{constr_sync:prop:coding_morphisms} 
%is unique and we always have $\varphi(\Gamma^*) = C^*$.

Throughout this work, our alphabet is always finite,
the only exception being \cref{prelim:prop:UMP,constr_sync:prop:coding_morphisms} 
where the (source) alphabets can be infinite.

%Note that the source alphabet might be infinite in this context.

\begin{remark}
	Combining these notions, we find that $L \subseteq \Sigma^+$
	is a code if and only if the map $\hat{i} \colon L^* \to \Sigma^*$
	from \cref{prelim:sec:formallang:basics} is injective, i.e.,
	a coding morphism where~$L$ is considered as a (possibly infinite)
	alphabet.
	% \hat{i} factors through any Gamma* to \Sigma* where Gamma and L bijective. Damit folgt vorherige Prop.
	Thus we can identify words formed by concatenating words from a code~$C$ with tuples
	over~$C$ in a unique way. This yields that \cref{prelim:prop:UMP}
	transfers to codes, i.e., every function from $C$ into a monoid~$M$
	extends uniquely into a monoid morphism from~$C^*$ into~$M$.	
\end{remark}

The next two results are well-known in combinatorics on words. \Cref{constr_sync:lem:two-words-code}
characterizes codes of size two and \cref{constrsync:prop:conjugate-words} gives a factorization result for triples of words $u,v,w$
such that $uv = wu$. 
For a proof of the former result see~\cite[Cor.~1.2.6, Prop.~1.3.1 \& 1.3.2]{lothaire_1997}
or~\cite[Lem.~5.1]{Ginsburg1964BOUNDEDAL} or~\cite[Lem.~3 \& Cor.~4.1]{LyndonSchutzenberger1962} for the latter three equivalences,
and for the latter result see~\cite[Prop.~1.3.4]{lothaire_1997}
or~\cite[Lem.~2]{LyndonSchutzenberger1962}.

\begin{proposition}
	\label{constr_sync:lem:two-words-code}
	\label{constr_sync:lem:common_power}
	Let $u,v \in \Sigma^*$ be distinct. Then the following are equivalent.
	\begin{enumerate}
		\item $\{u,v\}$ is not a code.
		
		\item $u,v \in w^*$ for some $w^*$.
		
		\item $u^i = v^j$ for some $i,j \ge 0$.
		
		\item $uv = vu$.
	\end{enumerate}	
\end{proposition}

\begin{remark}
	By \cref{constr_sync:lem:common_power}, with \cref{constrsync:ex:comma-free-single-word} from \cref{constrsync:ex:binarysolidcode}, 
	$\{w\}$ for $w \in \Sigma^+$ is comma-free if and only if $w = v^i$ for $i \ge 1$
	implies $i = 1$, 
\end{remark}

% Beweis von einem Teil:
% uv = vu
% u = u_1 .... u_n
% Angenommen u_1 = u_{n - m + 1}, ..., u_m = u_n und u_{m+1} = u_1, ..., u_n = u_{n - m} für ein m.
% Dann potenz von u_1 ... u_m. Folgt ganz einfach, da man die "Gleichungen" einfach nacheinander liest
\begin{proposition} \label{constrsync:prop:conjugate-words}
	\label{constrsync:lem:schuetz-lyndon}
	Let $u,v,w \in \Sigma^*$
	with $w \ne \varepsilon$ such that $uv = wu$.
	Then $w = xy$, $u = (xy)^kx$, and $v = yx$
	for some $x,y \in \Sigma^*$ and $k \ge 0$.		
\end{proposition}

\subsection{More on Prefix Codes}
\label{constrsync:sec:more-on-prefix-codes}

Here we give further properties of prefix codes. We begin with another characterization in \cref{constrsync:prop:char-prefix-code}. \Cref{constrsync:prop:suffix-factorizations} gives additional conditions equivalent to part of this characterization. As pointed out in \cref{constrsync:rem:prefix-codes-general-factorization}, these statements characterize prefix codes in a way that generalizes the defining property of codes by including certain suffixes in addition to code words. We then derive several consequences that will be needed later. Finally, in \cref{constr_sync:prop:aut_prefix_free}, we characterize regular prefix codes via PDFAs.

\begin{proposition} \label{constrsync:prop:char-prefix-code}
	A language $C \subseteq \Sigma^+$ is a prefix code if and only if it is a code
	such that $C^* \cap C^*u \ne \emptyset$ for $u \in \Sigma^*$
	implies $u \in C^*$.
\end{proposition}
\begin{proof}
 First, suppose~$C$ is a prefix code.
 Then it is a code.
 The remaining claim is equivalent to the fact that $vu \in C^*$ for $v \in C^*$
 implies $u \in C^*$.
 We show this by induction on the length of a product of~$v$ into words from~$C$.
 For $v = \varepsilon$, the claim clearly holds.
 Otherwise, we have $v = xx'$ and $vu = yy'$ with $x,y \in C$
 and $x',y' \in C^*$.
 Since~$C$ is a prefix code and $yy' = vu = xx'u$,
 we find $y = x$ and $y' = x'u$.
 Since $y',x' \in C^*$, inductively, we find $u \in C^*$.
 Thus, by induction, we obtain the claim. 
 
 Next, suppose that $C$ is not a prefix code. Then there exist $u \in C$ and a non-empty word $v \in \Sigma^+$ such that $uv \in C$. Thus either $C$ is not a code, or, if $C$ is a code, necessarily $v \notin C^+$. Indeed, otherwise the word $uv \in C$ would have the alternative factorization
 $
 uv = u \cdot v
 $
with $u \in C$ and $v \in C^+$, contradicting the unique factorization property of a code.
\end{proof}

The latter condition in \cref{constrsync:prop:char-prefix-code} is equivalent to two other conditions concerning the uniqueness of certain suffixes, as stated below. Notably, these equivalences hold for arbitrary languages in $\Sigma^+$, 
not just codes. For example, they apply to~$(\Sigma^i)^+$ for $i \ge 1$ and $\{ab,b\}$, although neither of these languages is a code.

\begin{proposition}  \label{constrsync:prop:suffix-factorizations}
	Let $L \subseteq \Sigma^+$ be any language not containing the empty word. Then the following are equivalent.
	\begin{enumerate}
		\item \label{constrsync:prop:suffix-factorizations:1} For $u,v \in \Sigma^* \setminus L\Sigma^*$ with $L^* u \cap L^* v \ne \emptyset$ we have $u = v$. 
		
		\item \label{constrsync:prop:suffix-factorizations:2} $L^* \cap L^* v \ne \emptyset$ for $v \in \Sigma^*$ implies $v \in L^*$.
		
		\item \label{constrsync:prop:suffix-factorizations:3} For $u,v \in \Sigma^+ \setminus L\Sigma^+$ with $L^* u \cap L^* v \ne \emptyset$ we have $u = v$.
	\end{enumerate}
\end{proposition}
\begin{proof}
	$\ref{constrsync:prop:suffix-factorizations:1} \Rightarrow \ref{constrsync:prop:suffix-factorizations:2}$: 
	Let $u, uv \in L^*$. Since $\varepsilon \notin L$, we can write $v = xy$
	with $x \in L^*$ and $y \notin L\Sigma^*$.
	Then we have $uv \in L^* \cap L^*y$. By \cref{constrsync:prop:suffix-factorizations:1}, we find $y = \varepsilon$
	and thus $v \in L^*$.
	
	$\ref{constrsync:prop:suffix-factorizations:2} \Rightarrow \ref{constrsync:prop:suffix-factorizations:3}$: 
	Let $w,w' \in L^*$ such that $wu = w'v$ with $u,v \in \Sigma^+ \setminus L\Sigma^+$.
	Assume $u \ne v$. Then, without loss of generality,
	we have $u = xv$ for some non-empty $x \in \Sigma^+$.
	This gives $w' = wx$. By assumption, we find $x \in L^+$
	and thus $u \in L\Sigma^+$, but $u \notin L \Sigma^+$ by choice of~$u$.
	Thus, we must have $u = v$.
	
	$\ref{constrsync:prop:suffix-factorizations:3} \Rightarrow \ref{constrsync:prop:suffix-factorizations:1}$: Suppose \cref{constrsync:prop:suffix-factorizations:3} holds true.
	Let $L^* u \cap L^* v \ne \emptyset$ for some $u,v \notin L\Sigma^*$.
	Then there exists $w \in L^*$
	such that $wu \in L^*v$.
	Assume $u \ne v$. 
	Then, by \cref{constrsync:prop:suffix-factorizations:3},
	we have $u,v \notin \Sigma^+ \setminus L \Sigma^+$.
	Since $u,v \notin L \Sigma^+$,
	this yields $u = \varepsilon$ or $v = \varepsilon$.
	Without loss of generality, let $u = \varepsilon$.	
	Write $w = xx'$ and $x'x' = yy'$ with $x$ and $y$ being the longest proper prefixes of~$w$ and $x'x'$, respectively, in~$L^*$. Then we have 
	\[
	x', y' \in \Sigma^+ \setminus L\Sigma^+
	\text{ and }
	wx' \in L^* x' \cap L^* y'.
	\]
	By assumption, this gives $x' = y'$.
	Since $x'x' = yy'$ we find $x' = y$.
	Thus $x' \in L^*$ since $y \in L^*$. With $x' \notin L\Sigma^+$ and $|x| > 0$,
	we get $x' \in L$.	
	We have $w \in L^* x' \cap L^* v$ where $v \in \Sigma^+$ since $u \ne v$.
	By assumption, we find $x' = v$, contradicting $v \notin L \Sigma^*$ since $x' \in L$.
	Thus, we must have $v = \varepsilon$.
\end{proof}

Note that $u$ and~$v$ must be non-empty in \cref{constrsync:prop:suffix-factorizations:3} of
\cref{constrsync:prop:suffix-factorizations}.
For example, consider any non-empty language $L \subseteq \Sigma^+$.
Then for $w \in L$ we have $w \in L^*\varepsilon \cap L^*w$,
but $w \ne \varepsilon$ since~$L$ is a subset of~$\Sigma^+$.

Combining \cref{constrsync:prop:char-prefix-code,constrsync:prop:suffix-factorizations},
we obtain that prefix codes extend the unique factorization property of codes
in the following sense.

\begin{corollary}  \label{constrsync:rem:prefix-codes-general-factorization}
	 \label{contrsync:prop:prefix-code-suffix-unique}
	A language $C \subseteq \Sigma^+$ is a prefix code if and only if
	for $u_1, \ldots, u_{n-1} \in C$, $v_1, \ldots, v_{m-1} \in C$,
	and $u_n, v_m \in \Sigma^+ \setminus C\Sigma^+$
	such that
	\[
	u_1 \cdots u_{n-1} u_n = v_1 \cdots v_{m-1} v_m
	\]
	we have $n = m$ and $u_i = v_i$ for all $i \in \{1,\ldots,n\}$.
\end{corollary}
\begin{proof}
	Suppose~$C$ is a prefix code and let $u_1, \ldots, u_n$
	and $v_1, \ldots, v_m$ as in the statement
	with $u_1 \cdots u_{n-1} u_n = v_1 \cdots v_{m-1} v_m$.
	Thus $C^* u_n \cap C^* v_m \ne \emptyset$.
	By \cref{constrsync:prop:char-prefix-code,constrsync:prop:suffix-factorizations}
	this gives $u_n = v_m$.
	Thus $u_1 \cdots u_{n-1} = v_1 \cdots v_{m-1}$.
	Since $C$ is a code, we find $n = m$ and $u_1 = v_1, \ldots, u_{n-1} = v_{n-1}$.	
	
	Conversely, this condition yields that~$C$ is a code
	since it generalizes the unique factorization property.
	Next, suppose $C^* u \cap C^* v$ for some $u,v \in \Sigma^+ \setminus C\Sigma^+$.
	Then there exist $u_1, \ldots, u_n \in C$
	and $v_1, \ldots, v_m \in C$
	such that $u_1 \cdots u_n u = v_1 \cdots v_m v$.
	By assumption, we find $n = m$, $u_i = v_i$ for $i \in \{1,\ldots,n\}$,
	and $u = v$.
	Thus, by \cref{constrsync:prop:char-prefix-code,constrsync:prop:suffix-factorizations},
	we conclude that~$C$ is a prefix code.
\end{proof}

Note that for a prefix code, we have
$C \subseteq \Sigma^+ \setminus C\Sigma^+$, i.e., the condition from
\cref{constrsync:rem:prefix-codes-general-factorization}
also allows $u_n$ and $v_m$ to belong to~$C$.

With our result, a variant of
\cref{constrsync:rem:prefix-codes-general-factorization}
can also be formulated with
$u_n, v_m \in \Sigma^* \setminus C\Sigma^*$
instead of
$u_n, v_m \in \Sigma^+ \setminus C\Sigma^+$.

\begin{corollary}
	Let $C \subseteq \Sigma^+$ be a prefix code.
	Then the factorization witnessing a word to
	be in $\pref{C^*} = C^*( \pref{C} \setminus C )$
	is unique, i.e., 
	for $u \in \pref{C^*}$
	there exists precisely one factorization
	of the form $u = u_1 \cdots u_{n-1} u_n$
	with $u_1, \ldots, u_{n-1} \in C$
	and $u_n \in \pref{C}\setminus C$.	
\end{corollary}
\begin{proof}
	This follows directly from \cref{contrsync:prop:prefix-code-suffix-unique}
	since $\pref{C} \cap C\Sigma^+ = \emptyset$ for a prefix code~$C$.
\end{proof}





\begin{corollary} \label{contrsync:lem:prefix-code-uv}
	Let $C \subseteq \Sigma^+$ be a prefix code.
	Then $\pref{C^*} \cap C^*u \ne \emptyset$ for $u \in \Sigma^*$
	implies $u \in \pref{C^*}$.	
\end{corollary}
\begin{proof}
	Let $vu \in \pref{C^*}$ with $v \in C^*$.	
	Write $vu = xx'$ and $u = yy'$ with $x \in C^*$ and $x'$ being a proper prefix of some word in~$C$
	and $y$ being a longest proper prefix of~$u$ in~$C^*$.
	Since~$C$ is a prefix code, we have $x' \notin C\Sigma^+$
	and, by choice of~$y$, we have $y' \in \Sigma^+ \setminus C\Sigma^+$.
	For $x' = \varepsilon$, we have $vu \in C^*$
	and by \cref{constrsync:prop:char-prefix-code} we find $u \in C^*$.
	Next, suppose $x' \ne \varepsilon$.
	We have $xx' = vu = vyy'$ and so $vu \in C^*x' \cap C^*y'$
	and by \cref{constrsync:prop:char-prefix-code,constrsync:prop:suffix-factorizations} we find $x' = y'$.
	Since $x' \in \pref{C^*}$,
	this gives $u = yx' \in C^*\pref{C^*} = \pref{C^*}$.	
\end{proof}

\Cref{contrsync:lem:prefix-code-uv} fails for general codes. For example, for $C = \{a,ab\}$
we have $ab \in \pref{C^*}$, $a \in C$, but $b \notin \pref{C^*}$.

\begin{corollary}
	Let $C \subseteq \Sigma^+$ be a prefix code and $u \in \Sigma^*$.
	Suppose $C^*u \cap C^* ( \pref{C} \setminus C ) \ne \emptyset$
	and that~$u$ is not a proper prefix of some word from~$C$ and $u \notin C\Sigma^+$.
	Then we have $u \in C$.
\end{corollary}
\begin{proof}
	Since $C^*w \subseteq \pref{C^*}$,
	we have $C^* u \cap \pref{C^*} \ne \emptyset$.
	Thus, by \cref{contrsync:lem:prefix-code-uv}, 
	we find $u\in \pref{C^*}$.
	Since $u \notin C\Sigma^+$, this gives $u \in \pref{C}$.
	Then, since $u \notin \pref{C}\setminus C$,
	we must have $u \in C$.
\end{proof}

\begin{corollary}
	Let $C \subseteq \Sigma^+$ be a prefix code and $L \subseteq \Sigma^*$
	such that $C^* L \subseteq L$.
	Then we have the following:
	\begin{align}  
		C^* ( L \setminus \pref{C^*} ) & = L \setminus \pref{C^*} \label{constrsync:eqn:complementprefCstar} \\	
		( \pref{C} \setminus C ) \Sigma	\cap \pref{C^*} & = ( \pref{C} \setminus C ) \Sigma	\cap \pref{C}, \label{constrsync:eqn:capprefCstar} 
	\end{align}
\end{corollary}
\begin{proof}
	\Cref{constrsync:eqn:complementprefCstar}: Since $\varepsilon \in C^*$
	we have $L \setminus \pref{C^*} \subseteq 	C^* ( L \setminus \pref{C^*} )$.
	Next, consider $u \in C^*$
	and $v \in L \setminus \pref{C^*}$.
	Now, \cref{contrsync:lem:prefix-code-uv} 
	would yield $v \in \pref{C^*}$
	for $uv \in \pref{C^*}$.
	Thus, we must have $uv \in  L \setminus \pref{C^*}$.
	
	\Cref{constrsync:eqn:capprefCstar}:	
	Let $ux = v_1 \cdots v_n$ for $u \in \pref{C} \setminus C$, $x \in \Sigma$,
	and $v_1, \ldots, v_{n-1} \in C$
	and $v_n \in \pref{C}$.
	We have $ux \notin C\Sigma^+$.
	Thus $n = 1$ and so $ux \in \pref{C}$.
\end{proof}

Lastly, we collect some facts on PDFAs accepting prefix codes. Related statements to Item~\ref{constr_sync:prop:aut_prefix_free:item:trim} and~\ref{constr_sync:prop:aut_prefix_free:item:min}
from \ref{constr_sync:prop:aut_prefix_free} have been given in~\cite[Prop.~2.2]{Eilenberg1974}.

\begin{proposition}
	\label{constr_sync:prop:aut_prefix_free}
	\label{constr_sync:prop:aut_prefix_free_exact_sync}
	Let $L \subseteq \Sigma^+$ be a regular language.
	
	\begin{enumerate}
		\item \label{constr_sync:prop:aut_prefix_free:item:trim} $L$ is a prefix code if and only if for every trim PDFA $\automat{B}{P}{\Sigma}{\mu}{p_0}{E}$
		with $L = L(\mathcal B)$ the states in $E$ have no outgoing transitions, i.e., $\mu(p, x)$ is undefined
		for each $p \in E$ and $x \in \Sigma$.
		
		% det und PDFA equal wenn maximal, sonst determinstic eien non-final trap zustand
		\item  \label{constr_sync:prop:aut_prefix_free:item:min} If $L$ is a prefix code, then the complete and deterministic minimal automaton $\minaut{L}$
		and a PDFA with the least number of states accepting $L$
		both have a single final state.
		
		\item  \label{constr_sync:prop:aut_prefix_free:item:sync} If $L$ is a prefix code and $\automat{B}{P}{\Sigma}{\mu}{p_0}{\{e\}}$ a  trim PDFA 
		with $L = L(\mathcal B)$, 
		then there exists an exactly synchronizing word $w \in L(\mathcal B)$
		with $\mu(P, w) = \{e\}$.
	\end{enumerate}
\end{proposition} 
\begin{proof}
	\ref{constr_sync:prop:aut_prefix_free:item:trim}: Let $\automat{B}{P}{\Sigma}{\mu}{p_0}{E}$ be a trim PDFA with $L = L(\mathcal B)$.
	First, suppose $\mu(p, x)$ is defined for some $p \in E$ and $x \in \Sigma$.
	As $\mathcal B$ is trim, 
	there exist $u,v \in \Sigma^*$
	with $\mu(\mu(p, x), u) \in E$ and $\mu(p_0, v) = p$.
	As $v \in L$ is a proper prefix of $vxu \in L$, $L$ is not a prefix code.
	Conversely, suppose $L$ is not a prefix code and $L = L(\mathcal B)$.
	Then there exist $u,v \in L$
	such that $v = uxu'$ for some $x \in \Sigma$ and $u' \in \Sigma^*$.
	Then $\mu(p_0, u) \in E$ and $\mu(p_0, ux)$ is defined.
	
	\ref{constr_sync:prop:aut_prefix_free:item:min}: 
	For $u,v \in L$ and $w \in \Sigma^+$
	we have $uw, vw \notin L$ since~$L$ is a prefix code,
	which gives $u \equiv_L v$. Therefore we obtain the claim
	for the complete and deterministic minimal automaton.
	The PDFA with the least number of states (up to renaming of the states)
	is obtained from the minimal automaton by omitting
	all non-coaccessible states (which form a set of size at most one), i.e.,
	by omitting the state $[u]_{\equiv_L}$ for $u \notin \pref{L}$
	if such a word $u$ exists. Formally, 
	let $\automat{B}{P}{\Sigma}{\mu}{p_0}{F}$
	where $P = \{ [u]_{\equiv_L} \mid u \in \pref{L} \}$, $p_0 = [\varepsilon]_{\equiv_L}$,
	$F = \{ [u]_{\equiv_L} \mid u \in L \}$,
	and $\mu([u]_{\equiv_L}, x) = [ux]_{\equiv_L}$ for $x \in \Sigma$ and $u \in \pref{L}$. Since~$L$ is regular, $P$ is finite.
	We have $L = L(\mathcal B)$ and $F$ is a singleton, which gives the claim.
		
	\ref{constr_sync:prop:aut_prefix_free:item:sync}: Set $n = |P|$. By Item~\ref{constr_sync:prop:aut_prefix_free:item:trim}
	the state $e$ has no outgoing transitions.
	Inductively, define a sequence of subsets of states $P_0, P_1, P_2, \ldots$, all containing the state $e$, % |P_i| \le n - i für i \ge n - 1
	by setting $P_0 = P$
	and, for $i > 0$, $P_i = \mu(P_{i-1}, u_i)$
	where $u_i$ is some word with $\mu(p, u_i) = e$ if there exists $p \in P_{i-1} \setminus \{e\}$
	and $u_i = \varepsilon$ otherwise.
	Then, for $|P_i| > 1$, as $u_i \ne \varepsilon$ yields that $\mu(e, u_i)$ is undefined,
	it is $|P_{i+1}| < |P_i|$.
	So $\mu(P, u_0 u_1 \cdots u_{n-1}) = \{e\}$ and there exist $p \in P$ and $v \in \Sigma^*$
	such that $\mu(p, u_0 u_1 \cdots u_{n-1}) = e$ and $\mu(p_0, v) = p$.
	Then for $w = vu_0 u_1 \cdots u_{n-1}$ we have $w \in L(\mathcal B)$ and $\mu(P, w) = \{e\}$.				
\end{proof}

\begin{example}\label{constrsync:ex:exact-sync}
	Consider the prefix code $C = \{b,ab,aab\}$.
	Then we obtain a trim PDFA $\mathcal B$ 
	accepting $C$ by setting $\mathcal B = (P, \{a,b\}, \mu, \varepsilon, \{e\})$ with $P = \pref{C}\setminus C \cup \{e\}$, i.e., the proper prefixes
	of~$C$ and~$e$ are considered as states of $\mathcal B$,
	and 
	\[
	\mu(u, x) = \begin{cases}
		e   & \text{if } ux \in C, \\
		ux  & \text{if } ux \in \pref{C} \setminus C.
	\end{cases}
	\]
	We have $\mu(P, b) = \{e\}$ (we also obtain this word by following the procedure from the proof of 
	\crefitem{constr_sync:prop:aut_prefix_free}{constr_sync:prop:aut_prefix_free:item:sync}).
	Other exactly synchronizing words are $aa$ and $aab$ with $\mu(P, aa) = \{aa\}$ and $\mu(P, aab) = \{e\}$.
\end{example}

\subsection{Relations to the Levenshtein-Romanov Mapping}
\label{constrsync:sec:codes:levensthein-romanov}

Here, we state characterizations and conditions for prefix, suffix, bifix, and infix codes
and overlap-free languages and combinations of these properties. We collect
several conditions here, but our focus is on those involving the Levenshtein-Romanov mapping.

\begin{proposition} \label{constrsync:prop:prefix-code}
	Let $L \subseteq \Sigma^+$ be non-empty. Then the following are equivalent.
	\begin{enumerate}
		\item \label{constrsync:prop:prefix-code:1} $L$ is a prefix code.
		
		\item \label{constrsync:prop:prefix-code:1a} $\pref{L} \setminus L$ equals the
		set of proper prefixes $\bigcup_{u \in L} ( \pref{u} \setminus \{u\} )$ of words from~$L$.
		
		\item \label{constrsync:prop:prefix-code:3} We have $\sigma_{\pref{L} \setminus \{u\}}(u) \ne u$ 
		for all $u \in L$.\footnote{With \cref{constrsync:eqn:sigma_prefL}, this is equivalent to $L$ being a singleton or, otherwise, to $\sigma_{L\setminus \{u\}}(u) \ne u$ for all $u \in L$.}
		
		\item \label{constrsync:prop:prefix-code:4} For all $u \in L$ and $v \in \Sigma^+$
		we have $\sigma_L(uv) \ne uv$.	
	\end{enumerate}
\end{proposition}
\begin{proof}
	$\ref{constrsync:prop:prefix-code:1} \Rightarrow \ref{constrsync:prop:prefix-code:1a}$:
	The set $\pref{L} \setminus L$ is always a subset of the set of proper prefixes of words from~$L$.
	Conversely, if~$u \in \pref{L}$ such that $uv \in L$ for some non-empty $v\in \Sigma^+$,
	then $u \notin L$ by the defining condition of a prefix code.	
		
	$\ref{constrsync:prop:prefix-code:1a} \Rightarrow \ref{constrsync:prop:prefix-code:3}$: 
	Let $u \in L$ with $\sigma_{\pref{L}\setminus{u}}(u) = u$.
	Then $u$ is a proper prefix of a word $v \in L$ with $u \ne v$.
	We have $u \notin \pref{L} \setminus L$
	and thus $\pref{L} \setminus L$ does not equal the set of proper prefixes.
		
	$\ref{constrsync:prop:prefix-code:3} \Rightarrow \ref{constrsync:prop:prefix-code:4}$: Suppose $\sigma_L(uv) = uv$ for $u \in L$
	and $v\in \Sigma^+$. Then we have $uv \in \pref{L}$
	where $uv \ne u$. Thus $L$ is not a singleton set 
	and we find $\sigma_{L\setminus\{u\}}(uv) = uv$.
	
	$\ref{constrsync:prop:prefix-code:4} \Rightarrow \ref{constrsync:prop:prefix-code:1}$: 
	If $\sigma_L(uv) = uv$
	for $u \in L$ and $v \in \Sigma^+$, then $uv\in \pref{L}$
	and thus $uvw \in L$ for some $w \in \Sigma^*$.
	Since $vw \ne \varepsilon$, $L$ is not a prefix code in this case.	
\end{proof}

\begin{proposition} \label{constrsync:prop:suffix-code}
	Let $L \subseteq \Sigma^+$ be non-empty. Then the following are equivalent.
	\begin{enumerate}
		\item \label{constrsync:prop:suffix-code:1} $L$ is a suffix code.\footnote{By \crefitem{constrsync:prop:suffsigmaB}{constrsync:prop:suffsigma:3B}, $L$ being a suffix code
		is also equivalent to $\suff{\sigma_{\pref{L}\setminus\{u\}}} \cap ( L \setminus \{u\} ) = \emptyset$
		for all $u \in L$.}
		
		\item \label{constrsync:prop:suffix-code:1a} $\suff{L} \setminus L$ equals the set of proper suffixes
		$\bigcup_{u \in L} ( \suff{u}\setminus\{u\} )$ of words from~$L$.
				
		\item \label{constrsync:prop:suffix-code:4} 
		%We have $uv \notin L$ for all $v \in L$ and $u \in \Sigma^+$.
		We have $\sigma_L(uv) \notin L \setminus\{v\}$ for all $v \in L$ and $u \in \Sigma^+$.
	\end{enumerate}
\end{proposition}
\begin{proof}
	$\ref{constrsync:prop:suffix-code:1} \Rightarrow \ref{constrsync:prop:suffix-code:1a}$: The set $\suff{L}\setminus L$
	is contained in the set of proper suffixes for any language. Conversely, if~$u$ is a proper suffix of a word from~$L$,
	then it is a suffix not contained in~$L$ if~$L$ is a suffix code, i.e., it is contained in $\suff{L}\setminus L$.
			
	$\ref{constrsync:prop:suffix-code:1a} \Rightarrow \ref{constrsync:prop:suffix-code:4}$: 
	Suppose $\sigma_L(uv) \in L \setminus \{v\}$ for some $v \in L$
	and $u \in \Sigma^+$. Then $\sigma_L(uv)$ is a word in~$L$
	having~$v$ as a proper suffix, but we have $v \notin \suff{L} \setminus L$ since $v \in L$.
		
	$\ref{constrsync:prop:suffix-code:4} \Rightarrow \ref{constrsync:prop:suffix-code:1}$: For $uv \in L$ with $v \in L$ and $u \in \Sigma^+$ we have $\sigma_L(uv) = uv \in L \setminus \{ v \}$.
\end{proof}

\begin{example}
	 $L = \{aaab,ab,aabb\}$ is not a suffix code and $\sigma_{L\setminus\{u\}}(u) \notin L$ for all $u \in L$.
%	\begin{enumerate}
%		\item $L = \{baa,aaab,a\}$ is not a suffix code
%		and we have $\sigma_{L\setminus\{u\}}(u) \notin L \setminus \{u\}$
%		for all $u \in L$.
%	\end{enumerate}
\end{example}

\begin{proposition} \label{constrsync:prop:bifix-code}
	Let $L \subseteq \Sigma^+$ be non-empty. Then the following are equivalent.
	\begin{enumerate}
		\item \label{constrsync:prop:bifix-code:1} $L$ is a bifix code.
		
		\item \label{constrsync:prop:bifix-code:2} We have $u,v \notin \suff{u}\cap \pref{v}$ for all $u,v \in L$ with $u\ne v$.
		
		\item \label{constrsync:prop:bifix-code:3} We have $\suff{\sigma_{\pref{L}\setminus\{u\}}(u)} \cap L = \emptyset$	
		for all $u \in L$.
	\end{enumerate}
\end{proposition}
\begin{proof}
$\ref{constrsync:prop:bifix-code:1} \Leftrightarrow \ref{constrsync:prop:bifix-code:2}$: 
For distinct $u,v \in L$ we have $u \in \suff{u} \cap \pref{v}$ (resp.~$v \in \suff{u} \cap \pref{v}$)
if and only if $u$ is a proper prefix of~$v$ (resp.~$v$ is a proper suffix of~$v$) in~$L$,
which gives the claim.

$\ref{constrsync:prop:bifix-code:2} \Leftrightarrow \ref{constrsync:prop:bifix-code:3}$: 
By \crefitem{constrsync:prop:suffsigmaB}{constrsync:prop:suffsigma:3B}, we have $v \in \suff{\sigma_{\pref{L}\setminus\{u\}}(u)} \cap L$ if and only if $v$ is a proper suffix of $u$
or $v = u$ and $u$ is a proper prefix of some word from~$L$,
which gives the claim.
\end{proof}

Clearly, $L \subseteq \Sigma^+$ is a prefix (resp.~suffix) code
if and only if $u \notin \suff{u} \cap \pref{v}$ (resp.~$u \notin \suff{v} \cap \pref{u}$) for distinct $u,v \in L$, 
giving similar characterizations as \cref{{constrsync:prop:bifix-code:2}} of \cref{constrsync:prop:bifix-code}.

\begin{proposition} \label{constrsync:prop:infix-code}
	Let $L \subseteq \Sigma^+$ be non-empty. Then the following are equivalent.
	\begin{enumerate}	
		\item \label{constrsync:prop:infix-code:1} $L$ is an infix code.	
		
		\item \label{constrsync:prop:infix-code:3} $L \Sigma^+ \cap \suff{L} = \Sigma^+ L \cap L = \emptyset$.
		
		\item \label{constrsync:prop:infix-code:4} $\Sigma^* L \Sigma^+ \cap L = \Sigma^+ L \Sigma^* \cap L = \emptyset$.
		
		\item \label{constrsync:prop:infix-code:5} $\Sigma^* L \Sigma^+ \cap L = \Sigma^+ L \cap L = \emptyset$.		
		
		% entsprechende chars wie die folgenden gehen auch für prefix und suffix codes, im fliestext danach
		\item \label{constrsync:prop:infix-code:2} $\Sigma^+ L \cap \pref{L} = L \Sigma^+ \cap L = \emptyset$.
		
		\item \label{constrsync:prop:infix-code:6} $L$ is a prefix code such that $\sigma_L(uv) = v$ for all $u \in \Sigma^*$ and $v \in L$.
			
		\item \label{constrsync:prop:infix-code:8} We have $uv \notin \pref{L}$
		and $vu \notin L$ for all $u \in \Sigma^+$ and $v \in L$.
				
		\item \label{constrsync:prop:infix-code:9} We have $\suff{\sigma_{\pref{L}\setminus\{u\}}(u)} \cap L = \emptyset$ for all $u \in \pref{L}$.
	\end{enumerate}
\end{proposition}
\begin{proof}
The equivalences from \cref{constrsync:prop:infix-code:1,constrsync:prop:infix-code:3,constrsync:prop:infix-code:4,constrsync:prop:infix-code:5,constrsync:prop:infix-code:2} are clear.
	
$\ref{constrsync:prop:infix-code:2}\Rightarrow \ref{constrsync:prop:infix-code:6}$:
The condition $\Sigma^+L \cap L = \emptyset$ implies that~$L$ is a prefix code.
Since $v \in L$ we have $\sigma_L(uv) \in \Sigma^*v \cap \pref{L}$.
By assumption, $\Sigma^+ v \cap \pref{L}$ is empty. Thus $\sigma_L(uv) = v$.

$\ref{constrsync:prop:infix-code:6}\Rightarrow \ref{constrsync:prop:infix-code:8}$:
Let $v \in L$ and $u \in \Sigma^+$.
For $uv \in \pref{L}$, by \cref{constr_sync:eqn_sigma_pref1} we have $\sigma_L(uv) = uv \ne v$.
For $vu \in L$, $L$ is not a prefix code.

$\ref{constrsync:prop:infix-code:8}\Rightarrow \ref{constrsync:prop:infix-code:9}$:
Let $v \in \suff{\sigma_{\pref{L}\setminus\{u\}}(u)} \cap L$.
By \cref{constrsync:prop:suffsigma:3B} of \cref{constrsync:prop:suffsigmaB}, 
$v$ is a proper suffix of $u$ or we have $u = v$ and $u$ is a proper prefix of some word from~$L$.
In the first case, we have $wv \in L$ for some $w \in \Sigma^+$
and in the latter case $vw \in L$ for some $w \in \Sigma^+$.
Thus, \cref{constrsync:prop:infix-code:8} does not hold.

$\ref{constrsync:prop:infix-code:9}\Rightarrow \ref{constrsync:prop:infix-code:1}$:
Let $v \in L$ and $uvw \in L$ with $u \ne \varepsilon$ or $w \ne \varepsilon$.
For $u = \varepsilon$, $v$ is a proper prefix of some word from~$L$.
Thus we have $v = \sigma_{\pref{L}\setminus\{v\}}(v)$ and $v \in \pref{L}$.
So,  \cref{constrsync:prop:infix-code:9}
does not hold.
Next, suppose $u \ne \varepsilon$.
Since $v \in L$ and $uv\ne v$
we have $\sigma_{\pref{L}\setminus\{uv\}}(uv) = w'v$ for some $w' \in \Sigma^*$.
Thus $v \in \suff{\sigma_{\pref{L}\setminus\{uv\}}(uv)} \cap L$.
Since $uv \in \pref{L}$, \cref{constrsync:prop:infix-code:9}
does not hold.
\end{proof}

\begin{example} 	
	For $L = \{a,ab\}$, we have $\sigma_L(uv) = v$ for all $u \in \Sigma^*$ and $v \in L$, but $L$ is not a prefix code. Hence, the assumption that $L$ is a prefix code in \crefitem{constrsync:prop:infix-code}{constrsync:prop:infix-code:6} is necessary.
\end{example}

Similar characterizations as in \cref{constrsync:prop:infix-code:2,constrsync:prop:infix-code:3,constrsync:prop:infix-code:4,constrsync:prop:infix-code:5} 
from \cref{constrsync:prop:infix-code} can also be easily given for prefix and suffix codes.

For $u,v \in \Sigma^*$, $\sigma_{\{v\}}(u)$ is the longest word $w$
such that $u \in w\Sigma^*$ and $v \in \Sigma^* w$. This suggests that we can relate
overlap-freeness to the Levenshtein-Romanov mapping. The next statements give some such relations.

\begin{proposition} \label{constrsync:prop:overlap-free}
	Let $L \subseteq \Sigma^+$ be non-empty. Then the following are equivalent.
	\begin{enumerate}
		\item \label{constrsync:prop:overlap-free:1} $L$ is overlap-free.
		
		\item \label{constrsync:prop:overlap-free:3} For all $u \in \Sigma^+$ we 
		have $u\Sigma^+\cap L = \emptyset$ or $\Sigma^+ u \cap L = \emptyset$.
		
		\item \label{constrsync:prop:overlap-free:2} For all $u,v \in L$ we 
		have $\suff{u} \cap \pref{v} \subseteq \{\varepsilon, u, v\}$.
		
		\item \label{constrsync:prop:overlap-free:2a} For all $u,v \in L$
		we have $\sigma_{\{v\}}(u) \in \{\varepsilon, u, v\}$ and $\suff{u}\cap \pref{u} \subseteq \{\varepsilon, u\}$.	
				
		\item \label{constrsync:prop:overlap-free:2aa} For all $u,v \in L$ we have $\sigma_{\pref{v}\setminus\{v\}}(u) \in \{\varepsilon, u\}$.
		
		\item \label{constrsync:prop:overlap-free:3a} We have $uv \notin L$ for all non-empty and proper prefixes $u,v$ of words from~$L$.
	\end{enumerate}
\end{proposition}
\begin{proof}
$\ref{constrsync:prop:overlap-free:1} \Rightarrow \ref{constrsync:prop:overlap-free:3}$:  Suppose we have $u\Sigma^+ \cap L \ne \emptyset$
and $\Sigma^+u \cap L \ne \emptyset$ for some $u \in \Sigma^+$.
Then $uv \in L$ and $wu \in L$ for some $v,w \in \Sigma^+$ and so~$L$ is not overlap-free.

$\ref{constrsync:prop:overlap-free:3} \Rightarrow \ref{constrsync:prop:overlap-free:2}$: 
Let $w \in \suff{u} \cap \pref{v}$
with $w \notin\{\varepsilon,u,w\}$ for $u,v \in L$.
Then $u = xw$ and $v = wy$ for some $x,y \in \Sigma^+$ and so $u \in \Sigma^+w$ and $v \in w\Sigma^+$.

$\ref{constrsync:prop:overlap-free:2} \Rightarrow \ref{constrsync:prop:overlap-free:2a}$: 
By \cref{constrsync:eqn:sigmaL-in-suffpref} and the assumption we have $\sigma_{\{v\}}(u) \in \suff{u} \cap \pref{v} \subseteq \{\varepsilon, u, v\}$.
From \cref{constrsync:prop:overlap-free:2} we also get $\suff{u}\cap \pref{u} \subseteq \{\varepsilon, u\}$.

$\ref{constrsync:prop:overlap-free:2a} \Rightarrow \ref{constrsync:prop:overlap-free:2aa}$: 
Suppose $\sigma_{\pref{v}\setminus\{v\}}(u) = w$ for some $w \notin \{\varepsilon, u\}$.
Then we have
\[
\sigma_{\{v\}}(u) = 
\begin{cases}
u & \text{if $u$ is a prefix of~$v$}, \\
w & \text{otherwise.}
\end{cases}
\]
In the first case, we have $w \in \suff{u}\cap \pref{u}$, in the second 
case we have $\sigma_{\{v\}}(u) \notin \{\varepsilon, u, v\}$.

$\ref{constrsync:prop:overlap-free:2aa} \Rightarrow \ref{constrsync:prop:overlap-free:3a}$:
Suppose we have $uv \in L$ for non-empty and proper prefixes $u,v$ of words from~$L$.
Then there exists $w \in \Sigma^+$ such that $vw \in L$.
We have $\sigma_{\pref{vw}\setminus\{vw\}}(uv) = xv$ for some $x \in \Sigma^*$.
By assumption, $x = u$. By \cref{constr_sync:eqn_sigma_pref1}, in this case, $uv$ is a prefix of $vw$.
Thus, $uv$ has $v$ as a proper prefix and 
we find that $\sigma_{\pref{uv}\setminus\{uv\}}(uv)$
has $v$ as a suffix.
Since $uv$ is not a prefix of a word from $\pref{uv}\setminus\{uv\}$,
we also have $\sigma_{\pref{uv}\setminus\{uv\}}(uv) \ne uv$.
Thus \cref{constrsync:prop:overlap-free:2aa} does not hold.

$\ref{constrsync:prop:overlap-free:3a} \Rightarrow \ref{constrsync:prop:overlap-free:1}$: 
 If~$L$ is not overlap-free, then $uv, vw \in L$ for some non-empty $u,v,w \in \Sigma^+$ where
 $u,v$ are non-empty and proper prefixes of words from~$L$.
 Thus, \cref{constrsync:prop:overlap-free:3a} does not hold.
\end{proof}

\begin{remark}
	Note that $\sigma_{\pref{v}\setminus\{v\}}(u) = \varepsilon$
	implies $\sigma_{\{v\}}(u) \in \{\varepsilon, v\}$,
	but the converse implication does not hold. For example, 
	we have $\sigma_{\{aa\}}(baa) = aa$, but $\sigma_{\{a\}}(baa) = a$.
\end{remark} 

\begin{corollary} \label{constrsync:rem:overlap-free}
	Let $u \in \Sigma^*$. Then the following are equivalent.
	\begin{enumerate}
	\item $\{u\}$ is overlap-free.
	
	\item $\sigma_{\pref{u}\setminus\{u\}}(u) = \varepsilon$.
	
	\item $\suff{u}\cap \pref{u} \subseteq \{\varepsilon, u\}$.
	\end{enumerate}
\end{corollary}

\begin{proposition} \label{constr_sync:cor:sigma_L_prefix-eps}
	Let $L \subseteq \Sigma^+$ and $u \in L$ where $L$ contains at least two words.
	Then we have $\sigma_{\pref{L} \setminus \{u\}}(u) = \varepsilon$
	if and only if $\{u\}$ is overlap-free and $\sigma_{L \setminus \{u\}}(u) = \varepsilon$.
\end{proposition}
\begin{proof}
	First, suppose $\sigma_{\pref{L} \setminus \{u\}}(u) = \varepsilon$.
	Since $u \in L$ we have $\pref{L}\setminus \{u\} = ( L \setminus \{u\} ) \cup ( \pref{u} \setminus \{u\})$.
	Thus, by \crefitem{constr_sync:prop:sigma_L_prefix}{constr_sync:prop:sigma_L_prefix:7},
	we find $\sigma_{L \setminus\{u\}}(u) = \varepsilon$
	and $\sigma_{\pref{u}\setminus\{u\}}(u) = \varepsilon$, and thus~$u$ is overlap-free by \cref{constrsync:rem:overlap-free}.
	Conversely, if $\{u\}$ is overlap-free and $\sigma_{L\setminus\{u\}}(u) = \varepsilon$,
	then by \cref{constrsync:rem:overlap-free,constrsync:eqn:sigma_prefL} 
	we obtain $\sigma_{\pref{L} \setminus \{u\}}(u) = \varepsilon$.
\end{proof}

\begin{example}
	Let $L = \{a, ab,ba\}$. Then $L$ is not overlap-free. 
	Note that here we have $uv \notin L$ for all $u,v \in \pref{L} \setminus L = \{\varepsilon, b\}$.
\end{example}

Suppose $uv, vw \in L$ for non-empty $u,v,w \in \Sigma^+$.
Then $uvw \in \left(\bigcup_{x \in L} \pref{x}\setminus\{x\} \right) L \cap L\Sigma^+$. Thus, if the latter
set is empty, then~$L$ is overlap-free. The following statement also shows that when~$L$
is a prefix, this condition characterizes overlap-freeness of~$L$.

\begin{proposition} \label{constrsync:prop:overlap-free-prefix}
Let $L \subseteq \Sigma^+$ be non-empty. Then the following are equivalent.
\begin{enumerate}		
\item \label{constrsync:prop:overlap-free-prefix:1xx} $L$ is an overlap-free prefix code.

\item \label{constrsync:prop:overlap-free-prefix:2xx} For all $u,v \in L$ we have $\suff{u} \cap \pref{v} \subseteq \{\varepsilon, v\}$.

\item  \label{constrsync:prop:overlap-free-prefix:2ax} For all $u,v \in L$ we have $\sigma_{\pref{u}\setminus\{u\}}(v) = \varepsilon$.		

\item  \label{constrsync:prop:overlap-free-prefix:2bx} For all $u \in L$, the language $\{u\}$ is overlap-free and 
$\sigma_{\{v\}}(u) \in \{\varepsilon, v\}$ for distinct $u,v \in L$.

\item \label{constrsync:prop:overlap-free-prefix:8}  For all $u \in L$, $\{u\}$ is overlap-free
and $\sigma_{L\setminus\{u\}}(u) \in (L \setminus\{u\}) \cup \{\varepsilon\}$
if $L \setminus\{u\} \ne \emptyset$.		

\item \label{constrsync:prop:overlap-free-prefix:5ax} For all $u \in L$ we have 
$\sigma_{\pref{L} \setminus \{u\}}(u) \in ( L \setminus \{u\} ) \cup \{\varepsilon\}$.

\item \label{constrsync:prop:overlap-free-prefix:3yx} 
We have $uv \notin L$ for all $u,v \in \Sigma^+$ such that $u \in L$ or $v \in \pref{L} \setminus L$.

\item \label{constrsync:prop:overlap-free-prefix:2axy} $L$ is a prefix code
such that $\left(\bigcup_{u \in L} \pref{u}\setminus\{u\} \right) L \cap L\Sigma^+ = \emptyset$.
\end{enumerate}
\end{proposition}
\begin{proof}	
$\ref{constrsync:prop:overlap-free-prefix:1xx} \Rightarrow \ref{constrsync:prop:overlap-free-prefix:2xx}$: 
This follows from \cref{constrsync:prop:overlap-free} and the fact that 
$u \notin \pref{v}$ for distinct $u,v \in L$ when~$L$ is a prefix code.
	
$\ref{constrsync:prop:overlap-free-prefix:2xx} \Rightarrow \ref{constrsync:prop:overlap-free-prefix:2ax}$: 
We have $\sigma_{\pref{u}\setminus\{u\}}(v) \in \suff{v} \cap ( \pref{u} \setminus\{u\} ) \subseteq \{\varepsilon\}$ by assumption and the definition of~$\sigma_L$.
	
$\ref{constrsync:prop:overlap-free-prefix:2ax} \Rightarrow \ref{constrsync:prop:overlap-free-prefix:2bx}$:
By \cref{constrsync:rem:overlap-free}, we find that $\{u\}$ is overlap-free for each $u \in L$.
Next, for $u,v \in \Sigma^*$ we have
\[
\sigma_{\{v\}}(u) = \begin{cases}
	\sigma_{\pref{v}\setminus\{v\}}(u) & \text{if } v \notin \suff{u}, \\
	v & \text{otherwise.}	
\end{cases}
\]
Thus, we get $\sigma_{\{v\}}(u) \in \{\varepsilon,v\}$
if $u,v \in L$ by \cref{constrsync:prop:overlap-free-prefix:2ax}.

$\ref{constrsync:prop:overlap-free-prefix:2bx}\Rightarrow\ref{constrsync:prop:overlap-free-prefix:8}$: We have $\sigma_{L \setminus \{u\}}(u) = \sigma_{\{v\}}(u)$
for some $v \in L$ with $v \ne u$. By assumption, $\sigma_{\{v\}}(u) \in \{\varepsilon,v\}$
and this gives the claim.

$\ref{constrsync:prop:overlap-free-prefix:8}\Rightarrow\ref{constrsync:prop:overlap-free-prefix:5ax}$: This follows from \cref{constrsync:eqn:sigma_prefL,constrsync:rem:overlap-free}.

$\ref{constrsync:prop:overlap-free-prefix:5ax} \Rightarrow \ref{constrsync:prop:overlap-free-prefix:3yx}$:
Let $u$ be a shortest non-empty word such that $uv \in L$
for some $v \in \Sigma^+$
with $u \in L$ or $v \in \pref{L}\setminus L$.
If $u \in L$,
then $\sigma_{\pref{L}\setminus\{u\}}(u) = u$.
Next, suppose $v \in \pref{L}\setminus L$.
Then $\sigma_{\pref{L}\setminus\{uv\}}(uv) = xv$
for some $x \in \Sigma^*$.
Assume $\sigma_{\pref{L}\setminus\{uv\}}(uv) \in L$.
Since $v \notin L$, this gives $x \ne \varepsilon$.
By the minimal choice of~$u$, we find $x = u$.
Thus we have $\sigma_{\pref{L}\setminus\{uv\}}(uv) = uv$.
This gives the claim.

$\ref{constrsync:prop:overlap-free-prefix:3yx} \Rightarrow \ref{constrsync:prop:overlap-free-prefix:2axy}$:
If $L$ is not a prefix code, then $uv \in L$
for some $u \in L$ and $v \in \Sigma^*$.
Next, suppose~$L$ is a prefix code
and $uv = wx$ for $u$ being a proper prefix of some word from~$L$
and where $v,w \in L$ and $x \in \Sigma^+$.
Since~$L$ is a prefix code,
we have $w = uy$ and $v = yx$ for some $y \in \Sigma^+$
where $y \notin L$. Then we have $uy \in L$ with $y \in \pref{L}\setminus L$.

$\ref{constrsync:prop:overlap-free-prefix:2axy} \Rightarrow 
\ref{constrsync:prop:overlap-free-prefix:1xx}$:
Suppose $uv \in L$ and $vw \in L$ for some $u,v,w \in \Sigma^+$.
Since~$L$ is a prefix code, we have $u \notin L$.
Thus 
$uvw \in \left(\bigcup_{u \in L} \pref{u}\setminus\{u\} \right) L \cap L\Sigma^+$.
\end{proof}

\begin{remark}
	Let $L = \{ aabb, bba, abb \}$ and $u = aabb$.
	Then we have $\sigma_{\pref{L} \setminus \{u\}}(u) = abb$
	where $abb \in L \setminus \{u\}$,
	but $\sigma_{\{bba\}}(u) = bb \notin \{\varepsilon, bba\}$.	
	Thus the equivalences of \cref{constrsync:prop:overlap-free-prefix:2bx} 
	and \cref{constrsync:prop:overlap-free-prefix:5ax}
	in \cref{constrsync:prop:overlap-free-prefix}
	do not hold for a fixed $u \in L$, but only when quantified
	over all~$u \in L$.	
\end{remark}

\begin{corollary}\label{lem:overlap-free-prefix}
	Let $C \subseteq \Sigma^+$ be a prefix code.
	Then $C$ is overlap-free if and only if $(\pref{C} \setminus C) C \cap C \Sigma^+ = \emptyset$.
\end{corollary} 
\begin{proof}
	Let~$C$ be a prefix code.
	Then its set of proper prefixes
	$\bigcup_{u \in L} \pref{u}\setminus\{u\}$
	equals $\pref{C}\setminus C$.
	Thus the claim follows from \cref{constrsync:prop:overlap-free-prefix}.
\end{proof}

\begin{remark}
	Let $L \subseteq \Sigma^*$ be non-empty and $u \in \Sigma^*$.
	We have
	\[
	\sigma_{\pref{L}\setminus \{u\}}(u) \notin \bigcup_{v \in L} ( \pref{v} \setminus \{\varepsilon,v\} )
	\Rightarrow
	\sigma_{\pref{L} \setminus \{u\}}(u) \in ( L \setminus \{u\} ) \cup \{\varepsilon\}.
	\]
	The converse implication holds if no proper suffix of~$u$ contained in~$L$
	is a proper prefix of some other word from~$L$ (which holds if~$L$ is a prefix code).
	This yields that \cref{constrsync:prop:overlap-free-prefix:5ax} 
	from \cref{constrsync:prop:overlap-free-prefix}
	can also be stated by saying the image of the Levenshtein-Romanov mapping
	is not in $\bigcup_{v \in L} ( \pref{v} \setminus \{\varepsilon,v\} )$
	since either condition, when quantified over all words $u \in L$,
	yields that~$L$ is a prefix code.
\end{remark}

\begin{proposition} \label{constrsync:prop:overlap-suffix}
Let $L \subseteq \Sigma^+$. Then the following are equivalent.
\begin{enumerate}
\item \label{constrsync:prop:overlap-suffix:1} $L$ is an overlap-free suffix code.

\item \label{constrsync:prop:overlap-suffix:2} For all $u,v \in L$ we have $\suff{u} \cap \pref{v} \subseteq \{\varepsilon, u\}$.

\item \label{constrsync:prop:overlap-suffix:3} $\{u\}$ is overlap-free for $u \in L$ and $\sigma_{\{v\}}(u) \in \{\varepsilon, u\}$ for $u,v \in L$.

\item \label{constrsync:prop:overlap-suffix:4} $uv \notin L$ for all $u,v \in \pref{L} \setminus \{\varepsilon\}$.	
\end{enumerate}
\end{proposition}
\begin{proof}
$\ref{constrsync:prop:overlap-suffix:1} \Leftrightarrow \ref{constrsync:prop:overlap-suffix:2}$: This follows from \cref{{constrsync:prop:overlap-free:2}} of \cref{constrsync:prop:overlap-free}
and the fact that for distinct $u,v \in L$ we have $v \notin \suff{u}$ if and only if~$L$ is a suffix code.

$\ref{constrsync:prop:overlap-suffix:1} \Leftrightarrow \ref{constrsync:prop:overlap-suffix:3}$:
This follows from \cref{{constrsync:prop:overlap-free:2a}} of \cref{constrsync:prop:overlap-free}
and \cref{constr_sync:eqn_sigma_suff0}.

$\ref{constrsync:prop:overlap-suffix:1} \Leftrightarrow \ref{constrsync:prop:overlap-suffix:4}$: This follows from \cref{constrsync:prop:overlap-free:3a} of
\cref{constrsync:prop:overlap-free}
and the definition of a suffix code.	
\end{proof}

\begin{proposition} \label{constrsync:prop:overlap-free-bifix}
	Let $L \subseteq \Sigma^+$ be non-empty. Then the following are equivalent.	
	\begin{enumerate}
		
		\item \label{constrsync:prop:overlap-free-bifix:1} $L$ is an overlap-free bifix code.
		
		\item \label{constrsync:prop:overlap-free-bifix:2} For $u,v \in L$ we have
		\[
		\suff{u} \cap \pref{v} = \begin{cases}
			\{\varepsilon\} & \text{if } u \ne v, \\
			\{\varepsilon, u\} & \text{if } u = v.	
		\end{cases}
		\]
		
		\item \label{constrsync:prop:overlap-free-bifix:3} For $u \in L$ we have $\sigma_{\pref{L}\setminus\{u\}}(u) = \varepsilon$.\footnote{As stated
			in \cref{constr_sync:cor:sigma_L_prefix-eps},
			this is equivalent to $\{u\}$ being overlap-free
			and $\sigma_{L\setminus\{u\}}(u) = \varepsilon$
			if $L \setminus \{u\} \ne \emptyset$.}

		\item \label{constrsync:prop:overlap-free-bifix:4} 
		We have $uv \notin L$ for all $u,v \in \Sigma^+$ with $u \in L$ or $v \in \pref{L}$.
	\end{enumerate}
\end{proposition}
\begin{proof}
$\ref{constrsync:prop:overlap-free-bifix:1} \Leftrightarrow \ref{constrsync:prop:overlap-free-bifix:2}$: This follows from \cref{constrsync:prop:overlap-free,constrsync:prop:bifix-code}.

$\ref{constrsync:prop:overlap-free-bifix:1} \Leftrightarrow \ref{constrsync:prop:overlap-free-bifix:3}$:
First, suppose $L$ is an overlap-free bifix code.
By \crefitem{constrsync:prop:overlap-free-prefix}{constrsync:prop:overlap-free-prefix:5ax}
we have that $\sigma_{\pref{L}\setminus\{u\}}(u)$ is in $( L \setminus \{u\} ) \cup \{\varepsilon\}$.
Since $\sigma_{\pref{L}\setminus\{u\}}(u)$ is a suffix of $u$, 
it cannot be contained in $L \setminus \{u\}$ since~$L$ is a suffix code.
Conversely, by \crefitem{constrsync:prop:overlap-free-prefix}{constrsync:prop:overlap-free-prefix:5ax}
this condition implies that~$L$ is an overlap-free prefix code.
Next, if $uv \in L$ for $u \in \Sigma^+$ and $v \in L$,
then $\sigma_{\pref{L}\setminus\{uv\}}(uv)$ has the non-empty word $v$ as a suffix.
Thus, $L$ must be a suffix code.

$\ref{constrsync:prop:overlap-free-bifix:1} \Leftrightarrow \ref{constrsync:prop:overlap-free-bifix:4}$: 
This follows from \crefitem{constrsync:prop:overlap-free-prefix}{constrsync:prop:overlap-free-prefix:3yx}
and the definition of a suffix code.
\end{proof}

Note that $L \subseteq \Sigma^+$ is a prefix (resp.~suffix, bifix, infix, overlap-free) code/language
if and only if $\{u,v\}$ is a prefix (resp.~suffix, bifix, infix, overlap-free) code/language
for all $u,v \in L$. So, the conditions from the above statements characterizing those language classes
can be written as conditions that must be true for all pairs (of possibly equal) words in $L$.

The following further characterizations were given in~\cite{DBLP:journals/tcs/JurgensenS21}
for non-empty $L \subseteq \Sigma^+$: 
\begin{enumerate}
	\item $\pref{L}\cap \Sigma^*L = L$ if and only if $\sigma_L(u) \in L$ for all $u \in \Sigma^*L$,
	
	\item $L$ is an overlap-free prefix code if and only if $\sigma_L(uv) = \sigma_L(v)$
	for all $u \in L$ and $v \in \Sigma^+$,
	
	\item $L$ is a suffix code such that $\pref{L} \cap \Sigma^*L = L$ if and only if $\sigma_L(uv) = v$
	for all $u \in \Sigma^+$ and $v \in L$ (this condition also appears in \crefitem{constrsync:prop:infix-code}{constrsync:prop:infix-code:6}).
\end{enumerate}

In the next section we will investigate overlap-free infix codes, i.e., solid codes, in more detail.

We always have $L \cup \{\varepsilon\} \subseteq \pref{L} \cap \suff{L}$
for any language $L \subseteq \Sigma^*$. Lastly, we connect
the converse inclusion to overlap-freeness.

\begin{proposition}  \label{constrsync:prop:pref-suff-overlap-free-inclusion-eqn}
	Let $C \subseteq \Sigma^+$ be a prefix or a suffix code.
	If $\pref{C} \cap \suff{C} \subseteq C \cup \{\varepsilon\}$,
	then $C$ is overlap-free.
\end{proposition}
\begin{proof}
	Let $uv, vw \in C$ for non-empty $u,v,w$.
	Then we have $v \in \pref{C} \cap \suff{C}$
	and since~$C$ is a prefix or suffix code, we have $v \notin C$, i.e.,
	$\pref{C} \cap \suff{C}$
	is not a subset of~$C$.
\end{proof}

\begin{proposition} \label{constrsync:prop:overlap-free-inclusion-eqn}
	Let $L \subseteq \Sigma^*$.
	If $L$ is overlap-free,
	then $\pref{L} \cap \suff{L} \subseteq L \cup \{\varepsilon\}$.
\end{proposition}
\begin{proof}
	Let $u \in \pref{L}\cap \suff{L}$ with $u \ne \varepsilon$.
	Then $uv, wu \in L$ for some $v,w \in \Sigma^*$.
	If~$L$ is overlap-free, then $v = \varepsilon$
	or $w = \varepsilon$ and each case gives $u \in L$,
	which shows the claim.	
	%Let $v \in \pref{L}\cap \suff{L}$ with $v \ne \varepsilon$ and $v \notin L$.
	%Then we have $uv, vw \in L$ for some non-empty $u,w \in \Sigma^+$, i.e., $L$ 
	%is not overlap-free.
\end{proof}

Thus, we find that a prefix or suffix code~$C$
is overlap-free if and only if $\pref{C} \cap \suff{C} \subseteq C \cup \{\varepsilon\}$.

\subsection{Characterizing Solid Codes}
\label{constrsync:sec:codes:solid}

%!TEX root = ./../../springer-mono-format/book.tex
%!TEX root = ./../../thesis-trier.tex
 
\IN{strongly regular codes}
Solid codes provide outstanding fault-tolerance when used for information transmission through a noisy channel w.r.t.~symbol substitutions and synchronization errors. In this context,
they were introduced by Levenshtein~\cite{Lev64} under the name \emph{strongly regular codes}, 
and further investigated by Romanov~\cite{Rom66}.
Shyr and Yu~\cite{ShyrYu90} introduced solid codes to investigate disjunctive domains, 
but did not recognize the connection to the earlier work by Levenshtein and Romanov. 
Later, Jürgensen and Yu~\cite{DBLP:journals/eik/JurgensenY90} observed that strongly 
regular codes describe precisely the solid codes.
For further information, we refer to the surveys~\cite{DBLP:journals/iandc/Jurgensen08,DBLP:reference/hfl/JurgensenK97},
or the introductory section of~\cite{DBLP:journals/tcs/JurgensenS21}. 
Recall that \cref{constrsync:ex:binarysolidcodeitem,constrsync:ex:constr-solid-code}
from \cref{constrsync:ex:binarysolidcode} gave an infinite binary solid code 
and a method to construct a solid codes, respectively.

In our investigation on the constrained synchronization problem and synchronization conditions
for automata obtained in reductions, we discovered (to the best of our knowledge)
new characterizations of solid codes, which we present in the following
(see the discussion section,
\cref{constr_sync:sec:gen_results:discussion}, page~\pageref{discussion:strongly-self-sync-codes}, 
for some background on how this characterization arose in the context
of the constrained synchronization problem).

Here, first, in \cref{thm:char-solid-codes}, we recap known characterizations. 
\Cref{thm:solidcode} collects further characterizations of solid codes
which are new and connect solid codes to various language equations.
Lastly, \cref{thm:solidcode-char2,cor:solidcode-char3,constrsync:prop:sigma-solid-code-char} 
give two other new characterizations.

\begin{theorem} \label{thm:char-solid-codes}
Let $L \subseteq \Sigma^+$ with $L \ne \emptyset$. Then the following are equivalent.
\begin{enumerate}
	
\item  \label{thm:char-solid-codes:1} $L$ is a solid code.
	
\item  \label{thm:char-solid-codes:2} (Jürgensen \& Yu~\cite{DBLP:journals/eik/JurgensenY90}) For every word $w \in \Sigma^*$ 
 we have $w = u_0 w_1 u_1 w_2 u_2 \cdots w_n u_n$
 for some $n \ge 0$ and unique words $u_0, u_1, \ldots, u_n \in \Sigma^* \setminus \Sigma^* L \Sigma^*$
 and $w_1, w_2, \ldots, w_n \in L$.
	
\item \label{thm:char-solid-codes:3} (Staiger~\cite{DBLP:journals/tcs/JurgensenS21}, Romanov~\cite{Rom66}, Levenshtein~\cite{Lev64})
 There exists a state-invariant decoder without look-ahead for $L$ (see \cite[Def.~9]{DBLP:journals/tcs/JurgensenS21}
 for a definition of these decoders for finite solid codes and extensions to infinite solid codes).

\item \label{thm:char-solid-codes:4} (Liu, Guo \& Shum~\cite[Thm.~2.4]{Liu2021}) $L$ is an infix code and $\{ u_1\cdots u_n \}$ is overlap-free for 
 all sequences of distinct words $u_1, \ldots, u_n \in L$ with $n \ge 1$.\footnote{In~\cite[Thm.~2.4]{Liu2021},
 	if~$L$ is finite, then overlap-freeness is only stated to hold for $\{u_1\cdots u_n\}$
 	if $2 \le n \le |L|$. However, this seems to be a typo,
 	as for example $L = \{ aba, bcb \}$
 	is an infix code
 	and $\{ ababcb\}$ and $\{ bcbaba \}$
 	are each overlap-free, while~$L$
 	is not overlap-free and so not a solid code.} 
 
\item \label{thm:char-solid-codes:5} (Liu, Guo \& Shum~\cite[Thm.~2.4]{Liu2021}) $L$ is an infix code such that $\{u\}$ is overlap-free for all $u \in L$ and $\{uv\}$ is overlap-free for $u,v \in L$ with $u \ne v$.
\end{enumerate}	
\end{theorem}

\begin{remark}
	Since for distinct $u_1, \ldots, u_n \in L$
	we have $u_i \cdots u_n \ne u_1 \cdots u_{n - i + 1}$
	for all $i \in \{1,\ldots,n\}$.
	Thus~$L$ is overlap-free
	if $\{ u_1 \cdots u_n \}$
	is overlap-free for distinct $u_1, \ldots, u_n \in L$.
	Next, if for distinct $u,v \in L$
	we have $u = xy$ and $v = yz$
	for non-empty $x,y,z \in \Sigma^+$,
	then $vu = yzxy$ and so $\{vu\}$
	is not overlap-free.
	Thus, if $\{u\}$ for $u \in L$ and $\{uv\}$
	for distinct $u,v \in L$ are overlap-free,
	then~$L$ is overlap-free.
	This shows that the additional conditions in Items~\ref{thm:char-solid-codes:4} and~\ref{thm:char-solid-codes:5} of \cref{thm:char-solid-codes}, beyond being an infix code, are equivalent to overlap-freeness (without any additional assumptions on~$L$).
\end{remark}

\begin{theorem} \label{thm:solidcode}
Let $L \subseteq \Sigma^+$. Then the following are equivalent:
\begin{enumerate}
\item  \label{thm:solidcode:1} $L$ is a solid code.

\item  \label{thm:solidcode:1a} $L$ is a bifix code and $(\pref{L} \setminus L) L \cap \Sigma^* L \Sigma^+ = \emptyset$.

\item  \label{thm:solidcode:1b}  $L$ is a suffix code and $\left( \bigcup_{u \in L} ( \pref{u}\setminus\{u\} ) \right)L \cap \Sigma^*L\Sigma^+ = \emptyset$.

\item \label{thm:solidcode:2} $L$ is an infix code and $(\pref{L} \setminus L) L \cap L \Sigma^+ = \emptyset$.

\item \label{thm:solidcode:3} $L$ is an infix code and $\pref{L} \cap \suff{L} = L \cup \{\varepsilon\}$.	
\end{enumerate}
\end{theorem}
\begin{proof}
$\ref{thm:solidcode:1}\Rightarrow \ref{thm:solidcode:1a}$: A solid code is an infix code, and so a bifix code, and overlap-free.
Suppose $vw = xuy$
with $x \in \Sigma^*$, $y \in \Sigma^+$, $u, w \in L$, and
$v \in \pref{L} \setminus L$.
By \cref{lem:overlap-free-prefix}
we have $|x| > 0$.
For $|x| \ge |v|$ (resp.~$|y| \ge |w|$), we have $u \in \factor{w}$ (resp.~$u \in \factor{v}$).
For $0 < |x| < |v|$ and $0 < |y| < |w|$,
we have $u = v'w'$ for a non-empty
and proper suffix $v'$ (resp.~prefix~$w'$)
of $v$ (resp.~$w$), so that $L$ is not overlap-free.
Hence, for $(\pref{L} \setminus L) L \cap \Sigma^+ L \Sigma^+  \ne \emptyset$,
we get that $L$ is not a solid code, which gives the claim.

$\ref{thm:solidcode:1a}\Rightarrow \ref{thm:solidcode:1b}$: Since~$L$ is a prefix code,
 its set of proper prefixes equals $\pref{L}\setminus L$, which gives the claim.

$\ref{thm:solidcode:1b}\Rightarrow \ref{thm:solidcode:2}$: Suppose~$L$ is a suffix code as stated.
Assume we have $uvw \in L$ with $u \in \Sigma^*$, $v \in L$,
and $w \in \Sigma^+$.
Let $x$ be a proper prefix of some word in~$L$.
Then $xuvw \in \Sigma^* L \Sigma^+$
and this is excluded by assumption.
Thus, $L$ is an infix code.
The assumption yields that $( \pref{L}\setminus L ) L \cap L \Sigma^+ = \emptyset$.

$\ref{thm:solidcode:2} \Rightarrow \ref{thm:solidcode:3}$: 	By \cref{lem:overlap-free-prefix}, $L$ is overlap-free.
Thus, \cref{constrsync:prop:overlap-free-inclusion-eqn}
gives $\pref{L} \cap \suff{L} \subseteq L \cup \{\varepsilon\}$.
The inclusion
$L \cup \{\varepsilon\} \subseteq \pref{L} \cap \suff{L}$
holds for any language.	

$\ref{thm:solidcode:3} \Rightarrow \ref{thm:solidcode:1}$: 
By \cref{constrsync:prop:pref-suff-overlap-free-inclusion-eqn}, $L$
is overlap-free and thus a solid code.	
\end{proof}

\begin{remark}
	As far as we are aware of,
	the condition $(\pref{L} \setminus L) L \cap \Sigma^* L \Sigma^+ = \emptyset$
	was never investigated in the literature before (apart
	from our own work~\cite{DBLP:conf/fct/Hoffmann21} 
	whose results are presented and extended
	in the present thesis).
	Note that $\pref{L}L \subseteq L\Sigma^+$ holds for every language $L \subseteq \Sigma^*$.
\end{remark}

We mention an implication for languages  $L \subseteq \Sigma^*$
such that $(\pref{L} \setminus L) L \cap \Sigma^* L \Sigma^+ = \emptyset$.

\begin{proposition}
	Let $L \subseteq \Sigma^*$
	such that $(\pref{L} \setminus L) L \cap \Sigma^* L \Sigma^+ = \emptyset$.
	Then $\pref{L}\subseteq L$
	or $uvw \in L$ with $u,v \in \Sigma^*$ and $v \in L$ implies $w = \varepsilon$.
\end{proposition}
\begin{proof}
	Suppose we have  $x \in \pref{L} \setminus L$
	for some $x \in \Sigma^*$.
	Let $v \in L$
	with $uvw \in L$ for some $u,w \in \Sigma^*$.
	Assume $|w| > 0$.
	Then we have $xuvw \in (\pref{L}\setminus L) L \cap \Sigma^* L \Sigma^+$.
	Thus, we must have $|w| = 0$.
\end{proof}

\begin{theorem} \label{thm:solidcode-char2}
	Let $L \subseteq \Sigma^+$.
	Then $L$ is a solid code
	if and only if for all $u \in \pref{L}$ and $v \in L$
	such that $uv = xwy$ for  $x \in \Sigma^*$ and $y \in \Sigma^+$ and $w \in L$
	we have $x = \varepsilon$ and $u = w$.			
\end{theorem}
\begin{proof}
	First, suppose $L$ is a solid code.
	Let $u \in \pref{L}$, $v, w \in L$,
	and $uv = xwy$ for $x \in \Sigma^*$ and $y \in \Sigma^+$.	
	By \crefitem{thm:solidcode}{thm:solidcode:1a}, we must have $u \in L$.
	By \crefitem{constr_sync:prop:selfsync_implies_avoiding}{constr_sync:prop:item:sc-in-cf}, 
	$L$ is comma-free.
	So $u, v \in L$	imply $x = \varepsilon$ and, as $L$ is a prefix code, $w = u$.	
	Conversely, suppose the stated condition holds.
	Then, clearly, $L$ is comma-free.
	Next, suppose $u \in \pref{L} \setminus L$ and $v,w \in L$
	with $uv = xwy$ for some $x \in \Sigma^*$ and $y \in \Sigma^+$.
	By assumption, $x = \varepsilon$ and $u = w$.
	However, as $u \notin L$ we have $u \ne w$.
	So we must have $(\pref{L} \setminus L)L \cap L \Sigma^+ = \emptyset$.
\end{proof}

\begin{corollary} \label{cor:solidcode-char2}
	Let $C \subseteq \Sigma^+$ be a solid code.
	Then $\pref{C}C \cap \Sigma^+ C \Sigma^+ = \emptyset$.
\end{corollary}

The condition $\pref{C}C \cap \Sigma^+ C \Sigma^+ = \emptyset$
clearly implies that~$C$ is comma-free,
but not every comma-free code fulfills it.
For example, for the comma-free code $C = \{aba,cd,bc\}$
we have $abcd \in (\pref{C}\setminus C)C \cap \Sigma bc \Sigma$.

The next examples show that $\pref{C}C \cap \Sigma^+ C \Sigma^+ = \emptyset$
does not characterize solid codes.

\begin{example}
	\begin{enumerate}
		\item For the comma-free code $C = \{abb,aba\}$ from \crefitem{constrsync:ex:binarysolidcode}{constrsync:ex:comma-free-not-solid}
		that is not a solid code we have $\pref{C}C \cap \Sigma^+ C \Sigma^+ = \emptyset$.
		
		\item Let $C = \{ab,bc\}$. Then $C$ is not a solid code,
		but we have $\pref{C}C \cap \Sigma^+ C \Sigma^+ = \emptyset$.
	\end{enumerate} 	
\end{example}

\begin{proposition} \label{constrsync:prop:prefLintersect}
	Let $L \subseteq \Sigma^+$.
	Then $\pref{L}L \cap \Sigma^+ L \Sigma^+ = \emptyset$
	if and only if for all $u \in \pref{L}$ and $v \in L$
	such that $uv = xwy$ for $x \in \Sigma^+$ and $y \in \Sigma^*$ and $w \in L$
	we have $y = \varepsilon$ and $w = v$.
\end{proposition}
\begin{proof}
	First, suppose $\pref{L}L \cap \Sigma^+ L \Sigma^+ = \emptyset$.
	In this case, we must have $y = \varepsilon$ since
	otherwise $uv \in \Sigma^+ L \Sigma^+$.
	Since~$L$ is a suffix code, this implies $w = v$.
	
	Conversely, the condition 
	implies $\pref{L}L \cap \Sigma^+ L \Sigma^+ = \emptyset$.	
\end{proof}

\begin{corollary} \label{cor:solidcode-char3}
	Let $L \subseteq \Sigma^+$.
	Then $L$ is a solid code
	if and only if for all $u \in \pref{L}$ and $v \in L$,
	if $uv = xwy$ for  $w \in L$ and $x,y \in \Sigma^*$,
	then $x = \varepsilon$ and $w = u$
	or $y = \varepsilon$ and $w = v$.
\end{corollary}
\begin{proof}
	\Cref{thm:solidcode-char2,cor:solidcode-char2,constrsync:prop:prefLintersect} show 
	that a solid code fulfills the condition.
	Conversely, it implies the condition
	from \cref{thm:solidcode-char2} if $|y| > 0$
	and thus that~$L$ is a solid code.
\end{proof}

Lastly, we give a characterization using the Levenshtein-Romanov mapping. It basically combines results
from \cref{constrsync:sec:codes:levensthein-romanov}.

\begin{proposition} \label{constrsync:prop:sigma-solid-code-char}
	Let $L \subseteq \Sigma^+$ be non-empty.
	Then $L$ is a solid code if and only if
	for all $u,v \in \Sigma^*$ the following is true:
	\begin{enumerate}
		\item \label{constrsync:prop:sigma-solid-code-char:1} 
		$\sigma_L(uv) = uv$ implies $u \notin L$ or $v = \varepsilon$,
		
		\item \label{constrsync:prop:sigma-solid-code-char:2} $\sigma_L(uv) \in L$ implies $v \notin \pref{L}\setminus L$ or $v = \varepsilon$,
		
		\item \label{constrsync:prop:sigma-solid-code-char:3} $\sigma_L(uv) \ne v$ implies $v \notin L$. 	
	\end{enumerate}
\end{proposition}
\begin{proof}
	First, suppose~$L$ is a solid code.	
	Assume $\sigma_L(uv) = uv$.
	Since~$L$ is a prefix code,
	by \cref{constrsync:prop:prefix-code:4} of \cref{constrsync:prop:prefix-code}
	we find that~$u \notin L$
	or $v = \varepsilon$.
	Next, let $\sigma_L(uv) \in L$ and $v \in \pref{L}\setminus L$.
	Then $\sigma_L(uv) = wv$ for some $w \in \Sigma^*$
	with $|w| > 0$ since $wv \in L$.
	By overlap-freeness, this yields $v = \varepsilon$.	
	\Cref{constrsync:prop:sigma-solid-code-char:3}
	follows from \cref{constrsync:prop:infix-code:6} of \cref{constrsync:prop:infix-code}
	since~$L$ is an infix code.
		
	Conversely, \cref{constrsync:prop:sigma-solid-code-char:1} 
	yields that~$L$ is a prefix code by \crefitem{constrsync:prop:prefix-code}{constrsync:prop:prefix-code:4}.
	Then, \cref{constrsync:prop:sigma-solid-code-char:3} 
	yields that~$L$ is an infix code by \cref{constrsync:prop:infix-code:6} 
	of \cref{constrsync:prop:infix-code}.
	Next, assume $uv, vw \in L$ for non-empty $u,v,w \in \Sigma^+$.
	We have $v \in \pref{L}\setminus L$ since~$L$ is a prefix code
	and $\sigma_L(uv) = uv$. Thus, \cref{constrsync:prop:sigma-solid-code-char:2} 
	does not hold. So, if all these conditions hold,
	we find that~$L$ is overlap-free.
\end{proof}

\begin{remark}
	By \cref{constrsync:prop:sigma-solid-code-char:3} 
	from \cref{constrsync:prop:sigma-solid-code-char},
	for a solid code~$C$ and $v \in C$,
	we have $\sigma_C(uv) = v$ for $u \in \Sigma^*$.
	By \cref{constrsync:prop:sigma-solid-code-char:2},
	we also have $\sigma_C(uw) \notin L$
	for every proper and non-empty prefix $w$
	of a word from~$C$.
	However, note that we do not have $\sigma_C(uw) = w$
	in this case. For example, for the solid code $C = \{abb, aac\}$
	we have $\sigma_C(uw) = uw$ for $u = a$
	and $w = a$.
\end{remark}

%\section{Characterizing Solid Codes}
%\label{constr_sync:subsec:charsolidcode}

\section{Bounded Languages}
\label{constr_sync:sec:bounded} 

%!TEX root = ./../../springer-mono-format/book.tex
%!TEX root = ./../../thesis-trier.tex

\IN{bounded language}
\IN{letter-bounded language}
\IN{strictly bounded language}
A language $L \subseteq \Sigma^*$ is \emph{bounded}~\cite{Ginsburg1964BOUNDEDAL}
if there exist words $u_1, \ldots, u_r \in \Sigma^*$ such that
$
L \subseteq u_1^* \cdots u_r^*.
$
The language $L$ is called \emph{letter-bounded} if each $u_i$ is a letter, i.e.,
$u_i \in \Sigma$ for all $i \in \{1,\ldots,r\}$.
It is called \emph{strictly bounded} if it is letter-bounded and the letters
$u_1,\ldots,u_r$ are pairwise distinct.

For example, finite languages are letter-bounded.
The bounded and letter-bounded languages
are closed under finite unions and concatenation
since for $L \subseteq u_1^* \cdots u_r^*$
and $L' \subseteq v_1^* \cdots v_l^*$
both $L\cup L'$ and $LL'$ are contained in $u_1^* \cdots u_r^*v_1^* \cdots v_l^*$.

Our first result characterizes subset of $w^*$. It
was also stated in \cite[Lem.~5.2]{Ginsburg1964BOUNDEDAL}.

\begin{proposition}\label{constrsync:prop:slender}
	Let $L \subseteq \Sigma^*$.
	Then there exists $w \in \Sigma^*$
	such that $L \subseteq w^*$
	if and only if, for all $u,v \in L$, there exists $w \in \Sigma^*$ such that $\{u,v\}\subseteq w^*$.
\end{proposition}
\begin{proof}
	If  $L \subseteq w^*$ for some $w \in \Sigma^*$,
	then in particular $\{u,v\} \subseteq w^*$ for all $u,v \in L$.	
	Conversely, suppose for all $u,v \in L$
	there exists $w \in \Sigma^*$ with $\{u,v\} \subseteq w^*$.
	Assume $L \not\subseteq w^*$ for all $w \in \Sigma^*$.
	Then for each $w \in \Sigma^*$
	with $v \in w^*$ for some $v \in L$
	we have $u \notin w^*$ for some $u \in L$.
	Now, pick a shortest word $w_1 \in \Sigma^*$  
	such that $u \in L \setminus w_1^*$
	and $v \in L \cap w_1^*$ for some $u,v \in \Sigma^+$.	
	By assumption there exists $w_2 \in \Sigma^*$
	such that $\{u, v\} \subseteq w_2^*$.
	So $w_1^i = v = w_2^j$ for some $i,j \ge 0$.
	We have	$i,j \ge 1$ and $w_1, w_2 \ne \varepsilon$ as~$|v| > 0$.
	By \cref{constr_sync:lem:two-words-code} there exists $w_3 \in \Sigma^*$
	such that $w_1, w_2 \in w_3^*$.
	Let $k \ge 1$ be such that $w_1 = w_3^k$.
	As $u \in w_2^* \subseteq w_3^*$ and $u \notin w_1^*$
	we have $k \ge 2$ so that $0 < |w_3| < |w_1|$.
	By the choice of $w_1$ and as $v \in L \cap w_3^*$ we find $L \subseteq w_3^*$, contradicting the 
	assumption.\footnote{For a similar inductive argument: Let $L = \{u_0, u_1, u_2, \ldots \}$.
		Inductively, for $i \ge 1$
		construct words $w_i \in \Sigma^*$
		with $w_i \in w_{i+1}^*$ and $\{u_{i-1},u_i\} \subseteq w_i^*$
		in the following way:
		Let $w_1 \in \Sigma^*$
		be a word such that $\{u_0, u_1\} \subseteq w_1^*$.
		For $i \ge 2$ we have $\{u_i, u_{i+1}\} \subseteq w^*$ 
		for some $w \in \Sigma^*$.
		Thus, $u_i \in w_i^*$ (by induction hypothesis) and $u_i \in w^*$.
		By \cref{constr_sync:lem:two-words-code} 
		we have $w_i, w \in w_{i+1}^*$ for some $w_{i+1} \in \Sigma^*$.
		Now, having such a sequence, $w_i \in w_{i+1}^*$
		implies $|w_{i+1}| \le |w_i|$.
		Thus there exists $N \ge 0$ such that $|w_N| = |w_{N+i}|$ for all $i \ge 0$.
		But then $w_i \in w_{i+1}^*$,
		gives $w_N = w_{N+i}$ for all $i \ge 0$.
		This yields $L \subseteq w_N^*$.}
\end{proof}

For regular subsets of $w^*$ we have the following representation.

\begin{proposition}\label{constr_sync:prop:bounded_normal_form_single_word}
	Let $L \subseteq w^*$ for some $w \in \Sigma^*$.
	Then $L$ is regular if and only 
	if $L = I \cdot (w^p)^* \cup F$
	for finite sets $I, F \subseteq w^*$ and some $p > 0$.
\end{proposition}
\begin{proof} 
	First, suppose $L \subseteq w^*$ is regular. By \cref{prelim:thm:regular_iff_quotients_finite}
	there exist $0 \le i < j$ such that
	$(w^{i})^{-1} L = (w^j)^{-1} L$. 
	Set $p = j - i$. Then, for all $k \ge i$ this gives
	\[
	w^k \in L \Leftrightarrow w^{k + p} \in L.
	\]
	This yields that for $F = \{ w^k \mid 0 \le k < i \land w^k \in L \}$ and
	$I = \{ w^k \mid i \le k < j \land w^k \in L \}$
	we have $L = I\cdot (w^p)^* \cup F$.		
	Conversely, a language of the form $L = I \cdot (w^p)^* \cup F$
	for finite sets $I, F \subseteq w^*$ and some $p \ge 0$
	is clearly regular and contained in $w^*$.
\end{proof}

The next representation theorem for bounded regular sets is from~\cite[Thm.~1.2]{GinsburgSpanier66}.

\begin{theorem}[{Ginsburg and Spanier~\cite[Thm.~1.2]{GinsburgSpanier66}}]
	\label{constr_sync:thm:bounded_normal_form}
	Let $w_1, \ldots, w_r \in \Sigma^*$.
	A~language $L \subseteq w_1^* \cdots w_r^*$ is regular if and only if
	it is a finite union of sets of the form $L_1 \cdots L_r$,
	where each $L_i$ is a regular language in $w_i^*$ for $i \in \{1,\ldots,r\}$.
\end{theorem} 

\begin{corollary}\label{constr_sync:cor:bounded_normal_form_single_state_aut}
	$L \subseteq w_1^* \cdots w_r^*$ 
	for $w_1, \ldots, w_r \in \Sigma^*$
	is regular if and only if
	it is a finite union of sets of the form $w_1^{i_1} (w_1^{p_1})^* \cdots w_r^{i_r} (w_r^{p_r})^*$
	for numbers $i_1, \ldots, i_r, p_1, \ldots, p_r \ge 0$.
\end{corollary}
\begin{proof}
	By \crefitem{prelim:langeqn}{prelim:langeqn:concatcup} 
	we can write  $I \cdot (w^p)^* \cup F$ 
	for finite $I$ and~$F$ as a finite union
	of sets of the form $w^i (w^q)^*$
	with $i, q \ge 0$ (note that $q = 0$ is allowed).
	Thus, the claim follows
	from \cref{constr_sync:thm:bounded_normal_form,constr_sync:prop:bounded_normal_form_single_word}.
\end{proof}

\IN{equal-letter run}
\label{def:rho}
The function $\rho \colon \Sigma^* \to \mathbb N_0$ given, for $w \in \Sigma^*$ and $a,b \in \Sigma$, by
\[ 
\rho(wab) = \begin{cases} 
	\rho(wa) + 1 & \text{if } a \ne b, \\
	\rho(wa)     & \text{if } a = b,
\end{cases}
\]
and $\rho(\varepsilon) = \rho(a) = 0$ counts the number of \emph{symbol switches}, 
or equivalently the number of \emph{equal-letter runs} minus one.

\begin{example}
\begin{enumerate}
\item $\rho(abab) = \rho(aaabbab) = 3$.

\item For $L = (ab)^+$
we have $\{ \rho(u) \mid u \in L \} = \{ n \mid n \text{ is odd} \}$.
\end{enumerate}	
\end{example}

We list a few properties of $\rho$.

\begin{proposition}
\label{constrsync:prop:rho}
\begin{enumerate}
\item $\rho(a^i) = 0$ for $a \in \Sigma$ and $i \ge 0$.

\item \label{constrsync:prop:rho:2} For $w \in \Sigma^+$ we have $\rho(w) = |\{ i \in \{1,\ldots,n-1\} : w_{i} \ne w_{i+1} \}|$
 where $w = w_1 \cdots w_n$ with $w_i \in \Sigma$ for $i \in \{1,\ldots,n\}$ and $n \ge 1$.

\item \label{constrsync:prop:rho:3} For $w \in \Sigma^+$
 we have $\rho(w) = k - 1$
 if and only if  there exist letters $a_1, \ldots, a_k \in \Sigma$, $k \ge 1$, such that $w \subseteq a_1^+ \cdots a_k^+$ 
 and $a_i \ne a_{i+1}$ for $i \in \{1,\ldots,k-1\}$.

\item \label{constrsync:prop:rho:4} $\rho(w) = \min\{ k \ge 1 \mid \exists a_1, \ldots, a_k \in \Sigma \colon w \subseteq a_1^* \cdots a_k^* \} - 1$ for $w \in \Sigma^*$.

\item \label{constrsync:prop:rho:item:uv} For $u,v \in \Sigma^*$ we have
$
\rho(u) + \rho(v) \le \rho(uv) \le \rho(u) + \rho(v) + 1.
$
In particular, $\rho(u) \le |u|$.
\end{enumerate}
\end{proposition}
\begin{proof}
	The first four equations follow readily. 
	For~\cref{constrsync:prop:rho:item:uv}, let $u,v \in \Sigma^*$ be two words.  We can assume $u,v \ne \varepsilon$.
	By~\cref{constrsync:prop:rho:3}
	there exist $a_1, \ldots, a_n, b_1, \ldots, b_m \in \Sigma$ for $n,m \ge 1$
	such that $\rho(u) = n - 1$ and $\rho(v) = m - 1$ where $u \subseteq a_1^+ \cdots a_n^+$
	with $a_i \ne a_{i+1}$ for $i \in \{1,\ldots,n-1\}$
	and $v \subseteq b_1^+ \cdots b_{m}^+$
	with $b_j \ne b_{j+1}$ for $j \in \{1,\ldots,m-1\}$.
	Then $uv \subseteq a_1^+ \cdots a_n^+ b_1^+ \cdots b_m^+$
	and we get
	\[
	\rho(uv) =
	\begin{cases}
		n + m - 1 & \text{if } a_n \ne b_1, \\
		n + m - 2 & \text{if } a_n = b_1.
	\end{cases}
	\]
	For an alternative argument,
	write $u = u_1 \cdots u_n$ and $v = v_1 \cdots v_m$
	for $n,m \ge 0$ and $u_i \in \Sigma$, $v_j \in \Sigma$ for $i \in \{1,\ldots,n\}$ and $j \in \{1,\ldots,m\}$.
	Then, by \cref{constrsync:prop:rho:2}, we find
	$\rho(uv) = |\{ i \in \{1,\ldots,n-1\} \mid u_i \ne u_{i+1} \}| + |\{ n \mid u_n \ne v_1 \}|
	+ |\{ i \in \{1,\ldots,m\} \mid v_{i} \ne v_{i+1}\}|
	= \rho(u) + \rho(v) +  |\{ n \mid u_n \ne v_1 \}|$.
\end{proof}

Next, we characterize letter-bounded languages with the mapping $\rho$.

\begin{proposition}
	\label{constr_sync:prop:letter_bounded_iff_equal_letter_runs_bounded}
	A language $L \subseteq \Sigma^*$ is letter-bounded if and only if there
	exists $n \ge 0$ such that $\applyfunc{u}{\rho} \le n$
	for all $u \in L$.
\end{proposition}
\begin{proof}  
	First, suppose $\rho(u) \le n$ for all $u \in L$.
	Let $f_1, \ldots, f_m$ be some enumeration
	of all functions  $\{1,\ldots,n\} \to \Sigma$
	where $m = |\Sigma|^n$.
	Let $u \in L$.
	We have $u \in a_1^* \cdots a_n^*$
	for some $a_1, \ldots, a_n \in \Sigma$.
	Thus $u \in \prod_{i=1}^{m} f_i(1)^* \cdots f_i(n)^*$.
	Hence $L$ is letter-bounded.	
	Conversely, suppose $L \subseteq \Sigma^*$ is letter-bounded. Then there exist $a_1, \ldots, a_n \in \Sigma$
	such that $L \subseteq a_1^* \cdots a_n^*$. So, we have $\rho(u) \le n - 1$ for all $u \in L$
	by \crefitem{constrsync:prop:rho}{constrsync:prop:rho:4}.
\end{proof}

\begin{corollary} \label{syncmax:lem:necessary-letter-bounded}
	Let $L \subseteq \Sigma^*$.
	Suppose there exists $u \in  \Sigma^* a \Sigma^* b \Sigma^*$
	for distinct $a,b \in \Sigma$
	such that $\Sigma^* u^n \Sigma^* \cap L \ne \emptyset$
	for infinitely many $n$.
	Then $L$ is not letter-bounded.
\end{corollary}
\begin{proof}
	We have $\rho(u) \ge 1$.
	Thus for $vu^nw \in L$ with $n \ge 1$
	we have $\rho(v u^n w) \ge \rho(u^n) \ge n$
	by \crefitem{constrsync:prop:rho}{constrsync:prop:rho:item:uv}.
	Whence, by~\cref{constr_sync:prop:letter_bounded_iff_equal_letter_runs_bounded},
	the language~$L$ is not letter-bounded.
\end{proof}

For bounded languages, the condition from \cref{syncmax:lem:necessary-letter-bounded} also characterizes letter-boundedness, as we show next.

\begin{proposition} \label{syncmax:lem:non-letter-bounded}
	Let $L \subseteq \Sigma^*$ be bounded.
	Then~$L$ is letter-bounded 
	if and only if for all $u \in \Sigma^*$
	such that $\Sigma^* u^n \Sigma^* \cap L \ne \emptyset$ for infinitely many $n$
	we have $u \subseteq a^*$ for some $a \in \Sigma$.
\end{proposition}
\begin{proof}
	First, suppose~$L$ is letter-bounded. Then, by \cref{syncmax:lem:necessary-letter-bounded},
	for all $u \in \Sigma^*$ such that $\Sigma^* u^n \Sigma^* \cap L \ne \emptyset$ for infinitely many $n$
	we must have $u \subseteq a^*$ for some $a \in \Sigma$.
	
	Conversely, suppose the stated condition holds.
	Let $L \subseteq u_1^* \cdots u_k^*$ for some $u_1, \ldots, u_k \in \Sigma^*$. 
	Suppose we have $u_j \in \Sigma^* a \Sigma^* b \Sigma^*$ for some $j \in \{1,\ldots,k\}$
	with distinct $a,b \in \Sigma$
	and such that for each $N \ge 0$
	there exist $i_1, \ldots, i_k \ge 0$ with $i_j \ge N$
	such that $u_1^{i_1} \cdots u_k^{i_k} \in L$. Then the the assumption is violated
	with $u = u_j$.	
	Thus, there exists $N \ge 0$
	such that $u_1^{i_1} \cdots u_k^{i_k} \in L$
	with $i_j > N$ for some $j \in \{1,\ldots,k\}$
	implies $u_{j} \subseteq a^*$ for some $a \in \Sigma^*$.
	Then pick all those indices $1 \le j_1 < \cdots < j_l \le k$ for some $l \in \{1,\ldots,k\}$
	where for $i \in \{1,\ldots,l\}$ we have $u_{j_i} \in \Sigma^* a \Sigma^* b \Sigma^*$ for distinct $a,b \in \Sigma$.
	Then for $u_1^{i_1} \cdots u_k^{i_k} \in L$,
	by assumption, we have $i_{j_r} \le N$ for $r \in \{1,\ldots,l\}$
	and $\rho(u_i) = 0$ for $i \notin \{j_1, \ldots, j_l\}$.
	With \cref{constrsync:prop:rho}, for $M =i_1 + \cdots + i_k + k $ this gives
	\[
	\rho( u_1^{i_1} \cdots u_k^{i_k} ) - M
	\le i_1 \rho(u_1) + \cdots + i_k \rho(u_k)
	\le \sum_{r = 1}^l i_{j_r} \rho(u_{j_r}) 
	\le N \cdot \sum_{r = 1}^l \rho(u_{j_r}).
	\]
	Thus, by \cref{constr_sync:prop:letter_bounded_iff_equal_letter_runs_bounded},
	the language~$L$ is letter-bounded.
\end{proof}

Next, we characterize regular letter-bounded languages.

\begin{proposition}\label{constr_sync:prop:letterboundedreg}
	Let $L \subseteq \Sigma^*$ be regular. Then the following are equivalent.
	\begin{enumerate}
		\item \label{constr_sync:prop:letterboundedreg:item:letterbounded} $L$ is letter-bounded.
		
		\item \label{constr_sync:prop:letterboundedreg:item:uvw} If $uv^*w \subseteq L$ with $u,v,w \in \Sigma^*$, then $v^* \subseteq a^*$ for some $a \in \Sigma$.
		
		\item \label{constr_sync:prop:letterboundedreg:item:PDFA} For every PDFA $\automat{B}{P}{\Sigma}{\mu}{p_0}{F}$
		with $L = L(\mathcal B)$ and state $p \in P$
		that is accessible and coaccessible, we have $L(\mathcal B_{p,\{p\}}) \subseteq a^*$ for some $a \in \Sigma$.
	\end{enumerate} 	
\end{proposition}
\begin{proof}
	$\ref{constr_sync:prop:letterboundedreg:item:letterbounded} \Rightarrow~\ref{constr_sync:prop:letterboundedreg:item:uvw}$: Suppose there exist $u,v,w \in \Sigma^*$
	with $uv^*w \subseteq L$ and $|v|_a, |v|_b > 0$ for two distinct $a,b \in \Sigma$.
	Then, by \cref{syncmax:lem:necessary-letter-bounded}, $L$ is not letter-bounded.
	
	$\ref{constr_sync:prop:letterboundedreg:item:uvw}\Rightarrow~\ref{constr_sync:prop:letterboundedreg:item:PDFA}$: 
	Let  $\automat{B}{P}{\Sigma}{\mu}{p_0}{F}$
	be a PDFA.
	Assume  $v \in L(\mathcal B_{p, \{p\}})$
	for some  accessible and coaccessible $p \in P$ and 
	$v \in \Sigma^*$ with $|v|_a > 0$ and $|v|_b > 0$
	for distinct $a,b \in \Sigma$.
	Then
	$\mu(p_0, u) = p$
	and $\mu(p, w) \in F$ for some $u,w \in \Sigma^*$.
	Hence $uv^*w \subseteq L(\mathcal B)$
	and~\cref{constr_sync:prop:letterboundedreg:item:uvw} 
	is not true, giving the claim.	
%	
%	
%	By assumption there exist $u,v \in \Sigma^*$ such that $\mu(p_0, u) = p$
%	and $\mu(p, v) \in F$.
%	Let $w \in L(\mathcal B_{p,\{p\}})$.
%	Then $uw^*v \subseteq L$ and so $w \in a^*$ for some $a \in \Sigma$.
%	Now, let $w' \in L(\mathcal B_{p,\{p\}})$.
%	Similarly, we obtain $w' \in b^*$ for some $b \in \Sigma$.
%	As $u(ww')^*v \subseteq L$ we conclude $a = b$
%	and, as $w'$ was arbitrarily chosen, we get $L(\mathcal B_{p,\{p\}}) \subseteq a^*$.
	
	$\ref{constr_sync:prop:letterboundedreg:item:PDFA} \Rightarrow~\ref{constr_sync:prop:letterboundedreg:item:letterbounded}$:	
	Suppose $L$ is not letter-bounded.
	Let $\automat{B}{P}{\Sigma}{\mu}{p_0}{F}$
	be a PDFA with $L = L(\mathcal B)$ and $n = |P|$.
	By \cref{constr_sync:prop:letter_bounded_iff_equal_letter_runs_bounded}
	there exists $u \in L(\mathcal B)$
	with $\rho(u) \ge 2n - 1$.
	By \crefitem{constrsync:prop:rho}{constrsync:prop:rho:3},
	we have $u = a_1^{i_1} \cdots a_{2n}^{i_{2n}}$
	for some $a_1, \ldots, a_{2n} \in \Sigma$ with $a_j \ne a_{j+1}$ for $j \in \{1,\ldots,2n-1\}$
	and $i_j \ge 1$ for $j \in \{1,\ldots,2n\}$.
	Set $p_j = \mu(p_0, a_1^{i_1} \cdots a_{i_{2j}}^{i_{2j}})$
	for $j \in \{0,1,\ldots,n\}$.\footnote{We set $a_1^{i_1} \cdots a_{0}^{i_{0}} = \varepsilon$.}
	By finiteness there exists $0 \le j < j' \le n$ such that $p_j = p_{j'}$.
	Set $w = a_{2j}^{i_{2j}} \cdots a_{2j'}^{i_{2j'}}$.
	Then we have $w^* \subseteq L(\mathcal B_{p_j, \{p_j\}})$.
	Since $a_{2j} \ne a_{2j+1}$ the language $w^*$ is not contained in $a^*$
	for each $a \in \Sigma$.
	The PDFA $\mathcal B$ accepts~$u$,
	and so the state $p_j$ is accessible and coaccessible.
	Thus \cref{constr_sync:prop:letterboundedreg:item:PDFA} does not hold.
\end{proof}

\section{Monoid Languages}
\label{constr_sync:subsec:monoid}

%!TEX root = ./../../thesis-trier.tex

\IN{monoid language}
A submonoid $L \subseteq \Sigma^*$
is also called a \emph{monoid (language)} for short.
The language $L \subseteq \Sigma^*$
is a monoid language if and only if % $LL \subseteq L$, % ne, neutrales element noch drin
$L^* \subseteq L$.
For any $L \subseteq \Sigma^*$, we have $\langle L \rangle_M = L^*$.

For a monoid language $L \subseteq \Sigma^*$
there exists a unique minimal generating set $X$, i.e., a set $X$ minimal for set-inclusion
with $L = X^*$. It is given as 
the set of non-empty words in $L$ that cannot be written as a product of two non-empty words from $L$, namely (see~\cite[p.~4]{lothaire_1997}):
\begin{equation}
	\label{constr_sync:eqn:generating_set_submonoid}
	X = ( L \setminus \{\varepsilon\} ) \setminus ( L \setminus \{\varepsilon\})^2.
\end{equation}
For every set $Y \subseteq \Sigma^*$
such that $Y^* = L$ we have $X \subseteq Y$ (more precisely,
we have $X = ( Y \setminus \{\varepsilon\} ) \setminus ( Y \setminus \{\varepsilon\})^2$).
This implies that~$L$ is generated by a prefix code 
if and only if~$X$ is a prefix code. The next examples
give languages where $X$ is not a code, and where is it a code, but not a prefix code.

\begin{example}
	For $L = \{a,ab,ba\}^*$ we have $X = \{a,ab,ba\}$, which is not a code since $a(ba) = (ab)a$.
\end{example}

\begin{example}
	For $L = \{a,ab\}^*$ we have $X = \{a,ab\}$. 	
\end{example}

\Cref{constr_sync:eqn:generating_set_submonoid} implies that $L$ is regular if and only if $X$ is regular. 

We always have $L \setminus ( L \Sigma^+ \cup \{\varepsilon \} ) \subseteq X$
since $(L \setminus \{\varepsilon\})^2 \subseteq L\Sigma^+$.
In case $X$ is a prefix code, 
the following equalities hold:\footnote{Since codes are defined here as non-empty subsets of $\Sigma^+$,
	for a prefix code~$X$ we have $X^* \ne \{\varepsilon\}$.}
\begin{equation}\label{constr_sync:eqn:prefix_code_gen}
	X = L \setminus ( L \Sigma^+ \cup \{\varepsilon \} )
	= \{ w \in L \setminus \{\varepsilon\} \mid \pref{w} \cap L = \{ w \} \}.
\end{equation} 

Regular languages that are monoids generated by prefix codes
admit an automata-theoretical characterization 
%(that we need in Section~\ref{constr_sync:sec:PSPACE-hardness})
as languages
accepted by PDFAs
whose start state is the only final state (a related statement
is given in \cite[Prop.~3.2]{Eilenberg1974}).

\begin{theorem}[{\cite[Prop.~3.2.5]{DBLP:books/daglib/0025093},\cite[Thm.~3.1]{DBLP:conf/litp/Restivo88}}]
	\label{constr_sync:thm:regular_prefix_codes_and_automata}
	Let $L \subseteq \Sigma^*$.
	Then we have $L = L(\mathcal B)$
	for some PDFA $\automat{B}{P}{\Sigma}{\mu}{p_0}{\{p_0\}}$
	if and only if  $L = \{\varepsilon\}$ or $L$ is a regular language that is a monoid generated by a prefix code that is regular.
\end{theorem}
\begin{proof}	
	We can assume $L \ne \{\varepsilon\}$. Let $L = L(\mathcal B)$
	for a PDFA $\automat{B}{P}{\Sigma}{\mu}{p_0}{\{p_0\}}$.
	Set $X = \{ u \in  L \mid \forall v \in \pref{u} \setminus \{u\} \colon \mu(p_0, v) \ne p_0 \}$.
	Then~$X$ is a regular prefix code (we obtain a PDFA accepting it by removing all outgoing transitions
	from final states in~$\mathcal B$) such that $L = X^*$.
	
	Conversely, let $L = X^*$ where $X$ is a regular prefix code.
	By \cref{constr_sync:prop:aut_prefix_free} 
	we obtain a trim PDFA $\automat{B}{P}{\Sigma}{\mu}{p_0}{\{q\}}$
	such that $q$ has no outgoing transitions.
	Thus we can merge the states $q_0$ and~$q$ to a single
	state $p$ to obtain a new PDFA where all transitions going out of~$p_0$ go out of this new state
	and all transitions ending in some state from $\{q, p_0\}$
	now end at~$p$.
	This new PDFA with start state~$p$ and final state set~$\{p\}$
	accepts~$L$.
\end{proof}

\Cref{constr_sync:thm:regular_prefix_codes_and_automata} has been extended in~\cite[Thm.~3.2]{DBLP:conf/litp/Restivo88}
to characterize submonoids that are regular languages as precisely those languages $L \subseteq \Sigma^*$
such that $L = \{ u \in \Sigma^* \mid \applytrans{\mu}{S}{u} = S \}$ for a PDFA $\automat{B}{P}{\Sigma}{\mu}{p_0}{F}$
and $S \subseteq P$.

\IN{literal automaton}
\begin{remark} \label{syncmax:rem:general-aut-prefix-code}
	The construction from the proof of
	\crefitem{constr_sync:prop:aut_prefix_free}{constr_sync:prop:aut_prefix_free:item:min}
	(used in the proof of \cref{constr_sync:thm:regular_prefix_codes_and_automata})
	is closely related to the rooted tree from
	\cref{syncmax:rem:root-tree-prefix-code}.
	In fact, this rooted tree can be considered as a (possibly infinite)
	PDFA accepting the prefix code.
	This PDFA is called the \emph{literal automaton}
	of the prefix code in~\cite[Sect.~3.2]{DBLP:books/daglib/0025093}.
	See also~\cite[Sect.~1.3 \& 4.3.8]{DBLP:phd/hal/Ryzhikov20},
	where this construction is used for finite maximal prefix codes.
	In the construction from
	\crefitem{constr_sync:prop:aut_prefix_free}{constr_sync:prop:aut_prefix_free:item:min},
	we additionally incorporate the Myhill--Nerode relation to obtain a finite automaton whenever the prefix code is regular.
\end{remark}

Next, we show that monoid languages in $u^*$ for some $u \in \Sigma^*$ are always regular.

\IN{Frobenius coin problem}
\begin{proposition}\label{constrsync:thm:sparsemonoid-regular}
	Le $L \subseteq u^*$ be a monoid language for some $u \in \Sigma^*$.
	Then there exist $d \ge 1$ and a multiple of $n$ of~$d$
	such that $L = u^n(u^d)^* \cup F$
	for a finite set $F \subseteq \{ u \}^{< n} \cap (u^d)^*$. 
	In particular, $L$ is regular.
\end{proposition}
\begin{proof}
	We can assume $L$ is infinite.
	Let $d = \min\{ \gcd(i_1, \ldots, i_r) \mid \{ u^{i_1}, \ldots, u^{i_r} \} \subseteq L\}$.		
	Then there exist numbers $i_1, \ldots, i_r \ge 0$ with $\{ u^{i_1}, \ldots, u^{i_r} \} \subseteq L$
	such that $d = \gcd( i_1, \ldots, i_r )$.
	Let $i \ge 0$ with $u^i \in L$.
	As $d = \gcd( a_1, \ldots, a_r, i )$ 
	the number~$d$ divides~$i$.
	This yields $L \subseteq (u^d)^*$.	
	By B\'{e}zout's identity 
	there exists a number $n \ge 0$ divisible by $d$ such that every
	number $m$ divisible by $d$
	with $m \ge n$ can be written as $m = a_1 i_1 + \cdots + a_r i_r$
	for numbers $a_1, \ldots, a_r \ge 0$.\footnote{We note that the \emph{Frobenius coin problem}
		asks for the smallest $n$ for given $i_1, \ldots, i_r$ when~$d = 1$, see~\cite{Alfonsis05}.}
	So for $u^m \in u^n(u^d)^*$ we have
	$u^m \in  (u^{i_1})^* \cdots (u^{i_r})^* \subseteq L$
	and we get $u^n (u^d)^* \subseteq L$.
	As $L \subseteq (u^d)^*$ and $n$ is a multiple of $d$, the language
	$F = L \setminus u^n (u^d)^*$ is finite and $F \subseteq \{u\}^{< n}$.
	Hence, $L = a^n (u^d)^* \cup F$.
\end{proof}

The next corollary generalizes the well-known folklore fact
that the Kleene star of any unary language is always regular.
Via the mapping $n \mapsto u^n$ for a fixed $u \in \Sigma^*$ 
it is in fact a reformulation of the statement
that a submonoid of~$\mathbb N_0$
is ultimately periodic~\cite[Prop.~2.2]{RosalesGarciaSanchez2009}.

\begin{corollary}
	Let $L \subseteq u^*$. Then $L^*$ is a regular language.
\end{corollary}
\begin{proof}
	As $L^*$ is a monoid and since $L^* \subseteq (u^*)^* = u^*$
	by \cref{prelim:langeqn:Kleenestar-idemp},
	the claim follows from 
	\cref{constrsync:thm:monoid-sparseregular-or-exp}.  
\end{proof} 

%https://www.google.com/search?q=kleene+star+of+unary+language+always+regular&rlz=1C1GCEU_deDE823DE823&oq=kleene+star+of+unary+language+always+regular&aqs=chrome..69i57j33i160.12614j0j7&sourceid=chrome&ie=UTF-8
%https://cs.stackexchange.com/questions/21765/kleene-star-of-an-infinite-unary-language-always-yields-a-regular-language
%https://cs.stackexchange.com/questions/10013/if-l-is-a-subset-of-0-then-how-can-we-show-that-l-is-regular
%https://cs.stackexchange.com/questions/164/what-are-the-possible-sets-of-word-lengths-in-a-regular-language/175#175
%https://mathworld.wolfram.com/FrobeniusNumber.html
%https://www.youtube.com/watch?v=H-GYBjDpT6U

\IN{semigroup language}
\begin{remark}
	A \emph{semigroup language} is a language $L \subseteq \Sigma^*$
	such that $L^+ \subseteq L$. Equivalently, $L$ forms a subsemigroup of~$\Sigma^*$.
	Every monoid language
	is also a semigroup language, but not vice versa. For example, $aaaa^*$
	is a semigroup language that is not a monoid language.
	However, a language $L \subseteq \Sigma^*$ is a semigroup language
	if and only if $L \cup \{\varepsilon\}$ is a monoid language.
	Similar statements can be given for semigroup languages and 
	our later results on the constrained synchronization problem w.r.t.~monoid
	languages are easily adaptable and also hold for semigroup languages.
\end{remark}

\section{The Growth of Languages}
\label{constr_sync:sec:growth} 

%!TEX root = ./../../thesis-trier.tex

\IN{sparse language}
\IN{exponential growth language}

A language $L \subseteq \Sigma^*$ is \emph{sparse} (resp.~has \emph{exponential growth})
if $|L \cap \Sigma^{\le n}|$ is bounded by a polynomial in~$n$ (resp.~there exists a real number $c > 1$ such that $|L \cap \Sigma^{\le n}| \ge c^n$ for infinitely many~$n$).\footnote{$|L \cap \Sigma^{\le n}|$ is bounded by a polynomial in~$n$ (resp.~has exponential growth)	if and only if $|L \cap \Sigma^n|$ is bounded by a polynomial in~$n$ (resp.~has exponential growth). So, we can equivalently use $|L \cap \Sigma^n|$ in both definitions.}

\begin{example}
	\label{ex:only_one_letter_infinitely}
	\begin{enumerate}
	\item The code $ab^*$ is sparse whereas the code $c\{a,b\}^*c$
	has exponential growth.
	
	\item \label{ex:only_one_letter_infinitely:2} Let $L \subseteq u^*$ with $u \in \Sigma^*$.
	Then we have $L \cap \Sigma^{\le n} \subseteq \{\varepsilon, u, \ldots, u^n\}$.
	Thus $|L \cap \Sigma^{\le n}| \le n + 1$,
	and so $L$ is sparse.
	
	\item  \label{ex:only_one_letter_infinitely:3}	Let $L \subseteq \Sigma^*$ and $a \in \Sigma$. If there
	exists $N \ge 0$ such that
	$w \in L$
	implies $|w|_b \le N$
	for all $b \in \Sigma \setminus\{a\}$,
	then the language $L$ is a sparse language.
	
	\item \label{ex:only_one_letter_infinitely:4} Let $u,v \in \Sigma^*$ and $L \subseteq \Sigma^*$.
	Then $L$ is sparse (resp.~has exponential growth) if and only if $uLv$
	is sparse (resp.~has exponential growth) since 
	$|uLv \cap \Sigma^n| = |L \cap \Sigma^{n - |u| - |v|}|$ for $n \ge |u| + |v|$. (in fact,
	the map $w \mapsto uwv$ is a bijection between $\Sigma^*$ and $u\Sigma^*v$.)
	\end{enumerate}
\end{example}

Next, generalizing \crefitem{ex:only_one_letter_infinitely}{ex:only_one_letter_infinitely:2}, we show that every bounded language is sparse, a proof is also given in \cite[Prop.~1]{DBLP:journals/tcs/Raz97}.

\begin{proposition}[{\cite[Prop.~1]{DBLP:journals/tcs/Raz97}}] \label{constr_sync:prop:sparse}
	A bounded language is sparse.
\end{proposition}
\begin{proof}
	Let $L \subseteq u_1^* \cdots u_k^*$ for $u_1, \ldots, u_k\in \Sigma^*$.
	Then $L \cap \Sigma^{\le n}$
	is contained in $\{ u_1^{i_1} \cdots u_k^{i_k} \mid i_1 |u_1| + \cdots + i_k |u_k| \le n \}$,
	which is contained
	in the set $\{ u_1^{i_1} \cdots u_k^{i_k} \mid i_1 |u_1| + \cdots + i_k |u_k| \land \forall j \in \{1,\ldots,k\} \colon 0 \le i_j \le n  \}$ of size at most $(n+1)^k$.
\end{proof}

Next, we give some examples of sparse languages that are not bounded.

\begin{example} \label{syncmax:growth:ex:sparse-not-bounded}
	\begin{enumerate}
		\item Let $L = \{ (a^n b^n)^n \mid n \ge 1 \}$.
		This language is sparse since
		\[
		L \cap \Sigma^i = 
		\begin{cases}
			\{ (a^n b^n)^n \} & \text{if $i = 2n^2$ for some $n \ge 1$,} \\
			\emptyset         & \text{otherwise.}  	
		\end{cases}
		\]
		Assume~$L$ is bounded.
		We have $\rho( (a^n b^n)^n ) = 2n - 1$
		and thus~$L$ is not letter-bounded by \cref{constr_sync:prop:letter_bounded_iff_equal_letter_runs_bounded}.
		By \cref{syncmax:lem:non-letter-bounded}
		there exists $u$ with $|u|_a > 0$ and $|u|_b > 0$
		such that $\Sigma^* u^i \Sigma^* \cap L \ne \emptyset$ 
		for infinitely many~$i$.
		Without loss of generality, we have $ab^ma \in \factor{uu}$
		for some $m \ge 0$.
	    Let $vu^iw \in L$ for $i \ge 2$ and $v,w \in \Sigma^*$;
	    and then $\Sigma^* u^j \Sigma^* \cap L \ne \emptyset$
	    for $j \ge |vu^iw| + 1$.
	    Thus~$L$ contains two distinct words with the factor $ab^ma$,
	    but this contradicts the definition of~$L$
	    since $(a^m b^m)^m$ is the only word with this factor.
		
		\item Let $u_n = b a^{i_1} b a^{i_2} \cdots b a^{i_n}b$
		for some strictly increasing sequence $1 \le i_1 < i_2 < \cdots$.
		Set $L = \{ u_1, u_2, \ldots\}$.
		Then $L$ contains at most one word of each length
		since the lengths of the $u_n$ strictly increase.
		Thus $L$ is sparse.
		Assume $L$ is bounded. Since $\rho(u_n) = 2n$, $L$ is not letter-bounded.
		By \cref{syncmax:lem:non-letter-bounded}
		there exists $u$ with $|u|_a > 0$ and $|u|_b > 0$
		such that $\Sigma^* u^i \Sigma^* \cap L \ne \emptyset$ 
		for infinitely many~$i$.
		No word in~$L$ contains the factor $bb$.
		Thus a factor of the form $ba^jb$ for some $j \ge 0$ appears
		at least twice in $uuu$.
		We have $\Sigma^* uuu \Sigma^* \cap L \ne \emptyset$.
		However each factor of the form $ba^jb$ for some $j \ge 1$
		appears at most once in each word from~$L$,
		so this is not possible.
		
		\item \label{syncmax:growth:ex:sparse-not-bounded:3} (Pin~\cite[Sect.~XII.4.2]{Pin2022}) Consider
		a bijection $f \colon \{a,b\}^* \to \mathbb N_0$
		such that $f(u) \ge |u|$ for all $u \in \{a,b\}^*$.
		One can construct such a bijection by letting $f(u)$
		be the position of~$u$ in the shortlex order: $\varepsilon$, $a$, $b$, $aa$, $ab$, $ba$, $bb$, $aaa$, $aab$, $aba$, $\ldots$. Now, let
		\[
		 L = \{ u c^{f(u) - |u|} \mid u \in \{a,b\}^*\}.
		\]
		Then for $n \ge 0$ 
		we have $L \cap \Sigma^n = \{ uc^{f(u) - |u|} \}$
		for the unique $u \in \{a,b\}^*$
		with $f(u) = n$ and thus~$L$ is sparse.		
		If $L \subseteq u_1^* \cdots u_k^*$,
		then $\pref{L} \subseteq \bigcup_{i=1}^{k} \bigcup_{v \in \pref{u_i}} u_1^* \cdots u_{i-1}^* u_i^* v$;
		and thus the set of prefixes of~$L$ is also bounded if~$L$ is bounded.
		However, $\{a,b\}^* \subseteq \pref{L}$
		and this language has exponential growth.
		So~$\pref{L}$ has exponential growth
		and thus cannot be bounded by \cref{constr_sync:prop:sparse},
		which implies that~$L$ is not bounded.
	\end{enumerate}	
\end{example}

Next, in \cref{constrsync:thm:reg-sparse-or-exp},
we show that a regular language is either bounded, and thus sparse,
or has exponential growth. But first, in \cref{constr_sync:cor:code_growth_of_star},
we give a sufficient condition for exponential growth. 

\begin{proposition} \label{constr_sync:cor:code_growth_of_star}
	Let $L \subseteq \Sigma^*$ be a language
	containing two words $u,v \in L$
	that are not powers of a common word
	and such that $\{u,v\}^* \subseteq L$.
	Then~$L$ has exponential growth.	
\end{proposition}
\begin{proof}
	By \cref{constr_sync:lem:two-words-code}, $C = \{u,v\}$
	is a code.
	By Proposition~\ref{constr_sync:prop:coding_morphisms} 
	there exists a coding morphism $\varphi \colon \Gamma^* \to \Sigma^*$
	such that $\varphi(\Gamma^*) = C^*$.
	We can suppose $\Gamma = \{a,b\}$ where $a \ne b$.	
	Since~$\varphi$ is injective, we have 
	$|\applyfunc{\{a,b\}^i}{\varphi}| = 2^i$
	for all $i \ge 0$.
	As $\varphi(\{a,b\}^{\lfloor n / m \rfloor}) \subseteq C^* \cap \Sigma^{\le n}$ 
	for
	$m = \max\{ |\applyfunc{a}{\varphi}|, |\applyfunc{b}{\varphi}| \}$
	and $n \ge 0$, $C^*$ has exponential growth.
	Since~$C^* \subseteq L$, also $L$ has exponential growth.
\end{proof}

As \crefitem{syncmax:growth:ex:sparse-not-bounded}{syncmax:growth:ex:sparse-not-bounded:3}
shows, $\{u,v\}^* \subseteq \factor{L}$ for a code $\{u,v\}$ does not imply exponential growth.
However, it implies that~$L$ cannot be bounded, as stated in~\cite[Lem.~5.3]{Ginsburg1964BOUNDEDAL}. 

\begin{corollary} \label{constrsync:thm:monoid-sparseregular-or-exp}
	A monoid language $L \subseteq \Sigma^*$ is either
	a regular subset of $u^*$ for some $u \in \Sigma^*$ or has exponential growth.
\end{corollary}
\begin{proof}
	First, suppose $L \subseteq u^*$ for some $u \in \Sigma^*$.
	Then, $L$ is regular by \cref{constrsync:thm:sparsemonoid-regular}.
	Otherwise, by~\cref{constrsync:prop:slender},
	there exist $v,w \in L$
	such that $\{v,w\} \not\subseteq u^*$ for each $u \in \Sigma^*$.
	Since~$L$ is a monoid, we have $\{v,w\}^* \subseteq L$.
	By \cref{constr_sync:cor:code_growth_of_star}, $L$ has exponential growth.	
\end{proof}

By \cref{constr_sync:lem:two-words-code,constrsync:thm:monoid-sparseregular-or-exp} a code 
either contains at most one word or its Kleene star has exponential growth.

\begin{theorem} \label{constrsync:thm:reg-sparse-or-exp}
	A regular language is either bounded or has exponential growth.
\end{theorem}
\begin{proof}	
	Let $L \subseteq \Sigma^*$ be regular and $\automat{A}{Q}{\Sigma}{\delta}{q_0}{F}$
	be a trim PDFA accepting~$L$.
	As a shorthand, for $q \in Q$, set $L_{q} = L(\mathcal A_{q, \{q\}})$.
	By \cref{constr_sync:thm:regular_prefix_codes_and_automata}
	these are monoid languages.	
	Suppose $L_q$ has exponential growth for some $q \in Q$.
	Since~$\mathcal A$ is trim, we have $\delta(q_0, u) = q$
	and $\delta(q, v) \in F$ for some $u,v \in \Sigma^*$.
	This yields $u L_q v \subseteq L$.
	By \crefitem{ex:only_one_letter_infinitely}{ex:only_one_letter_infinitely:4}
	the language $u L_q v$ has exponential growth.
	Thus, also $L$ has exponential growth.
	Otherwise, 
	by \cref{constrsync:thm:monoid-sparseregular-or-exp} 
	for each $q \in Q$ we have $L_q \subseteq v_q^*$
	for some $v_q \in \Sigma^*$.
	Let 
	\[
	E = \{ u_1 \cdots u_m \in L \mid u_1,\ldots,u_m \in \Sigma \land 
	\delta(q_0, u_1 \cdots u_i) \notin \{q_0, \delta(q_0, u_1), \ldots, \delta(q_0, u_1 \cdots u_{i-1}) \}
	\}.
	\]
	The words in~$E$ are precisely those
	words $u \in L$
	for which there do not exist $u_1, u_2, u_3 \in \Sigma^*$
	with $u = u_1 u_2 u_3$, $|u_2| > 0$, and $\mu(p_0, u_1) = \mu(p_0, u_1 u_2)$.
	Each word in~$E$ has length at most~$|Q|$ and so~$E$ is finite. 
	We have\footnote{Note that $\delta(q_0, u_1\cdots u_m) \in F$ in \cref{constrsync:eqn:reg-sparse-or-exp}
		since~$E \subseteq L$.}
	\begin{equation}\label{constrsync:eqn:reg-sparse-or-exp}
		L = \bigcup_{u_1 \cdots u_m \in E} L_{q_0} u_1 L_{\delta(q_0, u_1)} u_2
		\cdots u_m L_{\delta(q_0, u_1\cdots u_m)}.
	\end{equation}
	Thus $L$ is contained in a finite union	of languages of the form 
	$v_{q_0}^* u_1 v_{q_1}^* u_2 \cdots u_m v_{q_m}^*$
	for $u_1 \cdots u_m \in E$ and $q_i = \delta(q_0, u_1\cdots u_i)$ 
	and so $L$ is bounded.\footnote{\Cref{constrsync:eqn:reg-sparse-or-exp} can also be used
	to prove ``$\ref{constr_sync:prop:letterboundedreg:item:PDFA} \Rightarrow~\ref{constr_sync:prop:letterboundedreg:item:letterbounded}$''
	in \cref{constr_sync:prop:letterboundedreg} since
	with the assumption that, for each $q \in Q$, we have $L_q \subseteq a^*$ for some $a \in \Sigma^*$
	we can then show $\rho(u) \le \max\{ \rho(v) \mid v \in E \} + 2|Q|$
	for $u \in L$.} 
	%Thus we obtain $L \subseteq \bigcup_{u_1 \cdots u_m \in E} v_{q_0}^* u_1 v_{\delta(q_0, u_1)}^* u_2 \cdots u_m v_{\delta(q_0, u_1\cdots u_m)}^*$ and so~$L$ is bounded.
\end{proof}

\begin{corollary} \label{constr_sync:thm:cf_bounded_equals_sparse}
	A regular language is bounded if and only if it is sparse.
\end{corollary} 
\begin{proof}
	By \cref{constr_sync:prop:sparse}, every bounded language is sparse. Conversely, a sparse regular language cannot have exponential growth and is therefore bounded by \cref{constrsync:thm:reg-sparse-or-exp}.	
\end{proof}

The following corollary collects some conditions characterizing languages
in $u^*$ that follow from previous results.

\begin{corollary}
	\label{constr_sync:cor:kleene_star_growth}
	Let $L \subseteq \Sigma^*$. Then the following are equivalent.	
\end{corollary}
\begin{enumerate}	
	\item \label{constrsync:cor:kleenestar:item:Lustar} $L \subseteq u^*$ for some $u \in \Sigma^*$.	
	
	\item \label{constrsync:cor:kleenestar:item:Lstarustar} $L^* \subseteq u^*$ for some $u \in \Sigma^*$.
	
	\item \label{constrsync:cor:kleenestar:item:regular} $L^* = u^n(u^d)^* \cup F$ for some $u \in \Sigma^*$, $n \ge 0$, $d \ge 0$ and finite $F \subseteq \{ u \}^{< n}$.\footnote{Note that in case $L = \{\varepsilon\}$ we 
		have $d = 0$ where the first three equivalences hold for any $u \in \Sigma^*$.}
	
	\item \label{constrsync:cor:kleenestar:item:sparse} $L^*$ is sparse.
	
	\item \label{constrsync:cor:kleenestar:item:notexponential} $L^*$ has not exponential growth.
	
	\item \label{constrsync:cor:kleenestar:item:vu} For all $u,v \in L$ there exists $w \in \Sigma^*$ 
	such that $\{u, v\} \subseteq w^*$.
\end{enumerate}
\begin{proof}
$\ref{constrsync:cor:kleenestar:item:Lustar}\Rightarrow\ref{constrsync:cor:kleenestar:item:Lstarustar}$: This follows since $L^* \subseteq (u^*)^* = u^*$ by \crefitem{prelim:langeqn}{prelim:langeqn:Kleenestar-idemp}.

$\ref{constrsync:cor:kleenestar:item:Lstarustar}\Rightarrow\ref{constrsync:cor:kleenestar:item:regular}$: This follows from \cref{constrsync:thm:sparsemonoid-regular} since $L^*$ is a monoid language.

$\ref{constrsync:cor:kleenestar:item:regular}\Rightarrow\ref{constrsync:cor:kleenestar:item:sparse}$: In this case, $L^* \subseteq u^*$ and so it is sparse by \cref{constr_sync:prop:sparse}.

$\ref{constrsync:cor:kleenestar:item:sparse}\Rightarrow\ref{constrsync:cor:kleenestar:item:notexponential}$: This is clear.

$\ref{constrsync:cor:kleenestar:item:notexponential}\Rightarrow\ref{constrsync:cor:kleenestar:item:vu}$: Since $L^*$ is a monoid language, if it does not has exponential growth,
then $L^* \subseteq u^*$ for some $u \in \Sigma^*$ by \cref{constrsync:thm:monoid-sparseregular-or-exp}.

$\ref{constrsync:cor:kleenestar:item:vu}\Rightarrow\ref{constrsync:cor:kleenestar:item:Lustar}$: This follows from \cref{constrsync:prop:slender}.
\end{proof}

%\section{General Assumptions on Constraint Languages and Constraint PDFAs}
\section{Assumptions on Constraint Languages and Constraint PDFAs}
\label{constr_sync:sec:assumptions}
 
%!TEX root = ./../../thesis-trier.tex

Consider a constrained synchronization problem \constrsync{L}
with some regular constraint language $L \subseteq \Sigma^*$.

\begin{remark}
	\label{constrsync:rem:alphabet}	
	Every constraint language also specifies an alphabet.
	As we show next, up to $\eqACZEROm$-reductions,
	the specific alphabet does not matter and we can always
	chose a smallest possible alphabet over which to consider the constraint
	language.
	For let $\Gamma \subseteq \Sigma$ and $L \subseteq \Gamma^*$ be regular.
	Then $L \subseteq \Sigma^*$
	and \constrsync{L} with~$L$ over $\Sigma$
	and \constrsync{L} with~$L$ over $\Gamma$
	are equivalent for~$\leACZEROm$-reductions:
	for $\mathcal A$ over $\Sigma$
	we can consider~$\mathcal A_{|\Gamma}$
	and, in the other direction for $\mathcal A$ over~$\Gamma$,
	we can add self-loops for all letters in $\Sigma \setminus \Gamma$ 
	to $\mathcal A$ to obtain an automaton over $\Sigma$.
	Both reductions can be implemented by $\AC^0$-circuits
	by either ignoring certain input bits
	or adding the identity matrix as a transition matrix.
\end{remark}

\IN{constraint PDFA}
In this context, we call a PDFA $\mathcal B$ accepting $L$
a \emph{constraint PDFA}.
%As \constrsync{\emptyset} is solvable in constant time, 
%as no automaton has a synchronizing word in $\emptyset$.
%So, we exclude the empty set from consideration
%as a constraint language, and assume $L \ne \emptyset$
%and that $\mathcal B$ is trim.
Since the constraint languages is assumed to be fixed for a given constrained synchronization 
problem, we can assume $\mathcal B$ is trim.
Furthermore, with respect to Remark~\ref{constrsync:rem:alphabet},
%if not said otherwise
we can choose the alphabet $\Sigma$ 
as small as possible.
So, when nothing else is said, we make the following assumptions.

\begin{enumerate}
	\item The constraint PDFA $\mathcal B$ is trim.
	
	\item For a single regular constraint language $L \subseteq \Sigma^*$,
	every letter is used in $L$, i.e.,
	$\Sigma^* a \Sigma^* \cap L \ne \emptyset$
	for each $a \in \Sigma$.
	
	% nicht immer sinnvoll, z.B. reduktion b^* auf ab^*
	%\item When considering two problems
	%$\constrsync{L_1}$ and $\constrsync{L_2}$
	%over alphabets $\Sigma_{1}$ and $\Sigma_2$ respectively
	%that share some letters
	%and we do not distinguish explicitly the different alphabets,
	%then both are over a
	%a common alphabet, i.e., $\Sigma_1 = \Sigma_2$ 
	%(as we can take $\Sigma_1 \cup \Sigma_2$ otherwise).
\end{enumerate}

Here, PDFAs are only used to describe constraint languages. 
%Hence, all other automata and those appearing
%as inputs to our problems are assumed to be complete. 

Due to \cref{constrsync:rem:alphabet}, we do not specify an alphabet in the problem
specification of \constrsync{L}.
%and might choose whatever alphabet  is most convenient.
Changing the alphabet is actually a (very) special case
of a more general way to encode
the input by so-called solid codes, see Remark~\ref{constrsync:rem:alphabet_solid_special_case}.	


Later, in \cref{rem:PDFA-for-constraint}, we justify the extension of our model to 
PDFAs rather than using only the automaton model with a completely defined transition function introduced earlier.

\section{Discussion}
\label{constrsync:notions:sec:discussion}

%!TEX root = ./../../springer-mono-format/book.tex
%!TEX root = ./../../thesis-trier.tex

\paragraph{Partial Functions.} 
In principle, everything that can be formulated with partial functions
can be recast using only (total) functions, and every PDFA
can be converted into an automaton as defined before by adding a trap
state at which all previously undefined transitions end. 
The reader might therefore ask what is gained from this generalization. 
In theoretical computer science, especially in computability theory (see~\cite{HopUll79,DBLP:books/daglib/0088160}), where a notion of effectiveness (or computability) is associated with functions via input-output machines, it is often natural to allow the ``result'' to be undefined. In this setting, a function may be undefined for inputs on which a machine is not required to produce a result, or on which it may fail to terminate by entering an infinite loop.
In our context, we introduced PDFAs based on partial functions.  
As we will also see in the next paragraph, it is convenient to specify specific instances of the constrained
synchronization problem by PDFAs.

\paragraph{Partial Automata.} We will not be concerned with (carefully or exactly)
synchronizing PDFAs in this work. The reason for introducing partial automata
is that they give a convenient specification for constraint languages.
More specifically:
\begin{enumerate}
	\item If $\mathcal A$ is a complete automaton having only coaccessible states, 
	then $\constrsync{L(\mathcal A)}$
	is equivalent to~$\constrsync{\Sigma^*}$
	see~\cref{constrsync:sec:constraut:unconstrained}, 
	\crefitem{constrsync:cor:upperset}{constrsync:cor:upperset:item:complete}, and \cref{rem:PDFA-for-constraint} for details.
	But non-coaccessible states can be discarded to give
	a PDFA accepting the same language.
    So, instead of carrying around these useless states it is more economical 
	to use trim PDFAs, in particular in our discussions of small constraint automata later in \cref{constr_sync:sec:two-state,sec:constr_sync:three_states}. 
	
	\item There is some freedom on how to distribute the start state and the final states 
	without altering the computational complexity of the corresponding constrained synchronization problem (see \cref{constr_sync:prop:final_states}). 
	This might give rise to many different regular expressions 
	for the resulting languages while the underlying automaton, apart from its start state and set of final states,
	remains the same.
	Because of this, we regard PDFAs as a better specification mechanism compared to regular expressions.
\end{enumerate}

The fact mentioned in \cref{constrsync:rem:alphabet}
that we can ignore letters that are not used in any transitions up to $\AC^0$-reductions
was stated \cite[Lem.~11]{FernauGHHVW19} for $\lelogm$-reductions.


\IN{synchronizing prefix code}
\paragraph{Codes and Overlap-Free Languages.}
Codes, and a connection to synchronizing automata,
appeared in our introductory
example on error-prone communication from \cref{intro:sec:examples:communication}.
In this example, a decoder modeled by a synchronizing automaton
allows recovery from errors by sending a synchronizing word
that realigns the decoder with the encoder.
This connection to codes has been noted since the early days
of synchronizing automata,
for example, by Liu~\cite{Liu63} and Laemmel~\cite{Laemmel56,Laemmel63,Laemmel69};
see also the surveys~\cite{KariVolkov2021,Volkov22},
the textbook~\cite{DBLP:books/daglib/0025093} 
or the theses~\cite{Biskup2008,DBLP:phd/hal/Ryzhikov20}.

Without reference to automata or regular languages,
Schützenberger~\cite{DBLP:journals/iandc/Schutzenberger67} called a prefix code $C \subseteq \Sigma^+$
\emph{synchronizing} if there exists $u \in C^*$
such that $\Sigma^* u \subseteq C^*$.\footnote{This definition implies 
	maximality of the code~$C$ w.r.t.~set inclusion.
	For consider $v \in \Sigma^*$.
	By assumption, $vu \in C^*$.
	Thus this word has a factorization
	in words from~$C$, which excludes
	that $C \cup \{v\}$ can be a code
	since $u \in C^*$ implies that it also 
	has a factorization in words from $C \cup \{v\}$.
	For a more general definition of synchronization applying to a wider class of prefix codes (and even general codes), see~\cite{DBLP:books/daglib/0025093}.}
Almost every prefix code that is maximal for set-inclusion
is synchronizing~\cite{DBLP:journals/tit/FreilingJTZ03}. 

We can obtain a decoding automaton for a regular prefix code $C$
via the constructions outlined in
\cref{syncmax:rem:root-tree-prefix-code,constr_sync:prop:aut_prefix_free_exact_sync,syncmax:rem:general-aut-prefix-code}.
Each time the decoding automaton enters a final state, this indicates
the detection of a code word in the input stream.
According to \cref{constr_sync:prop:aut_prefix_free_exact_sync},
this automaton admits exactly synchronizing words, and each such word
can be extended to an exactly synchronizing word ending in the unique
final state.
In case the automaton is complete, each exactly synchronizing word
is a synchronizing word in the usual sense as defined in \cref{part:intro},
Moreover, those exactly synchronizing words ending in the final state
are precisely the synchronizing words of~$C$.
For example, consider the complete automaton obtained from
\cref{intro:subfig:decoder}, whose start state and unique final state
coincide at $q_{\varepsilon}$.
The word $0110$, which synchronizes every state to $q_{\varepsilon}$,
is also a synchronizing word of the code.
Apart from the literature already mentioned above, for more information
on the relation of codes to synchronizing automata,
see also~\cite{BeaBerPer2011,DBLP:journals/tcs/CarpiD17,DBLP:conf/dlt/GusevP11,DBLP:phd/hal/Ryzhikov20}.

Overlap-free languages were introduced in~\cite{DBLP:conf/cocoon/HanW06},
but with a different definition
by defining~$L \subseteq \Sigma^*$ overlap-free if for all distinct $u,v \in L$
we have that $\pref{u} \cap \suff{v}$ and $\suff{u} \cap \pref{v}$
both equal $\{\varepsilon\}$ and $\pref{u} \cap \suff{u} = \{\varepsilon, u\}$
for all $u \in \Sigma^*$.\footnote{In fact, the original definition
demands equality to the empty set of $\pref{u}\cap \suff{v}$, which does not seem to be the intended
definition since the empty word is always a common prefix and suffix
of any two words,  and also the condition
$\pref{u}\cap \suff{u} = \{\varepsilon, u\}$ is missing that excludes a non-empty proper overlap
of a word with itself only a condition on distinct words is imposed. See also the footnote in~\cite[p.~33]{DBLP:journals/tcs/JurgensenS21} which
remarks on the problematic definition from~\cite{DBLP:conf/cocoon/HanW06}.}
This implies overlap-freeness in our sense, but also that~$L$
is a bifix code since if, without loss of generality,
$v \in \pref{u}$ for distinct $u,v \in \Sigma^*$, then $v \in \pref{u} \cap \suff{v}$ since $v \in \suff{v}$.
For example, $L = \{aba, b\}$ fulfills this definition of overlap-freeness,
but $\{a,ab\}$, which is overlap-free according to our definition,
is excluded with this definition. In terms of \cref{constrsync:def:codes}
this wider definition corresponds to replacing the condition
``$u \ne \varepsilon$ and $w \ne \varepsilon$'' from our definition
by ``$u \ne \varepsilon$ or $w \ne \varepsilon$''. We have not seen
this definition of overlap-freeness elsewhere and the one we adopted
seems to be the standard definition used, for example, 
also in~\cite{DBLP:journals/ijfcs/HanS11,DBLP:journals/iandc/Jurgensen08,DBLP:reference/hfl/JurgensenK97,DBLP:journals/tcs/JurgensenS21,DBLP:journals/jalc/KonstantinidisY10}.
In fact, the authors of~\cite{DBLP:conf/cocoon/HanW06}, a conference contribution,
later adopted the definition we gave here in their extended journal version~\cite{DBLP:journals/ijfcs/HanS11}
of the conference contribution.

As noted in \cref{constrsync:sec:codes:solid}, the characterization of solid codes
from \crefitem{thm:solidcode}{thm:solidcode:2} appears to be new.

We note that for codes that are regular languages, often interesting relations hold among the classes
considered in the present work. For example, in~\cite{DBLP:journals/jalc/Konstantinidis20}
it was proved that a regular bifix code is a finite union of regular infix codes.



% dazu noch was sagen? auch noch zu solid codes mehr sagen?
%\paragraph{The Levenshtein-Romanov Mapping.}


%(in particular, we can regard $C$ as a set of abstract symbols and take $\Gamma = C$).
\IN{free submonoid (of $\Sigma^*$)}
\IN{defect theorem}
\paragraph{Monoid Languages.} The free monoid $\Sigma^*$ has a central role in theoretical computer science
as our primary means of representing and processing information. As $\Sigma^*$ is a monoid, languages
that are submonoids of~$\Sigma^*$ form a natural object of study. A submonoid
having a generating set $C$ that is a code is said to be a \emph{free \mbox{(sub-)}monoid} of $\Sigma^*$. 
By \cref{constr_sync:prop:coding_morphisms} a bijection between some set $\Gamma$
and~$C$ extends to a coding morphism.
%\footnote{In particular,
%	for $X = C$ we obtain an isomorphism between $C^*$ considered as 
%	the set of finite sequences	and the monoid generated by $C$ in $\Sigma^*$, and it is isomorphic to 
%	a free monoid over $X$ as defined in \SEC~\ref{sec:formal_language_theory}.}	 
Conversely, the image of a coding morphism
yields a submonoid generated by a code. So, free submonoids of $\Sigma^*$
are precisely the images of coding morphisms, i.e., isomorphic
images of free monoids over some set~$\Gamma$ as
introduced in \cref{sec:formal_language_theory} as sets of finite sequences,
which explains why we call them free monoids as well.
The relevance of these submonoids lies in the transmission, translation and encoding of information w.r.t.~different representations.
Every language over $\Sigma$
is contained in a smallest free submonoid of $\Sigma^*$~\cite{lothaire_1997}.
The \emph{defect theorem}~\cite[Thm. 1.2.5]{lothaire_1997}
states that in case $L \subseteq \Sigma^*$ is a finite language that is not a code,
there exists a finite code $C \subseteq \Sigma^+$
with $|C| \le |L| - 1$ such that $L \subseteq C^*$.

\IN{thin language}
\IN{slender language}
\IN{dense language}

 
\paragraph{The Growth Of Languages and Bounded Languages.}

Berman and Hartmanis~\cite{DBLP:journals/siamcomp/BermanH77}
introduced the notion of sparse languages to study encodings of \NP-complete problems.
They conjectured that, under the assumption $\PTIME \ne \NP$,
sparse languages cannot encode \NP-complete problems with respect to $\lePm$-reductions.
This conjecture was later proved by Mahaney~\cite{DBLP:journals/jcss/Mahaney82}
(see also the survey~\cite{DBLP:conf/mfcs/HartmanisM80}).
Surveys of sparse regular languages can be found
in~\cite[Chap.~2, Sect.~3.5]{DBLP:reference/hfl/Yu97}
and~\cite[Sect.~XII.4.2]{Pin2022}.

P\u{a}un and Salomaa~\cite{DBLP:journals/dam/PaunS95}
introduced \emph{thin} and \emph{slender}
languages where $|L \cap \Sigma^n| \le 1$ and $|L \cap \Sigma^n| \le c$, respectively,
for sufficiently large $n$.
Slender context-free languages were investigated in~\cite{DBLP:journals/tcs/Ilie94,DBLP:journals/tcs/Raz97}.
Bertoni et al.~\cite{DBLP:conf/stacs/BertoniGS87} investigated the computational complexity
of the function $n \mapsto |L \cap \Sigma^n|$
on context-free languages.
D'Alessandro et al.~\cite{DBLP:journals/tcs/DAlessandroIV06} studied this function for sparse context-free languages.

Ginsburg and Spanier~\cite{Ginsburg1964BOUNDEDAL}
introduced bounded languages and investigated bounded context-free languages.
They later studied bounded regular languages
in~\cite{GinsburgSpanier66}.
Several authors extended \cref{constrsync:thm:reg-sparse-or-exp}
to context-free languages~\cite{DBLP:journals/jcss/BridsonG02,DBLP:conf/stacs/IbarraR86,DBLP:journals/tcs/Incitti01,DBLP:journals/eik/LatteuxT84,DBLP:journals/tcs/Raz97,Trofimov81}.
Indexed languages do not satisfy this dichotomy: there
exist indexed languages of intermediate growth~\cite{DBLP:journals/tcs/GrigorchukM99}.
Ginsburg and Spanier~\cite[Thm.~5.2]{Ginsburg1964BOUNDEDAL}
proved that boundedness is decidable for context-free languages.
Consequently, the growth type of a context-free language is decidable.
Gawrychowski et al.~\cite{DBLP:journals/ijfcs/GawrychowskiKRS10}
gave a polynomial-time algorithm for deciding the growth type of a context-free language.

Ilie et al.~\cite{DBLP:journals/ita/IlieRS00}
characterized sparse context-free languages via special Dyck languages.
Szilard et al.~\cite[Cor.~1]{DBLP:conf/mfcs/SzilardYZS92}
characterized sparse regular languages by regular expressions of star-height one.
Shur~\cite{DBLP:conf/csr/Shur08}
related the growth of a regular language to structural properties of its accepting automata.

A language $L \subseteq \Sigma^*$ is called \emph{dense}
if $\Sigma^* w \Sigma^* \cap L \ne \emptyset$
for every $w \in \Sigma^*$~\cite[Sect.~XII.4.1]{Pin2022}.
Regular dense languages over alphabets of size at least two
are never sparse~\cite[Cor.~4.28]{Pin2022}.

Our terminology for letter-bounded languages is not standard.
For example,~\cite{DBLP:journals/iandc/IbarraR16}
uses the term letter-bounded for the strict variant in which all letters are distinct,
whereas~\cite{DBLP:journals/tcs/CarpiDIM21}
refers to this variant as distinct-letter-bounded.

The function~$\rho$ via equal-letter runs occurs naturally in text compression
algorithms, see~\cite{DBLP:conf/dlt/FrosiniMRRS22,Prezza16},
and combinatorics on words~\cite{DBLP:journals/tcs/BiagiCLR25}
since it is related to run-length encodings. In the theory of synchronizing automata,
number of equal-letter runs in synchronizing words have also 
been investigated in~\cite{DBLP:journals/jalc/DonZ19}.